% =============================================================================
% Chuỗi Markov - Tài liệu Song ngữ Việt-Anh
% Dựa trên: "Notes on Markov Chains" - Nicolas Privault (2011)
% Biên soạn cho nghiên cứu sinh PhD - Đại học Debrecen
% Biên dịch và trình bày lại bằng XeLaTeX
% =============================================================================
\documentclass[12pt,a4paper]{book}

% --- Font & Vietnamese support (XeLaTeX) ---
\usepackage{fontspec}
\usepackage{polyglossia}
\setdefaultlanguage{vietnamese}
\setotherlanguage{english}

% --- Math packages ---
\usepackage{amsmath,amssymb,amsthm}
\usepackage{mathtools}
\usepackage{bm}

% --- Layout for iPad reading ---
\usepackage[margin=2cm]{geometry}
\usepackage{setspace}
\onehalfspacing

% --- Colors and boxes ---
\usepackage[dvipsnames]{xcolor}
\usepackage{tcolorbox}
\tcbuselibrary{theorems,breakable,skins}

% --- Other ---
\usepackage{enumitem}
\usepackage{hyperref}
\hypersetup{
  colorlinks=true,
  linkcolor=MidnightBlue,
  urlcolor=RoyalBlue,
  citecolor=ForestGreen
}
\usepackage{bookmark}

% =============================================================================
% Custom theorem environments
% =============================================================================
\newtheoremstyle{vndef}
  {10pt}{10pt}{\itshape}{}{\bfseries}{.}{.5em}{}
\newtheoremstyle{vnthm}
  {10pt}{10pt}{\itshape}{}{\bfseries}{.}{.5em}{}
\newtheoremstyle{vnplain}
  {10pt}{10pt}{}{}{\bfseries}{.}{.5em}{}

\theoremstyle{vnthm}
\newtheorem{theorem}{Định lý}[chapter]
\newtheorem{proposition}[theorem]{Mệnh đề}
\newtheorem{corollary}[theorem]{Hệ quả}
\newtheorem{lemma}[theorem]{Bổ đề}

\theoremstyle{vndef}
\newtheorem{definition}[theorem]{Định nghĩa}

\theoremstyle{vnplain}
\newtheorem{example}[theorem]{Ví dụ}
\newtheorem{exercise}[theorem]{Bài tập}
\newtheorem{remark}[theorem]{Nhận xét}

% =============================================================================
% Custom tcolorbox environments
% =============================================================================

% Summary box (Tóm tắt)
\newtcolorbox{tomtat}{
  colback=blue!5!white,
  colframe=MidnightBlue,
  title={\textbf{Tóm tắt}},
  fonttitle=\bfseries,
  breakable,
  arc=2mm,
  boxrule=0.8pt
}

% Definition box
\newtcolorbox{dinhnghia}[1][]{
  colback=green!3!white,
  colframe=ForestGreen,
  title={\textbf{Định nghĩa} #1},
  fonttitle=\bfseries,
  breakable,
  arc=2mm,
  boxrule=0.8pt
}

% Theorem box
\newtcolorbox{dinhly}[1][]{
  colback=red!3!white,
  colframe=BrickRed,
  title={\textbf{Định lý} #1},
  fonttitle=\bfseries,
  breakable,
  arc=2mm,
  boxrule=0.8pt
}

% Key formula box
\newtcolorbox{congthuc}{
  colback=yellow!5!white,
  colframe=Goldenrod,
  breakable,
  arc=2mm,
  boxrule=0.8pt
}

% =============================================================================
% Shortcut commands
% =============================================================================
\newcommand{\E}{\mathbb{E}}
\newcommand{\PP}{\mathbb{P}}
\newcommand{\R}{\mathbb{R}}
\newcommand{\N}{\mathbb{N}}
\newcommand{\Z}{\mathbb{Z}}
\newcommand{\ind}{\mathbf{1}}
\newcommand{\en}[1]{\textit{#1}}  % English term in italics
\newcommand{\term}[1]{\textbf{#1}}  % Key Vietnamese term in bold

% =============================================================================
% Document
% =============================================================================
\begin{document}

% --- Title page ---
\begin{titlepage}
\centering
\vspace*{3cm}
{\Huge\bfseries Chuỗi Markov\\[0.5cm]}
{\Large\itshape Tài liệu Song ngữ Việt--Anh\\[0.3cm]}
{\large Bilingual Vietnamese--English Notes\\[2cm]}
{\large Dựa trên \textit{Notes on Markov Chains} của Nicolas Privault\\[0.5cm]}
{\normalsize MAS328 -- Nanyang Technological University, 2011\\[3cm]}
{\normalsize Biên soạn phục vụ nghiên cứu PhD\\
Đại học Debrecen (\textit{University of Debrecen})\\[1cm]}
\today
\end{titlepage}

% --- Table of contents ---
\tableofcontents
\newpage

% --- Chapters ---
\chapter{Giới thiệu: Quá trình ngẫu nhiên \& Bước đi ngẫu nhiên\\{\normalsize\textit{Introduction: Stochastic Processes \& Random Walks}}}
\label{ch:intro}

% ============================================================================
\section{Quá trình ngẫu nhiên là gì? (What is a Stochastic Process?)}
\label{sec:what-is-sp}

\textbf{Câu chuyện:} Hãy tưởng tượng bạn đang đứng trên cầu nhìn xuống một chiếc lá trôi trên dòng sông.
Vị trí chiếc lá thay đổi theo thời gian, nhưng bạn không thể dự đoán chính xác nó sẽ ở đâu ---
dòng nước, gió, và xoáy nước đều tạo ra sự ngẫu nhiên.
Nếu ta ghi lại vị trí chiếc lá tại mỗi thời điểm, ta có một \term{quá trình ngẫu nhiên} (\en{stochastic process}).

\begin{dinhnghia}[(Quá trình ngẫu nhiên -- \en{Stochastic process})]
\label{def:stochastic-process}
Một \term{quá trình ngẫu nhiên} (\en{stochastic process}) là một họ các biến ngẫu nhiên
\[
  \{X_t : t \in T\}
\]
được định nghĩa trên cùng một không gian xác suất $(\Omega, \mathcal{F}, \PP)$,
trong đó $T$ là \term{tập chỉ số thời gian} (\en{index set}) và mỗi $X_t$ nhận giá trị trong
một \term{không gian trạng thái} (\en{state space}) $S$.
\end{dinhnghia}

\begin{tomtat}
Nói đơn giản: một quá trình ngẫu nhiên là một ``đại lượng thay đổi theo thời gian một cách ngẫu nhiên''.
Tại mỗi thời điểm $t$, ta có một biến ngẫu nhiên $X_t$ --- và tập hợp tất cả chúng tạo thành quá trình.
\end{tomtat}

\textbf{Hai loại thời gian:}
\begin{itemize}[nosep]
  \item \textbf{Thời gian rời rạc} (\en{discrete-time}): $T = \{0, 1, 2, \ldots\} = \N$.
    Ví dụ: giá cổ phiếu đóng cửa mỗi ngày, số khách hàng mỗi giờ.
  \item \textbf{Thời gian liên tục} (\en{continuous-time}): $T = [0, +\infty)$.
    Ví dụ: nhiệt độ tại một địa điểm, mực nước sông theo thời gian thực.
\end{itemize}

\textbf{Ba loại không gian trạng thái:}
\begin{itemize}[nosep]
  \item \textbf{Hữu hạn} (\en{finite}): $S = \{1, 2, \ldots, N\}$.
    Ví dụ: thời tiết $\in \{\text{nắng, mưa, mây}\}$.
  \item \textbf{Đếm được} (\en{countable}): $S = \Z$ hoặc $S = \N$.
    Ví dụ: số khách trong hàng đợi, vị trí bước đi ngẫu nhiên.
  \item \textbf{Không đếm được} (\en{uncountable}): $S = \R$ hoặc $S = \R^d$.
    Ví dụ: giá cổ phiếu, nhiệt độ.
\end{itemize}

\begin{dinhnghia}[(Đường mẫu -- \en{Sample path})]
\label{def:sample-path}
Cho một kết quả cố định $\omega \in \Omega$, ánh xạ
\[
  t \mapsto X_t(\omega)
\]
được gọi là một \term{đường mẫu} (\en{sample path}) hay \term{quỹ đạo} (\en{trajectory}) của quá trình.
Mỗi đường mẫu là một ``kịch bản'' cụ thể --- một lần ``chạy'' của quá trình.
\end{dinhnghia}

\begin{example}[Các ví dụ thực tế]
\label{ex:real-world-sp}
\begin{enumerate}[nosep]
  \item \textbf{Giá cổ phiếu:} $X_n$ = giá đóng cửa ngày thứ $n$. Thời gian rời rạc, không gian trạng thái liên tục.
  \item \textbf{Thời tiết:} $X_n \in \{\text{nắng, mưa}\}$. Thời gian rời rạc, không gian trạng thái hữu hạn.
  \item \textbf{Hàng đợi:} $X_t$ = số khách hàng trong hệ thống tại thời điểm $t$. Thời gian liên tục, không gian trạng thái đếm được.
  \item \textbf{DNA:} Dãy nucleotide $X_1, X_2, \ldots \in \{A, T, G, C\}$. Chỉ số là vị trí, không gian trạng thái hữu hạn.
\end{enumerate}
\end{example}

% ============================================================================
\section{Hai họ lớn: Markov và Martingale (Two Big Families)}
\label{sec:two-families}

Trong thế giới các quá trình ngẫu nhiên, hai họ quan trọng nhất là:

\textbf{1. Quá trình Markov} (\en{Markov process}): Tương lai chỉ phụ thuộc hiện tại, không phụ thuộc quá khứ.
\[
  \PP(X_{n+1} = j \mid X_n = i, X_{n-1}, \ldots, X_0) = \PP(X_{n+1} = j \mid X_n = i).
\]
Đây là tính chất ``không nhớ'' (\en{memoryless}). Khi không gian trạng thái là rời rạc,
ta gọi đây là \term{chuỗi Markov} (\en{Markov chain}) --- chủ đề chính của tài liệu này.
Ta sẽ nghiên cứu chi tiết từ Chương~1.

\textbf{2. Martingale}: Quá trình ``trò chơi công bằng'' --- kỳ vọng của giá trị tương lai,
khi biết quá khứ, đúng bằng giá trị hiện tại:
\[
  \E[X_{n+1} \mid X_n, X_{n-1}, \ldots, X_0] = X_n.
\]
Ý nghĩa: không có lợi thế hay bất lợi trong dài hạn. Ví dụ: tổng tiền thắng/thua trong một trò chơi công bằng.

\begin{tomtat}
\begin{center}
\begin{tabular}{lll}
\textbf{Họ} & \textbf{Ý tưởng chính} & \textbf{Ứng dụng} \\
\hline
Markov & ``Không nhớ'' & Chuỗi DNA, PageRank, hàng đợi \\
Martingale & ``Trò chơi công bằng'' & Tài chính, lý thuyết cờ bạc \\
\end{tabular}
\end{center}
Tài liệu này tập trung vào \textbf{chuỗi Markov}. Martingale chỉ xuất hiện ngắn gọn khi cần.
\end{tomtat}

% ============================================================================
\section{Bước đi ngẫu nhiên (Random Walks)}
\label{sec:random-walks}

Bước đi ngẫu nhiên là ví dụ đơn giản nhất --- và quan trọng nhất --- của cả quá trình Markov lẫn martingale.
Nó là ``Hello World'' của quá trình ngẫu nhiên.

\begin{dinhnghia}[(Bước đi ngẫu nhiên -- \en{Random walk})]
\label{def:random-walk}
Cho $(X_k)_{k \geq 1}$ là dãy các biến ngẫu nhiên \term{độc lập và cùng phân phối} (\en{i.i.d.}).
\term{Bước đi ngẫu nhiên} (\en{random walk}) được định nghĩa bởi $S_0 = 0$ và
\begin{equation}\label{eq:rw-def}
  S_n = X_1 + X_2 + \cdots + X_n, \qquad n \geq 1.
\end{equation}
\end{dinhnghia}

\begin{dinhnghia}[(Bước đi Bernoulli -- \en{Bernoulli random walk})]
\label{def:bernoulli-walk}
Khi mỗi bước $X_k$ chỉ nhận hai giá trị $\{+1, -1\}$ với
\[
  \PP(X_k = +1) = p, \qquad \PP(X_k = -1) = q = 1 - p,
\]
ta gọi $(S_n)_{n \geq 0}$ là \term{bước đi Bernoulli} (\en{Bernoulli random walk}).

Khi $p = q = 1/2$, ta gọi đây là bước đi \term{đối xứng} (\en{symmetric random walk}).
\end{dinhnghia}

\begin{example}[Người say rượu -- \en{Drunkard's walk}]
\label{ex:drunkard}
Một người say rượu đứng tại gốc tọa độ trên một con đường thẳng.
Mỗi bước, anh ta bước sang phải ($+1$) hoặc sang trái ($-1$) với xác suất bằng nhau $p = q = 1/2$.
Sau $n$ bước, vị trí của anh ta là $S_n$.

Câu hỏi tự nhiên: anh ta có quay lại gốc không? Đi xa bao nhiêu? Ta sẽ trả lời trong các mục tiếp theo.
\end{example}

% --- 3.1 Mean and Variance ---
\subsection{Trung bình và phương sai (Mean and Variance)}
\label{subsec:mean-var}

Ta tính kỳ vọng và phương sai của $S_n$ cho bước đi Bernoulli.

Kỳ vọng của mỗi bước:
\[
  \E[X_k] = (+1) \cdot p + (-1) \cdot q = p - q.
\]

Do đó, bằng tính tuyến tính của kỳ vọng:
\begin{equation}\label{eq:mean-Sn}
  \E[S_n] = n\,\E[X_1] = n(p - q).
\end{equation}

Phương sai của mỗi bước:
\begin{align*}
  \operatorname{Var}[X_k] &= \E[X_k^2] - (\E[X_k])^2 \\
  &= 1 \cdot q + 1 \cdot p - (p-q)^2 \\
  &= 1 - (2p-1)^2 = 4pq = 4p(1-p).
\end{align*}

Do các $X_k$ độc lập:
\begin{equation}\label{eq:var-Sn}
  \operatorname{Var}[S_n] = n\,\operatorname{Var}[X_1] = 4npq.
\end{equation}

\begin{congthuc}
\textbf{Trung bình và phương sai của bước đi Bernoulli:}
\[
  \E[S_n] = n(p - q), \qquad \operatorname{Var}[S_n] = 4npq.
\]
\begin{itemize}[nosep]
  \item Khi $p > q$: bước đi ``lệch phải'' (\en{positive drift}), $\E[S_n] > 0$.
  \item Khi $p < q$: bước đi ``lệch trái'' (\en{negative drift}), $\E[S_n] < 0$.
  \item Khi $p = q = 1/2$: bước đi đối xứng, $\E[S_n] = 0$, $\operatorname{Var}[S_n] = n$.
    Trung bình bằng $0$ nhưng phương sai tăng tuyến tính --- bước đi ``lang thang'' ngày càng xa gốc.
\end{itemize}
\end{congthuc}

% --- 3.2 Distribution ---
\subsection{Phân phối (Distribution)}
\label{subsec:distribution}

Ta muốn tính $\PP(S_n = k)$ cho bước đi Bernoulli.

Nhận xét quan trọng: $S_n$ chỉ có thể nhận các giá trị cùng tính chẵn lẻ với $n$.
Cụ thể, $S_{2n}$ chỉ có thể nhận giá trị chẵn, và $S_{2n+1}$ chỉ nhận giá trị lẻ.

Nếu $S_{2n} = 2k$, tức là trong $2n$ bước có $l = n+k$ bước đi lên và $2n - l = n - k$ bước đi xuống
(vì $l - (2n-l) = 2l - 2n = 2k$).

Xác suất của một đường đi cụ thể có $n+k$ bước lên và $n-k$ bước xuống là $p^{n+k}q^{n-k}$.
Số cách chọn vị trí cho $n+k$ bước lên trong $2n$ bước là $\binom{2n}{n+k}$.

\begin{dinhly}[(Phân phối của $S_{2n}$)]
\label{thm:dist-S2n}
\begin{equation}\label{eq:dist-S2n}
  \PP(S_{2n} = 2k) = \binom{2n}{n+k} p^{n+k}\,q^{n-k}, \qquad -n \leq k \leq n.
\end{equation}
\end{dinhly}

Đặc biệt, xác suất quay lại gốc tại thời điểm $2n$ là:
\begin{equation}\label{eq:return-origin}
  \PP(S_{2n} = 0) = \binom{2n}{n} p^n\,q^n = \binom{2n}{n}(pq)^n.
\end{equation}

\begin{example}[Bước đi đối xứng $p = q = 1/2$]
\label{ex:symmetric-dist}
Khi $p = q = 1/2$:
\[
  \PP(S_{2n} = 0) = \binom{2n}{n}\frac{1}{2^{2n}} = \binom{2n}{n}\frac{1}{4^n}.
\]
Sử dụng \term{xấp xỉ Stirling} (\en{Stirling's approximation}) $n! \approx \sqrt{2\pi n}\,(n/e)^n$:
\begin{equation}\label{eq:stirling-return}
  \PP(S_{2n} = 0) \approx \frac{1}{\sqrt{\pi n}}, \qquad n \to \infty.
\end{equation}
Xác suất này tiến về $0$ khi $n \to \infty$, nhưng rất chậm (tỉ lệ $1/\sqrt{n}$).
\end{example}

% --- 3.3 First Returns to Zero ---
\subsection{Quay lại gốc lần đầu (First Returns to Zero)}
\label{subsec:first-return}

Đặt
\begin{equation}\label{eq:tau0}
  \tau_0 = \inf\{n \geq 1 : S_n = 0\}
\end{equation}
là \term{thời gian quay lại gốc lần đầu} (\en{first return time to zero}).

\begin{dinhly}[(Xác suất quay lại gốc -- Privault, Mệnh đề 3.11)]
\label{thm:first-return}
Xác suất mà bước đi quay lại gốc trong thời gian hữu hạn là:
\begin{equation}\label{eq:prob-return}
  \PP(\tau_0 < \infty) = 2\min(p, q).
\end{equation}
Đặc biệt:
\begin{itemize}[nosep]
  \item Khi $p = q = 1/2$ (đối xứng): $\PP(\tau_0 < \infty) = 1$.
    Bước đi \textbf{chắc chắn} quay lại gốc.
  \item Khi $p \neq q$: $\PP(\tau_0 < \infty) < 1$.
    Có xác suất dương rằng bước đi \textbf{không bao giờ} quay lại gốc.
\end{itemize}
\end{dinhly}

\begin{tomtat}
Đây là một kết quả bất ngờ: bước đi đối xứng luôn quay lại gốc (xác suất 1),
nhưng \textbf{thời gian chờ trung bình} là vô hạn!
\[
  p = q = 1/2 \implies \PP(\tau_0 < \infty) = 1, \quad \text{nhưng} \quad \E[\tau_0] = +\infty.
\]
Kết quả này liên quan trực tiếp đến \textbf{tính hồi quy} (\en{recurrence}) của bước đi ngẫu nhiên 1 chiều,
mà ta sẽ gặp lại dưới dạng Định lý Polya trong Chương~3.
\end{tomtat}

% ============================================================================
\section{Từ bước đi ngẫu nhiên đến chuỗi Markov (From Random Walks to Markov Chains)}
\label{sec:rw-to-mc}

Bước đi ngẫu nhiên là một chuỗi Markov! Ta kiểm tra tính chất Markov:
do các $X_k$ độc lập, giá trị $S_{n+1} = S_n + X_{n+1}$ chỉ phụ thuộc vào $S_n$ và $X_{n+1}$,
không phụ thuộc vào $S_0, S_1, \ldots, S_{n-1}$:
\[
  \PP(S_{n+1} = j \mid S_n = i, S_{n-1}, \ldots, S_0) = \PP(S_{n+1} = j \mid S_n = i).
\]

Tuy nhiên, chuỗi Markov tổng quát hơn nhiều: xác suất chuyển $\PP(X_{n+1} = j \mid X_n = i)$
có thể là \textbf{bất kỳ giá trị nào} trong $[0,1]$ (miễn là tổng hàng bằng $1$),
không nhất thiết phải đối xứng hay đồng nhất.

\begin{congthuc}
\textbf{Bản đồ các chương:}
\begin{center}
\begin{tabular}{cl}
\textbf{Chương} & \textbf{Nội dung} \\
\hline
0 (đây) & Quá trình ngẫu nhiên, bước đi ngẫu nhiên \\
1 & Chuỗi Markov thời gian rời rạc, ma trận chuyển \\
2 & Phân tích bước đầu tiên: xác suất chạm, thời gian chạm \\
3 & Phân loại trạng thái: hồi quy, tạm thời, tuần hoàn \\
4 & Phân phối dừng và giới hạn \\
5 & Chuỗi Markov thời gian liên tục \\
A & Phụ lục: nền tảng xác suất \\
\end{tabular}
\end{center}
\end{congthuc}

\textbf{Câu chuyện cầu nối:} Hãy tưởng tượng một con kiến trên bàn cờ $3 \times 3$.
Mỗi bước, con kiến chọn ngẫu nhiên một ô kề cạnh để di chuyển.
Nếu bàn cờ chỉ có 1 chiều (một đường thẳng), con kiến thực hiện \emph{bước đi ngẫu nhiên}.
Nhưng trên bàn cờ 2 chiều, xác suất chuyển phụ thuộc vào vị trí
(ô góc có 2 hàng xóm, ô cạnh có 3, ô giữa có 4) --- đây là chuỗi Markov tổng quát.
Chương~1 sẽ khám phá thế giới này.

% ============================================================================
\section{Bài tập (Exercises)}
\label{sec:exercises-intro}

\begin{exercise}[Đường mẫu]
\label{ex:sample-paths}
Xét bước đi Bernoulli đối xứng ($p = q = 1/2$) với $n = 6$ bước.
\begin{enumerate}[nosep]
  \item Có bao nhiêu đường mẫu có thể? (Gợi ý: mỗi bước có 2 lựa chọn.)
  \item Liệt kê tất cả các đường mẫu kết thúc tại $S_6 = 2$.
  \item Tính $\PP(S_6 = 2)$ bằng hai cách: đếm trực tiếp và dùng công thức~\eqref{eq:dist-S2n}.
\end{enumerate}
\end{exercise}

\begin{exercise}[Phân phối $S_4$]
\label{ex:dist-S4}
Xét bước đi Bernoulli $(S_n)_{n \in \N}$ với $S_0 = 0$, xác suất tiến $p$ và xác suất lùi $1-p$.
\begin{enumerate}[nosep]
  \item Liệt kê tất cả các đường đi có thể dẫn đến $S_4 = 2$.
  \item Chứng minh rằng $\PP(S_4 = 2) = \binom{4}{3}p^3(1-p)$.
  \item Chứng minh công thức tổng quát: nếu $n+k$ chẵn và $|k| \leq n$,
    \[
      \PP(S_n = k) = \binom{n}{(n+k)/2} p^{(n+k)/2}\,(1-p)^{(n-k)/2}.
    \]
\end{enumerate}
\en{(Adapted from Privault, Exercise 3.1)}
\end{exercise}

\begin{exercise}[Trung bình và phương sai]
\label{ex:mean-var-calc}
Xét bước đi ngẫu nhiên $S_n = X_1 + \cdots + X_n$ với $\PP(X_k = +1) = 0.6$, $\PP(X_k = -1) = 0.4$.
\begin{enumerate}[nosep]
  \item Tính $\E[S_{100}]$ và $\operatorname{Var}[S_{100}]$.
  \item Sử dụng bất đẳng thức Chebyshev, tìm cận trên cho $\PP(|S_{100} - 20| \geq 30)$.
  \item Giải thích trực quan tại sao phương sai tăng theo $n$.
\end{enumerate}
\end{exercise}

\begin{exercise}[Quay lại gốc]
\label{ex:return-to-zero}
Cho bước đi Bernoulli đối xứng ($p = q = 1/2$).
\begin{enumerate}[nosep]
  \item Tính $\PP(S_2 = 0)$, $\PP(S_4 = 0)$, $\PP(S_6 = 0)$ bằng công thức~\eqref{eq:return-origin}.
  \item Kiểm tra rằng các giá trị trên phù hợp với xấp xỉ Stirling~\eqref{eq:stirling-return}.
  \item Giải thích tại sao $\PP(S_{2n+1} = 0) = 0$ cho mọi $n$.
\end{enumerate}
\end{exercise}

\begin{exercise}[Tính chất Markov của bước đi ngẫu nhiên]
\label{ex:rw-markov}
Cho bước đi Bernoulli $S_n = X_1 + \cdots + X_n$ với các $X_k$ i.i.d.
\begin{enumerate}[nosep]
  \item Chứng minh rằng $(S_n)_{n \geq 0}$ có tính chất Markov.

    \emph{Gợi ý:} Sử dụng tính độc lập của $X_{n+1}$ với $(X_1, \ldots, X_n)$.
  \item Viết ma trận chuyển $P$ cho bước đi Bernoulli trên không gian trạng thái $\Z$.
    Hàng $i$, cột $j$: $P_{i,j} = ?$
  \item Bước đi đối xứng ($p = 1/2$) có phải martingale không? Kiểm tra $\E[S_{n+1} \mid S_n] = S_n$.
\end{enumerate}
\end{exercise}

\chapter{Chuỗi Markov thời gian rời rạc\\{\normalsize\textit{Discrete-Time Markov Chains}}}
\label{ch:discrete}

% ============================================================================
\section{Tại sao chuỗi Markov quan trọng? (Why Markov Chains Matter)}
\label{sec:why-markov}

\textbf{Câu chuyện:} Hãy tưởng tượng bạn đang chơi cờ tỷ phú (Monopoly).
Bạn đang ở ô ``Nhà ga''. Bạn sẽ đi đâu tiếp theo?
Điều đó chỉ phụ thuộc vào ô hiện tại và kết quả gieo xúc xắc ---
không quan trọng bạn đã đi qua những ô nào trước đó.
Đây chính là ý tưởng cốt lõi của chuỗi Markov: \textbf{tương lai chỉ phụ thuộc hiện tại, không phụ thuộc quá khứ}.

\textbf{Ứng dụng thực tế:}
\begin{itemize}[nosep]
  \item \textbf{Thời tiết:} Ngày mai mưa hay nắng chỉ phụ thuộc hôm nay, không phụ thuộc tuần trước.
  \item \textbf{Xếp hạng Google (PageRank):} Người lướt web ``nhảy'' ngẫu nhiên giữa các trang.
  \item \textbf{Tài chính:} Giá cổ phiếu ngày mai phụ thuộc giá hôm nay (mô hình đơn giản).
  \item \textbf{Sinh học:} Kích thước quần thể thế hệ sau phụ thuộc thế hệ hiện tại.
  \item \textbf{Truyền thông:} Xếp hạng tín dụng năm sau phụ thuộc xếp hạng năm nay.
\end{itemize}

% ============================================================================
\section{Tính chất Markov (Markov Property)}
\label{sec:markov-property}

\textbf{Câu chuyện:} Hãy tưởng tượng một con kiến trên bàn cờ. Con kiến chỉ biết ô nó đang đứng.
Nó không nhớ đã đi qua bao nhiêu ô, cũng không nhớ con đường đã đi.
Quyết định bước tiếp theo chỉ dựa vào vị trí hiện tại --- đây là ``con kiến Markov''.

\begin{dinhnghia}[(Tính chất Markov -- \en{Markov property})]
Một \term{quá trình ngẫu nhiên} (\en{stochastic process}) $(Z_n)_{n\in\N}$ nhận giá trị trong $\Z$
được gọi là \term{Markov}, hay có \term{tính chất Markov}, nếu với mọi $n \geq 1$,
phân phối xác suất của $Z_{n+1}$ chỉ được xác định bởi trạng thái $Z_n$ tại thời điểm $n$,
và không phụ thuộc vào các giá trị quá khứ $Z_k$ với $k \leq n-1$.

Nói cách khác, với mọi $n \geq 1$ và mọi $j, i_n, \ldots, i_1 \in \Z$:
\[
  \PP(Z_{n+1} = j \mid Z_n = i_n,\; Z_{n-1} = i_{n-1},\; \ldots,\; Z_0 = i_0)
  = \PP(Z_{n+1} = j \mid Z_n = i_n).
\]
\end{dinhnghia}

\begin{tomtat}
Nói đơn giản: nếu bạn biết trạng thái hiện tại $Z_n$,
thì biết thêm quá khứ $(Z_0, Z_1, \ldots, Z_{n-1})$ cũng không giúp dự đoán tương lai $Z_{n+1}$ tốt hơn.
Đây chính là tính chất ``không nhớ'' (\en{memoryless property}).

Ví dụ con kiến: biết con kiến đang ở ô $A$ là đủ. Biết thêm ``nó đi từ $C \to B \to A$'' không thay đổi xác suất bước tiếp theo.
\end{tomtat}

Đặc biệt ta có:
\[
  \PP(Z_{n+1} = j \mid Z_n = i_n,\; Z_{n-1} = i_{n-1}) = \PP(Z_{n+1} = j \mid Z_n = i_n),
\]
và:
\[
  \PP(Z_2 = j \mid Z_1 = i_1,\; Z_0 = i_0) = \PP(Z_2 = j \mid Z_1 = i_1).
\]

Mặt khác, các \term{xác suất chuyển bậc một} (\en{first order transition probabilities})
có thể dùng để tính toàn bộ phân phối của quá trình bằng quy nạp:
\begin{align}
  \PP(Z_n = i_n, Z_{n-1} = i_{n-1}, \ldots, Z_0 = i_0)
  &= P_{i_{n-1}, i_n} \cdots P_{i_0, i_1} \, \PP(Z_0 = i_0),
\end{align}
với $i_0, i_1, \ldots, i_n \in \Z$.

Theo \term{quy tắc xác suất toàn phần} (\en{law of total probability}):
\begin{equation}\label{eq:total-prob}
  \PP(Z_1 = i) = \sum_{j \in \Z} \PP(Z_1 = i \mid Z_0 = j)\,\PP(Z_0 = j), \qquad i \in \Z.
\end{equation}

% --- Example: Random walk ---
\begin{example}[Bước đi ngẫu nhiên -- \en{Random walk}]
Bước đi ngẫu nhiên (\en{random walk}):
\begin{equation}\label{eq:random-walk}
  S_n = X_1 + \cdots + X_n, \qquad n \in \N,
\end{equation}
trong đó $(X_n)_{n \geq 1}$ là dãy các biến ngẫu nhiên \term{độc lập} (\en{independent})
nhận giá trị trong $\Z$, là một chuỗi Markov thời gian rời rạc.

\textbf{Chứng minh:} Ta cần chỉ ra tính Markov. Với mọi $j, i_n, \ldots, i_1 \in \Z$:
\begin{align*}
  \PP(S_{n+1} = j \mid S_n = i_n,\; S_{n-1} = i_{n-1},\; \ldots,\; S_1 = i_1)
  &= \PP(S_n + X_{n+1} = j \mid S_n = i_n, \ldots) \\
  &= \PP(X_{n+1} = j - i_n \mid S_n = i_n, \ldots) \\
  &= \PP(X_{n+1} = j - i_n) \\
  &= \PP(S_{n+1} = j \mid S_n = i_n).
\end{align*}
Bước thứ ba dùng tính \term{độc lập} của $X_{n+1}$ với $S_1, \ldots, S_n$ (vì $X_{n+1}$ độc lập với $X_1, \ldots, X_n$).
\end{example}

Tổng quát hơn, mọi quá trình có \term{gia số độc lập} (\en{independent increments}) đều là chuỗi Markov. Tuy nhiên, không phải mọi chuỗi Markov đều có gia số độc lập.

% ============================================================================
\subsection{Ma trận chuyển (Transition Matrix)}
\label{subsec:transition-matrix}

\textbf{Câu chuyện:} Ma trận chuyển là ``bản đồ xác suất'': từ mỗi thành phố (trạng thái),
bạn biết xác suất đi đến các thành phố khác. Hàng $i$ trong ma trận cho bạn biết:
``Nếu đang ở thành phố $i$, xác suất đi đến mỗi thành phố khác là bao nhiêu?''

Sự tiến hóa ngẫu nhiên của quá trình Markov $(Z_n)_{n\in\N}$ được xác định bởi dữ liệu:
\[
  P_{i,j} := \PP(Z_{n+1} = j \mid Z_n = i), \qquad i, j \in \Z,
\]
trong đó ta giả sử xác suất $\PP(Z_{n+1} = j \mid Z_n = i)$ không phụ thuộc vào $n \in \N$.
Khi đó chuỗi Markov $(Z_n)_{n\in\N}$ được gọi là \term{đồng nhất theo thời gian} (\en{time homogeneous}).

\begin{dinhnghia}[(Ma trận chuyển -- \en{Transition matrix})]
Dữ liệu trên được mã hóa thành một ma trận, gọi là \term{ma trận chuyển} (\en{transition matrix}) của chuỗi Markov:
\[
  [P_{i,j}]_{i,j\in\Z} = [\PP(Z_{n+1} = j \mid Z_n = i)]_{i,j\in\Z}.
\]
\end{dinhnghia}

\textbf{Lưu ý quan trọng:} Thứ tự các chỉ số $(i,j)$ trong $\PP(Z_{n+1} = j \mid Z_n = i)$ và $P_{i,j}$ bị đảo ngược. Cụ thể, trạng thái ban đầu $i$ là \emph{số hàng} trong ma trận, còn trạng thái đích $j$ là \emph{số cột}.

\begin{dinhnghia}[(Trạng thái hấp thụ -- \en{Absorbing state})]
Trạng thái $k \in \Z$ được gọi là \term{hấp thụ} (\en{absorbing}) nếu $P_{k,k} = 1$.
\end{dinhnghia}

Do quan hệ:
\begin{equation}\label{eq:row-sum-1}
  \sum_{j \in \Z} \PP(Z_{n+1} = j \mid Z_n = i) = 1, \qquad i \in \N,
\end{equation}
các \emph{hàng} của ma trận chuyển thỏa mãn: $\sum_{j} P_{i,j} = 1$ với mọi $i$.

% --- Finite state space ---
Trong trường hợp chuỗi Markov $(Z_k)_{k\in\N}$ nhận giá trị trong không gian trạng thái hữu hạn $\{0, \ldots, N\}$, ma trận chuyển có dạng:
\[
  [P_{i,j}]_{0 \leq i,j \leq N}
  = \begin{bmatrix}
    P_{0,0} & P_{0,1} & P_{0,2} & \cdots & P_{0,N} \\
    P_{1,0} & P_{1,1} & P_{1,2} & \cdots & P_{1,N} \\
    P_{2,0} & P_{2,1} & P_{2,2} & \cdots & P_{2,N} \\
    \vdots  & \vdots  & \vdots  & \ddots & \vdots \\
    P_{N,0} & P_{N,1} & P_{N,2} & \cdots & P_{N,N}
  \end{bmatrix}.
\]

% --- Examples ---
\begin{example}[Quá trình cờ bạc -- \en{Gambling process}]
\label{ex:gambler}
Một người chơi bắt đầu với $i$ đô la. Mỗi ván, anh ta thắng $1$ đô la (xác suất $p$)
hoặc thua $1$ đô la (xác suất $q = 1-p$). Trò chơi dừng khi cháy túi ($0$ đô la) hoặc đạt mục tiêu ($N$ đô la).

Ma trận chuyển trên $\{0, 1, \ldots, N\}$ với các trạng thái hấp thụ $0$ và $N$:
\[
  P = [P_{i,j}]_{0 \leq i,j \leq N}
  = \begin{bmatrix}
    1   & 0 & 0 & \cdots & 0 & 0 & 0 \\
    q   & 0 & p & \cdots & 0 & 0 & 0 \\
    \vdots & \vdots & \vdots & \ddots & \vdots & \vdots & \vdots \\
    0   & 0 & 0 & \cdots & q & 0 & p \\
    0   & 0 & 0 & \cdots & 0 & 0 & 1
  \end{bmatrix}.
\]
\end{example}

\begin{example}[Chuỗi Ehrenfest -- \en{Ehrenfest chain}]
\label{ex:ehrenfest}
Hai bình khí (trái và phải) nối với nhau bằng một lỗ, chứa tổng cộng $N$ viên bi.
Tại mỗi bước, ta chọn ngẫu nhiên một viên bi và chuyển nó sang bình kia.
Gọi $X_n \in \{0,\ldots,N\}$ là số bi bên trái tại thời điểm $n$.
Xác suất chuyển:
\[
  \PP(X_{n+1} = k+1 \mid X_n = k) = \frac{N-k}{N}, \qquad
  \PP(X_{n+1} = k-1 \mid X_n = k) = \frac{k}{N}.
\]

\textbf{Ý nghĩa vật lý:} Mô hình này mô phỏng sự khuếch tán khí giữa hai ngăn.
Nếu bên trái có nhiều bi hơn ($k > N/2$), xác suất giảm ($k/N$) lớn hơn xác suất tăng ($(N-k)/N$)
--- hệ thống ``kéo'' về trạng thái cân bằng $k = N/2$.
\end{example}

\begin{example}[Ehrenfest $N = 4$ -- ma trận cụ thể]
\label{ex:ehrenfest-N4}
Với $N = 4$ viên bi, ma trận chuyển $5 \times 5$:
\[
  P = \begin{bmatrix}
    0 & 1 & 0 & 0 & 0 \\
    1/4 & 0 & 3/4 & 0 & 0 \\
    0 & 2/4 & 0 & 2/4 & 0 \\
    0 & 0 & 3/4 & 0 & 1/4 \\
    0 & 0 & 0 & 1 & 0
  \end{bmatrix}.
\]

\textbf{Đọc ma trận:}
\begin{itemize}[nosep]
  \item Hàng 0: nếu bên trái có 0 bi, chắc chắn tăng lên 1 (chọn bi nào cũng từ bên phải).
  \item Hàng 2: nếu bên trái có 2 bi, xác suất giảm = $2/4$, xác suất tăng = $2/4$ (cân bằng!).
  \item Hàng 4: nếu bên trái có 4 bi (tất cả), chắc chắn giảm xuống 3.
\end{itemize}
\end{example}

\begin{example}[Xếp hạng tín dụng -- \en{Credit rating}]
Trong tài chính, xếp hạng tín dụng (AAA, AA, A, BBB, \ldots, D) có thể được mô hình hóa
bằng chuỗi Markov. Ma trận chuyển ghi lại xác suất thay đổi hạng tín dụng sau mỗi năm.
Các hạng cao hơn có hệ số đường chéo lớn hơn (ổn định hơn).
\end{example}

% ============================================================================
\subsection{Xác suất chuyển nhiều bước (Multi-step Transition Probabilities)}

Ta muốn tính \term{xác suất chuyển hai bước} (\en{two-step transition probability})
$\PP(Z_{n+2} = j \mid Z_n = i)$, không có trực tiếp trong ma trận chuyển.

Bằng \term{phân tích bước đầu tiên} (\en{first step analysis}), ký hiệu $S$ là không gian trạng thái:
\begin{align*}
  \PP(Z_{n+2} = j \mid Z_n = i)
  &= \sum_{l \in S} \PP(Z_{n+1} = l \mid Z_n = i)\,\PP(Z_{n+2} = j \mid Z_{n+1} = l)\\
  &= \sum_{l \in S} P_{i,l}\,P_{l,j}
  = P^2_{i,j}, \qquad i, j \in S.
\end{align*}

Tổng quát hơn, với mọi $k \in \N$:

\begin{congthuc}
\textbf{Xác suất chuyển $k$ bước:}
\[
  [\PP(Z_{n+k} = j \mid Z_n = i)]_{i,j\in S} = [P^k_{i,j}]_{i,j\in S} = P^k.
\]
Nghĩa là: xác suất chuyển $k$ bước chính là phần tử tương ứng của lũy thừa bậc $k$ của ma trận chuyển.
\end{congthuc}

Quan hệ $P^{m+n} = P^m P^n$ chính là \term{phương trình Chapman--Kolmogorov} (\en{Chapman--Kolmogorov equation}):
\[
  P^{m+n}_{i,j} = \sum_{l \in S} P^m_{i,l}\,P^n_{l,j}, \qquad i,j \in S, \quad m, n \in \N.
\]

\begin{example}[Tính $P^2$ cho chuỗi 3 trạng thái]
\label{ex:P2-three-state}
Xét ma trận chuyển trên $S = \{0, 1, 2\}$:
\[
  P = \begin{bmatrix}
    0 & 1/2 & 1/2 \\
    1/3 & 1/3 & 1/3 \\
    1/2 & 0 & 1/2
  \end{bmatrix}.
\]

Tính $P^2 = P \cdot P$:
\begin{align*}
  P^2_{0,0} &= 0 \cdot 0 + \frac{1}{2} \cdot \frac{1}{3} + \frac{1}{2} \cdot \frac{1}{2} = \frac{1}{6} + \frac{1}{4} = \frac{5}{12}, \\
  P^2_{0,1} &= 0 \cdot \frac{1}{2} + \frac{1}{2} \cdot \frac{1}{3} + \frac{1}{2} \cdot 0 = \frac{1}{6}, \\
  P^2_{0,2} &= 0 \cdot \frac{1}{2} + \frac{1}{2} \cdot \frac{1}{3} + \frac{1}{2} \cdot \frac{1}{2} = \frac{5}{12}.
\end{align*}

Kiểm tra: $P^2_{0,0} + P^2_{0,1} + P^2_{0,2} = 5/12 + 1/6 + 5/12 = 5/12 + 2/12 + 5/12 = 12/12 = 1$. Đúng!

\textbf{Ý nghĩa:} Xuất phát từ trạng thái 0, sau 2 bước, xác suất ở trạng thái 0 là $5/12 \approx 0.42$.
\end{example}

\begin{example}[Tính $P^n$ cho chuỗi Ehrenfest $N = 4$]
\label{ex:Pn-ehrenfest}
Với ma trận $P$ ở Ví dụ~\ref{ex:ehrenfest-N4}, ta tính $P^2$:
\[
  P^2_{0,0} = 0 \cdot 0 + 1 \cdot \frac{1}{4} + 0 + 0 + 0 = \frac{1}{4}.
\]
Nghĩa là: nếu bên trái có 0 bi, sau 2 bước xác suất trở về 0 bi là $1/4$.

Thật vậy: từ 0 phải đi lên 1 (chắc chắn), rồi từ 1 quay về 0 (xác suất $1/4$).
\end{example}

\begin{tomtat}
Các xác suất chuyển $k$ bước được tính bằng cách lũy thừa ma trận chuyển: $P^k$.
Phương trình Chapman--Kolmogorov cho phép phân rã quá trình chuyển thành hai giai đoạn.

\textbf{Mẹo tính toán:} Đối với ma trận $n \times n$ lớn, dùng chéo hóa $P = M D M^{-1}$
sẽ cho $P^k = M D^k M^{-1}$ --- dễ tính hơn nhiều so với nhân ma trận $k$ lần.
\end{tomtat}

% ============================================================================
\section{Chuỗi Markov hai trạng thái (The Two-State Markov Chain)}
\label{sec:two-state}

\textbf{Câu chuyện:} Mô hình đơn giản nhất: thời tiết ngày mai chỉ có ``Nắng'' (0) hoặc ``Mưa'' (1).
Nếu hôm nay nắng, ngày mai mưa với xác suất $a$; nếu hôm nay mưa, ngày mai nắng với xác suất $b$.
Sau nhiều ngày, tỷ lệ ngày nắng/mưa sẽ ổn định --- đây là phân phối giới hạn.

Việc tính ma trận chuyển bậc $n$, tức $P^n$, nói chung là khó. Tuy nhiên, khi số trạng thái bằng 2, tức $S = \{0, 1\}$, ta có thể tính chính xác.

Ma trận chuyển có dạng:
\begin{equation}\label{eq:two-state}
  P = \begin{bmatrix} 1-a & a \\ b & 1-b \end{bmatrix},
\end{equation}
với $a \in [0,1]$ và $b \in [0,1]$.

\textbf{Ý nghĩa:}
\begin{itemize}[nosep]
  \item $\PP(Z_{n+1} = 1 \mid Z_n = 0) = a$: xác suất chuyển từ $0$ sang $1$.
  \item $\PP(Z_{n+1} = 0 \mid Z_n = 0) = 1-a$: xác suất ở lại trạng thái $0$.
  \item $\PP(Z_{n+1} = 0 \mid Z_n = 1) = b$: xác suất chuyển từ $1$ sang $0$.
  \item $\PP(Z_{n+1} = 1 \mid Z_n = 1) = 1-b$: xác suất ở lại trạng thái $1$.
\end{itemize}

Ta loại trừ trường hợp $a = b = 0$ (ma trận đơn vị, chuỗi hằng).

\subsection{Chéo hóa ma trận (Matrix Diagonalization)}

Ma trận $P$ có hai \term{trị riêng} (\en{eigenvalues}):
\[
  \lambda_1 = 1 \qquad \text{và} \qquad \lambda_2 = 1 - a - b,
\]
với các \term{vectơ riêng} (\en{eigenvectors}) tương ứng:
\[
  \begin{bmatrix} 1 \\ 1 \end{bmatrix}
  \qquad \text{và} \qquad
  \begin{bmatrix} -a \\ b \end{bmatrix}.
\]

Do đó $P$ có thể \term{chéo hóa} (\en{diagonalize}): $P = M D M^{-1}$, và:
\begin{equation}\label{eq:Pn-two-state}
  P^n = M D^n M^{-1}
  = \frac{1}{a+b}
  \begin{bmatrix}
    b + a\lambda_2^n & a(1 - \lambda_2^n) \\
    b(1-\lambda_2^n) & a + b\lambda_2^n
  \end{bmatrix}.
\end{equation}

\textbf{Chi tiết tính toán:}
\[
  M = \begin{bmatrix} 1 & -a \\ 1 & b \end{bmatrix}, \quad
  D = \begin{bmatrix} 1 & 0 \\ 0 & 1-a-b \end{bmatrix}, \quad
  M^{-1} = \frac{1}{a+b}\begin{bmatrix} b & a \\ -1 & 1 \end{bmatrix}.
\]

\subsection{Phân phối giới hạn (Limiting Distribution)}

Cho $n \to \infty$ trong \eqref{eq:Pn-two-state}, vì $|\lambda_2| = |1-a-b| < 1$ khi $(a,b) \neq (1,1)$:

\begin{dinhly}[(Phân phối giới hạn của chuỗi hai trạng thái)]
\[
  \lim_{n\to\infty} P^n = \frac{1}{a+b}
  \begin{bmatrix} b & a \\ b & a \end{bmatrix},
\]
do đó:
\begin{align}
  \lim_{n\to\infty} \PP(Z_n = 1 \mid Z_0 = 0) &= \lim_{n\to\infty} \PP(Z_n = 1 \mid Z_0 = 1) = \frac{a}{a+b}, \label{eq:limit-1}\\
  \lim_{n\to\infty} \PP(Z_n = 0 \mid Z_0 = 0) &= \lim_{n\to\infty} \PP(Z_n = 0 \mid Z_0 = 1) = \frac{b}{a+b}. \label{eq:limit-0}
\end{align}
\end{dinhly}

\term{Phân phối giới hạn} (\en{limiting distribution}):
\begin{equation}\label{eq:pi-two-state}
  \pi = [\pi_0, \pi_1] = \left[\frac{b}{a+b},\; \frac{a}{a+b}\right].
\end{equation}

\begin{tomtat}
Phân phối giới hạn $\pi$ \textbf{không phụ thuộc vào trạng thái ban đầu} $Z_0$.
Sau thời gian đủ lớn, xác suất ở trạng thái $1$ xấp xỉ $a/(a+b)$,
và ở trạng thái $0$ xấp xỉ $b/(a+b)$.

Phân phối $\pi$ là \term{bất biến} (\en{invariant}) theo $P$, tức $\pi P = \pi$.

\textbf{Ví dụ thời tiết:} Nếu $a = 0.3$ (nắng$\to$mưa) và $b = 0.4$ (mưa$\to$nắng),
thì dài hạn: $\pi_{\text{nắng}} = 0.4/0.7 \approx 57\%$, $\pi_{\text{mưa}} = 0.3/0.7 \approx 43\%$.
\end{tomtat}

\textbf{Trường hợp đặc biệt:} Khi $a + b = 1$ (hay $a = 1-b$):
\[
  P^n = \begin{bmatrix} b & a \\ b & a \end{bmatrix} = P, \qquad n \in \N,
\]
và $\PP(Z_n = 1 \mid Z_0 = 0) = \PP(Z_n = 1 \mid Z_0 = 1) = a$ với mọi $n \geq 1$.

\textbf{Trường hợp $a = b = 1$:} Giới hạn $\lim_{n\to\infty} P^n$ không tồn tại vì:
\[
  P^n = \begin{cases}
    \begin{bmatrix} 1 & 0 \\ 0 & 1 \end{bmatrix}, & n = 2k, \\[6pt]
    \begin{bmatrix} 0 & 1 \\ 1 & 0 \end{bmatrix}, & n = 2k+1.
  \end{cases}
\]
Chuỗi ``dao động'' mãi mãi --- đây là chuỗi \term{tuần hoàn} với chu kỳ $d = 2$.

\subsection{Tốc độ hội tụ (Rate of Convergence)}

Từ \eqref{eq:Pn-two-state}, sai số giữa $P^n_{i,j}$ và $\pi_j$ là:
\[
  |P^n_{i,j} - \pi_j| \leq C \cdot |\lambda_2|^n = C \cdot |1 - a - b|^n.
\]

Sai số giảm theo cấp số nhân. Hội tụ nhanh khi $|1-a-b|$ nhỏ, tức $a + b$ gần $1$.

\begin{example}[So sánh tốc độ hội tụ]
\begin{itemize}[nosep]
  \item $a = 0.5, b = 0.5$: $|\lambda_2| = |1-1| = 0$. Hội tụ ngay sau 1 bước!
  \item $a = 0.1, b = 0.1$: $|\lambda_2| = 0.8$. Hội tụ chậm (cần $\sim 30$ bước).
  \item $a = 0.01, b = 0.01$: $|\lambda_2| = 0.98$. Rất chậm (cần $\sim 300$ bước).
\end{itemize}
Trực quan: khi $a, b$ nhỏ, chuỗi ``dính'' ở mỗi trạng thái lâu, nên mất nhiều thời gian mới ``quên'' trạng thái ban đầu.
\end{example}

% ============================================================================
\section{Bài tập (Exercises)}
\label{sec:ch4-exercises}

\begin{exercise}
Trong mô hình đơn giản hóa của một chương trình truyền hình, người chơi đã thắng $k$ đô la,
sẽ có $k+1$ đô la ở lượt tiếp theo với xác suất $q$, hoặc bị loại (về $0$) với xác suất $p = 1-q$.
Giả sử người chơi bắt đầu với $1$ đô la. Hãy mô hình hóa tiền thắng sau $n$ lượt chơi
bằng chuỗi Markov và viết ma trận chuyển tương ứng.
\end{exercise}

\begin{exercise}
Xét chuỗi Markov trên $S = \{0, 1, 2\}$ với ma trận chuyển:
\[
  P = \begin{bmatrix}
    1/2 & 1/4 & 1/4 \\
    1/3 & 1/3 & 1/3 \\
    1/4 & 1/4 & 1/2
  \end{bmatrix}.
\]
\begin{enumerate}[label=(\alph*)]
  \item Tính $P^2$ (nhân ma trận).
  \item Tính $\PP(Z_2 = 0 \mid Z_0 = 1)$.
  \item Tính $\PP(Z_3 = 2 \mid Z_0 = 0)$ (gợi ý: dùng $P^3 = P^2 \cdot P$).
\end{enumerate}
\end{exercise}

\begin{exercise}
Cho chuỗi Ehrenfest với $N = 4$ viên bi (Ví dụ~\ref{ex:ehrenfest-N4}).
\begin{enumerate}[label=(\alph*)]
  \item Bắt đầu với 4 bi bên trái. Tính xác suất có 2 bi bên trái sau 2 bước.
  \item Chuỗi này có trạng thái hấp thụ không? Giải thích.
  \item Tìm phân phối dừng $\pi$ (gợi ý: so sánh với phân phối nhị thức $B(4, 1/2)$).
\end{enumerate}
\end{exercise}

\begin{exercise}[Chapman--Kolmogorov trực tiếp]
Xét chuỗi Markov hai trạng thái với $a = 0.4$, $b = 0.6$.
\begin{enumerate}[label=(\alph*)]
  \item Tính $P^2$ trực tiếp bằng nhân ma trận.
  \item Tính $P^2$ bằng công thức \eqref{eq:Pn-two-state} với $n = 2$. Xác minh kết quả trùng nhau.
  \item Tìm $n$ nhỏ nhất sao cho $|P^n_{0,1} - \pi_1| < 0.01$.
\end{enumerate}
\end{exercise}

\begin{exercise}
Chứng minh rằng nếu $(S_n)_{n\geq 0}$ là bước đi ngẫu nhiên với $S_n = X_1 + \cdots + X_n$
và $(X_k)$ độc lập cùng phân phối, thì $(S_n)$ thỏa mãn tính chất Markov.
(Gợi ý: sử dụng tính độc lập của $X_{n+1}$ với $(X_1, \ldots, X_n)$.)
\end{exercise}

\chapter{Phân tích bước đầu tiên\\{\normalsize\textit{First Step Analysis}}}
\label{ch:first-step}

\textbf{Câu chuyện:} Hãy tưởng tượng bạn đang ở tầng 3 của một tòa nhà 5 tầng.
Mỗi phút, bạn ngẫu nhiên đi lên hoặc xuống 1 tầng. Câu hỏi:
``Xác suất bạn đến tầng trệt (tầng 0) trước khi lên mái nhà (tầng 5)?''
và ``Trung bình mất bao lâu để bạn đến một trong hai?''

Kỹ thuật \term{phân tích bước đầu tiên} trả lời cả hai câu hỏi bằng cách:
``Nhìn bước đầu tiên --- bạn đi lên hay xuống? --- rồi từ vị trí mới, bài toán giống hệt nhưng nhỏ hơn.''

\begin{tomtat}
\textbf{Ý tưởng cốt lõi:} Điều kiện hóa theo bước đầu tiên $Z_1$, rồi dùng tính Markov
để quy về bài toán cùng dạng. Kết quả: hệ phương trình tuyến tính --- giải là xong!
\end{tomtat}

% ============================================================================
\section{Xác suất chạm (Hitting Probabilities)}
\label{sec:hitting-prob}

Cho chuỗi Markov $(Z_n)_{n \in \N}$ nhận giá trị trong không gian trạng thái hữu hạn
$S = \{0, 1, \ldots, N\}$ với ma trận chuyển $P = [P_{k,l}]_{k,l \in S}$.

\begin{dinhnghia}[(Thời gian chạm -- \en{Hitting time})]
Cho $A \subseteq S$. \term{Thời gian chạm} (\en{hitting time}) tập $A$ được định nghĩa là:
\[
  T_A := \inf\{n \geq 0 : Z_n \in A\},
\]
tức là thời điểm đầu tiên chuỗi Markov rơi vào tập $A$.
\end{dinhnghia}

\begin{dinhnghia}[(Xác suất chạm -- \en{Hitting probability})]
\term{Xác suất chạm} (\en{hitting probability}) trạng thái $l \in A$ xuất phát từ $k \in S$ là:
\[
  g_l(k) := \PP(Z_{T_A} = l \mid Z_0 = k), \qquad k \in S, \quad l \in A.
\]
Đây là xác suất mà chuỗi Markov sẽ đến trạng thái $l$ khi lần đầu tiên đi vào tập $A$,
biết rằng ban đầu chuỗi ở trạng thái $k$.
\end{dinhnghia}

\begin{dinhnghia}[(Thời gian chạm trung bình -- \en{Mean hitting time})]
\term{Thời gian chạm trung bình} (\en{mean hitting time}) tập $A$ xuất phát từ $k$ là:
\[
  h_A(k) := \E[T_A \mid Z_0 = k], \qquad k \in S.
\]
\end{dinhnghia}

% --- First step analysis derivation ---
\subsection{Phương pháp phân tích bước đầu tiên (First Step Analysis Method)}

Ý tưởng chính của \term{phân tích bước đầu tiên} (\en{first step analysis}) là điều kiện hóa
theo bước đầu tiên $Z_1$ của chuỗi Markov. Từ tính Markov và quy tắc xác suất toàn phần,
ta có:
\begin{align*}
  g_l(k) &= \PP(Z_{T_A} = l \mid Z_0 = k)\\
  &= \sum_{m \in S} \PP(Z_{T_A} = l \mid Z_1 = m,\; Z_0 = k)\,\PP(Z_1 = m \mid Z_0 = k)\\
  &= \sum_{m \in S} P_{k,m}\, g_l(m), \qquad k \in S \setminus A.
\end{align*}

Với \term{điều kiện biên} (\en{boundary condition}): nếu $k \in A$ thì chuỗi đã ở trong $A$
ngay tại $n = 0$, do đó:
\[
  g_l(k) = \ind_{\{k = l\}}, \qquad k \in A, \quad l \in A.
\]

\begin{congthuc}
\textbf{Hệ phương trình phân tích bước đầu tiên cho xác suất chạm:}
\[
  \boxed{
    g_l(k) = \sum_{m \in S} P_{k,m}\, g_l(m), \qquad k \in S \setminus A,
  }
\]
với điều kiện biên $g_l(k) = \ind_{\{k = l\}}$ khi $k \in A$.
\end{congthuc}

\begin{tomtat}
Phân tích bước đầu tiên biến bài toán tính xác suất chạm thành một hệ phương trình tuyến tính.
Ta ``mở ra'' bước đầu tiên $Z_0 \to Z_1$, sau đó sử dụng tính Markov để quy bài toán
về cùng dạng nhưng xuất phát từ trạng thái $Z_1$.
\end{tomtat}

% --- Absorbing set and matrix form ---
\subsection{Tập hấp thụ và dạng ma trận (Absorbing Set and Matrix Form)}

Ta nói $A$ là \term{tập hấp thụ} (\en{absorbing set}) nếu chuỗi Markov không thể rời khỏi $A$:
\[
  P_{k,l} = \ind_{\{k = l\}}, \qquad k, l \in A.
\]
Nghĩa là mọi trạng thái trong $A$ đều là \term{trạng thái hấp thụ} (\en{absorbing state}).

Khi $A = \{r+1, r+2, \ldots, N\}$ là tập hấp thụ, ta có thể sắp xếp ma trận chuyển thành dạng khối:
\begin{equation}\label{eq:block-matrix}
  P = \begin{bmatrix} Q & R \\ \mathbf{0} & I \end{bmatrix},
\end{equation}
trong đó:
\begin{itemize}[nosep]
  \item $Q$ là ma trận $(r+1) \times (r+1)$ chứa xác suất chuyển giữa các trạng thái \emph{không hấp thụ}
        (\en{transient states}) $S \setminus A = \{0, 1, \ldots, r\}$.
  \item $R$ là ma trận $(r+1) \times (N-r)$ chứa xác suất chuyển từ trạng thái không hấp thụ sang trạng thái hấp thụ.
  \item $I$ là ma trận đơn vị $(N-r) \times (N-r)$.
  \item $\mathbf{0}$ là ma trận không $(N-r) \times (r+1)$.
\end{itemize}

Trong dạng ma trận, hệ phương trình cho xác suất chạm trở thành:

\begin{congthuc}
\textbf{Dạng ma trận:}
\[
  \mathbf{g}_l = Q\, \mathbf{g}_l + R\, \mathbf{e}_l,
\]
trong đó $\mathbf{g}_l = [g_l(0), g_l(1), \ldots, g_l(r)]^\top$ và $\mathbf{e}_l$ là vectơ đơn vị
tương ứng với trạng thái $l$ trong $A$. Do đó:
\[
  \boxed{
    \mathbf{g}_l = (I - Q)^{-1} R\, \mathbf{e}_l.
  }
\]
\end{congthuc}

Ngoài ra, tổng xác suất chạm tất cả các trạng thái trong $A$ thỏa mãn:
\begin{equation}\label{eq:total-hitting}
  1 = \PP(T_A = +\infty \mid Z_0 = k) + \sum_{l \in A} g_l(k), \qquad k \in S.
\end{equation}
Khi $A$ là tập hấp thụ và chuỗi chắc chắn bị hấp thụ, thì $\PP(T_A = +\infty \mid Z_0 = k) = 0$
và $\sum_{l \in A} g_l(k) = 1$ với mọi $k \in S$.

% --- Example 1: 4-state chain ---
\begin{example}[Chuỗi 4 trạng thái -- \en{4-state chain}]
\label{ex:4state-hitting}
Xét chuỗi Markov trên $S = \{0, 1, 2, 3\}$ với các trạng thái hấp thụ $A = \{0, 3\}$.
Ma trận chuyển có dạng:
\[
  P = \begin{bmatrix}
    1 & 0 & 0 & 0 \\
    P_{1,0} & P_{1,1} & P_{1,2} & P_{1,3} \\
    P_{2,0} & P_{2,1} & P_{2,2} & P_{2,3} \\
    0 & 0 & 0 & 1
  \end{bmatrix}.
\]

Ta cần tính $g_0(1)$, $g_0(2)$: xác suất bị hấp thụ vào trạng thái $0$, xuất phát từ $1$ hoặc $2$.

Áp dụng phân tích bước đầu tiên:
\[
  \begin{cases}
    g_0(1) = P_{1,0} \cdot 1 + P_{1,1}\, g_0(1) + P_{1,2}\, g_0(2) + P_{1,3} \cdot 0,\\
    g_0(2) = P_{2,0} \cdot 1 + P_{2,1}\, g_0(1) + P_{2,2}\, g_0(2) + P_{2,3} \cdot 0,
  \end{cases}
\]
với điều kiện biên $g_0(0) = 1$ và $g_0(3) = 0$.

Đây là hệ hai phương trình tuyến tính hai ẩn, có thể giải trực tiếp.

Ví dụ cụ thể, nếu:
\[
  P = \begin{bmatrix}
    1 & 0 & 0 & 0 \\
    1/4 & 1/4 & 1/4 & 1/4 \\
    1/4 & 1/4 & 1/4 & 1/4 \\
    0 & 0 & 0 & 1
  \end{bmatrix},
\]
thì hệ trở thành:
\[
  \begin{cases}
    g_0(1) = \frac{1}{4} + \frac{1}{4}\, g_0(1) + \frac{1}{4}\, g_0(2),\\
    g_0(2) = \frac{1}{4} + \frac{1}{4}\, g_0(1) + \frac{1}{4}\, g_0(2).
  \end{cases}
\]
Từ tính đối xứng $g_0(1) = g_0(2)$, ta suy ra:
\[
  g_0(1) = \frac{1}{4} + \frac{1}{2}\, g_0(1)
  \quad\Longrightarrow\quad
  g_0(1) = g_0(2) = \frac{1}{2}.
\]
\end{example}

% ============================================================================
\section{Thời gian chạm và hấp thụ trung bình (Mean Hitting and Absorption Times)}
\label{sec:mean-hitting}

\begin{dinhnghia}[(Thời gian chạm trung bình -- \en{Mean hitting time})]
Cho $A \subseteq S$. \term{Thời gian chạm trung bình} tập $A$ xuất phát từ $k$ là:
\[
  h_A(k) := \E[T_A \mid Z_0 = k], \qquad k \in S.
\]
\end{dinhnghia}

Áp dụng phân tích bước đầu tiên tương tự như phần trước, nhưng lần này ta điều kiện hóa
kỳ vọng thay vì xác suất:
\begin{align*}
  h_A(k) &= \E[T_A \mid Z_0 = k]\\
  &= \E[1 + T_A \circ \theta_1 \mid Z_0 = k]\\
  &= 1 + \sum_{l \in S} P_{k,l}\, h_A(l), \qquad k \in S \setminus A,
\end{align*}
trong đó $\theta_1$ là toán tử dịch chuyển thời gian (\en{time shift operator}).

Điều kiện biên: $h_A(l) = 0$ với mọi $l \in A$ (nếu đã ở trong $A$ thì không cần thời gian nào).

\begin{congthuc}
\textbf{Hệ phương trình phân tích bước đầu tiên cho thời gian chạm trung bình:}
\[
  \boxed{
    h_A(k) = 1 + \sum_{l \in S} P_{k,l}\, h_A(l), \qquad k \in S \setminus A,
  }
\]
với điều kiện biên $h_A(l) = 0$ khi $l \in A$.
\end{congthuc}

Trong dạng ma trận, với $\mathbf{h}_A = [h_A(0), \ldots, h_A(r)]^\top$ chỉ chứa các trạng thái
không hấp thụ:
\begin{equation}\label{eq:mean-hitting-matrix}
  \mathbf{h}_A = \mathbf{1} + Q\, \mathbf{h}_A
  \quad\Longrightarrow\quad
  \boxed{
    \mathbf{h}_A = (I - Q)^{-1}\, \mathbf{1},
  }
\end{equation}
trong đó $\mathbf{1}$ là vectơ toàn $1$.

\begin{tomtat}
Thời gian chạm trung bình thỏa mãn hệ phương trình tuyến tính tương tự xác suất chạm,
chỉ khác ở vế phải: thêm $+1$ (đếm một bước) và điều kiện biên bằng $0$ tại $A$.
Ma trận $(I - Q)^{-1}$ được gọi là \term{ma trận cơ bản} (\en{fundamental matrix}) của chuỗi Markov hấp thụ.
\end{tomtat}

% --- Two-state example ---
\begin{example}[Chuỗi hai trạng thái -- \en{Two-state chain}]
\label{ex:two-state-hitting}
Xét chuỗi Markov hai trạng thái $S = \{0, 1\}$ với ma trận chuyển:
\[
  P = \begin{bmatrix} 1-a & a \\ b & 1-b \end{bmatrix}.
\]

\textbf{Tính $h_{\{0\}}(1)$:} Thời gian trung bình để đến trạng thái $0$ từ trạng thái $1$.

Áp dụng phân tích bước đầu tiên với $A = \{0\}$:
\[
  h_{\{0\}}(1) = 1 + P_{1,0}\, h_{\{0\}}(0) + P_{1,1}\, h_{\{0\}}(1)
  = 1 + b \cdot 0 + (1-b)\, h_{\{0\}}(1).
\]
Giải ra:
\[
  b\, h_{\{0\}}(1) = 1 \quad\Longrightarrow\quad h_{\{0\}}(1) = \frac{1}{b}.
\]

Tương tự, $h_{\{1\}}(0) = \dfrac{1}{a}$.
\end{example}

% --- Example 2: Chain with alpha, beta ---
\begin{example}[Chuỗi với tham số $\alpha$, $\beta$ -- \en{Chain with parameters $\alpha$, $\beta$}]
\label{ex:alpha-beta-hitting}
Xét chuỗi Markov trên $S = \{0, 1, 2, 3\}$ với trạng thái hấp thụ $A = \{3\}$ và ma trận chuyển:
\[
  P = \begin{bmatrix}
    1 - \beta & \beta & 0 & 0\\
    0 & 1 - \beta & \beta & 0\\
    0 & 0 & 1 - \beta & \beta\\
    0 & 0 & 0 & 1
  \end{bmatrix}.
\]
Ta cần tính thời gian trung bình để đến trạng thái $3$.

Hệ phương trình phân tích bước đầu tiên:
\[
  \begin{cases}
    h_3(0) = 1 + (1-\beta)\, h_3(0) + \beta\, h_3(1),\\
    h_3(1) = 1 + (1-\beta)\, h_3(1) + \beta\, h_3(2),\\
    h_3(2) = 1 + (1-\beta)\, h_3(2) + \beta \cdot 0.
  \end{cases}
\]

Từ phương trình thứ ba: $\beta\, h_3(2) = 1$, suy ra $h_3(2) = \dfrac{1}{\beta}$.

Thay vào phương trình thứ hai:
$\beta\, h_3(1) = 1 + \beta \cdot \dfrac{1}{\beta} = 2$, suy ra $h_3(1) = \dfrac{2}{\beta}$.

Thay vào phương trình thứ nhất:
$\beta\, h_3(0) = 1 + \beta \cdot \dfrac{2}{\beta} = 3$, suy ra $h_3(0) = \dfrac{3}{\beta}$.

Tổng quát hơn, với chuỗi ``sinh--tử'' (\en{birth chain}) tuyến tính đi từ $0$ đến $N$,
ta có $h_{\{N\}}(k) = \dfrac{N - k}{\beta}$.
\end{example}

% --- Generalization ---
\subsection{Tổng quát hóa (Generalization)}

Phương pháp phân tích bước đầu tiên có thể mở rộng để tính:
\[
  h_A(k) = \E\!\left[\sum_{i=0}^{T_A - 1} f(X_i) \;\middle|\; X_0 = k\right],
\]
trong đó $f : S \to \R$ là một hàm cho trước. Khi $f \equiv 1$, ta thu được thời gian chạm trung bình.

Hệ phương trình tương ứng:
\begin{equation}\label{eq:generalized-hitting}
  h_A(k) = f(k) + \sum_{l \in S} P_{k,l}\, h_A(l), \qquad k \in S \setminus A,
\end{equation}
với điều kiện biên $h_A(l) = 0$ khi $l \in A$.

Trường hợp đặc biệt khi $A = \{m\}$ chỉ gồm một trạng thái:
\[
  h_{\{m\}}(k) = 1 + \sum_{l \neq m} P_{k,l}\, h_{\{m\}}(l), \qquad k \neq m.
\]

% ============================================================================
\section{Thời gian quay lại lần đầu (First Return Times)}
\label{sec:first-return}

Có sự khác biệt quan trọng giữa \term{thời gian chạm} và \term{thời gian quay lại lần đầu}:
thời gian chạm cho phép $n = 0$ (tức nếu đã ở $j$ thì $T_{\{j\}} = 0$),
còn thời gian quay lại yêu cầu $n \geq 1$ (phải rời khỏi rồi quay lại).

\begin{dinhnghia}[(Thời gian quay lại lần đầu -- \en{First return time})]
\term{Thời gian quay lại lần đầu} trạng thái $j$ là:
\[
  \tau_j := \inf\{n \geq 1 : X_n = j\}.
\]
Lưu ý: $n \geq 1$, không phải $n \geq 0$.
\end{dinhnghia}

\begin{dinhnghia}[(Xác suất quay lại -- \en{Return probability})]
\term{Xác suất quay lại} trạng thái $j$ từ trạng thái $i$ là:
\[
  p_{ij} := \PP(\tau_j < \infty \mid X_0 = i).
\]
Khi $i = j$, đại lượng $p_{jj}$ là xác suất quay lại chính trạng thái $j$.
\end{dinhnghia}

\begin{dinhnghia}[(Phân phối thời gian quay lại -- \en{First return distribution})]
Xác suất quay lại trạng thái $j$ sau đúng $n$ bước là:
\[
  f_{ij}^{(n)} := \PP(\tau_j = n \mid X_0 = i), \qquad n \geq 1.
\]
\end{dinhnghia}

Tổng xác suất quay lại: $p_{ij} = \sum_{n=1}^{\infty} f_{ij}^{(n)}$.

\begin{dinhly}[(Quan hệ giữa $P_{i,i}^n$ và $f_{ii}^{(n)}$)]
\label{thm:Pn-fn-relation}
Xác suất chuyển $n$ bước quay lại trạng thái $i$ có thể phân tách theo thời gian quay lại lần đầu:
\begin{equation}\label{eq:Pn-fn}
  P_{i,i}^{(n)} = \sum_{k=1}^{n} f_{ii}^{(k)}\, P_{i,i}^{(n-k)}, \qquad n \geq 1,
\end{equation}
với quy ước $P_{i,i}^{(0)} = 1$.
\end{dinhly}

\textbf{Chứng minh.}
Ta phân hoạch sự kiện $\{X_n = i\}$ theo thời điểm quay lại lần đầu $\tau_i = k$:
\begin{align*}
  \PP(X_n = i \mid X_0 = i)
  &= \sum_{k=1}^{n} \PP(X_n = i \mid \tau_i = k,\; X_0 = i)\,\PP(\tau_i = k \mid X_0 = i)\\
  &= \sum_{k=1}^{n} P_{i,i}^{(n-k)}\, f_{ii}^{(k)}.
\end{align*}
\hfill$\square$

% --- Mean return time ---
\subsection{Thời gian quay lại trung bình (Mean Return Time)}

\begin{dinhnghia}[(Thời gian quay lại trung bình -- \en{Mean return time})]
\[
  \mu_j(i) := \E[\tau_j \mid X_0 = i].
\]
Khi $i = j$, ta ký hiệu $\mu_j := \mu_j(j) = \E[\tau_j \mid X_0 = j]$.
\end{dinhnghia}

Áp dụng phân tích bước đầu tiên:

\begin{congthuc}
\textbf{Thời gian quay lại trung bình:}
\[
  \boxed{
    \mu_j(i) = 1 + \sum_{l \neq j} P_{i,l}\, \mu_j(l), \qquad i \in S.
  }
\]
\end{congthuc}

\begin{example}[Chuỗi hai trạng thái -- thời gian quay lại]
\label{ex:two-state-return}
Với chuỗi hai trạng thái $S = \{0,1\}$ và ma trận chuyển
$P = \begin{bsmallmatrix} 1-a & a \\ b & 1-b \end{bsmallmatrix}$:

\textbf{Tính $\mu_0(0)$:} Thời gian quay lại trung bình trạng thái $0$, xuất phát từ $0$.
\[
  \mu_0(0) = 1 + P_{0,1}\, \mu_0(1) = 1 + a\, \mu_0(1).
\]
Tiếp tục:
\[
  \mu_0(1) = 1 + P_{1,1}\, \mu_0(1) = 1 + (1-b)\, \mu_0(1)
  \quad\Longrightarrow\quad
  \mu_0(1) = \frac{1}{b}.
\]
Do đó:
\[
  \mu_0(0) = 1 + \frac{a}{b} = \frac{a + b}{b}.
\]

Tương tự:
\[
  \mu_1(1) = \frac{a + b}{a}, \qquad
  \mu_1(0) = \frac{1}{a}, \qquad
  \mu_0(1) = \frac{1}{b}.
\]

\begin{tomtat}
Lưu ý mối liên hệ với phân phối giới hạn (Chương~\ref{ch:discrete}):
\[
  \pi_0 = \frac{b}{a+b} = \frac{1}{\mu_0(0)}, \qquad
  \pi_1 = \frac{a}{a+b} = \frac{1}{\mu_1(1)}.
\]
Đây không phải sự trùng hợp! Với chuỗi Markov \term{ergodic}, phân phối dừng
thỏa mãn $\pi_j = 1/\mu_j$ (\en{ergodic theorem}).
\end{tomtat}
\end{example}

% ============================================================================
\section{Số lần quay lại (Number of Returns)}
\label{sec:number-returns}

\begin{dinhnghia}[(Số lần thăm -- \en{Number of visits})]
Số lần chuỗi Markov thăm trạng thái $i$ (không kể thời điểm $n=0$) là:
\[
  R_i := \sum_{n=1}^{\infty} \ind_{\{X_n = i\}}.
\]
\end{dinhnghia}

\begin{dinhly}[(Phân phối số lần thăm)]
\label{thm:num-visits}
Với $m \geq 1$:
\begin{equation}\label{eq:num-visits}
  \PP(R_j = m \mid X_0 = i) = p_{ij}(1 - p_{jj})(p_{jj})^{m-1},
\end{equation}
và:
\[
  \PP(R_j = 0 \mid X_0 = i) = 1 - p_{ij}.
\]
\end{dinhly}

\textbf{Chứng minh.}
Để thăm trạng thái $j$ đúng $m$ lần:
\begin{enumerate}[nosep]
  \item Trước tiên, chuỗi phải đến $j$ lần đầu (xác suất $p_{ij}$).
  \item Sau đó, quay lại $j$ thêm $m-1$ lần nữa (mỗi lần xác suất $p_{jj}$).
  \item Cuối cùng, không quay lại $j$ nữa (xác suất $1 - p_{jj}$).
\end{enumerate}
Do tính Markov, các sự kiện này độc lập, nên:
\[
  \PP(R_j = m \mid X_0 = i) = p_{ij} \cdot (p_{jj})^{m-1} \cdot (1 - p_{jj}).
\]
\hfill$\square$

Đây chính là \term{phân phối hình học} (\en{geometric distribution}) trên $\{0, 1, 2, \ldots\}$.

\begin{dinhly}[(Trường hợp $p_{jj} = 1$)]
Nếu $p_{jj} = 1$ (chuỗi chắc chắn quay lại $j$), thì:
\[
  \PP(R_j = \infty \mid X_0 = j) = 1.
\]
Ngược lại, nếu $p_{jj} < 1$, thì $\PP(R_j < \infty \mid X_0 = j) = 1$.
\end{dinhly}

\begin{congthuc}
\textbf{Kỳ vọng số lần thăm:}
\[
  \boxed{
    \E[R_j \mid X_0 = i] = \frac{p_{ij}}{1 - p_{jj}},
  }
\]
khi $p_{jj} < 1$.
\end{congthuc}

\textbf{Chứng minh.}
\begin{align*}
  \E[R_j \mid X_0 = i]
  &= \sum_{m=1}^{\infty} m\, p_{ij}(1-p_{jj})(p_{jj})^{m-1}
  = p_{ij}(1-p_{jj}) \sum_{m=1}^{\infty} m\, (p_{jj})^{m-1}\\
  &= p_{ij}(1-p_{jj}) \cdot \frac{1}{(1-p_{jj})^2}
  = \frac{p_{ij}}{1 - p_{jj}}.
\end{align*}
\hfill$\square$

\begin{tomtat}
Nếu trạng thái $j$ là \term{thoáng qua} (\en{transient}), tức $p_{jj} < 1$, thì số lần thăm
$R_j$ là hữu hạn hầu chắc chắn và có phân phối hình học.
Nếu $j$ là \term{hồi quy} (\en{recurrent}), tức $p_{jj} = 1$, thì chuỗi quay lại $j$ vô hạn lần.
\end{tomtat}

% --- Maze problem ---
\begin{example}[Bài toán mê cung: Cá trong bể 9 ngăn -- \en{Maze problem: Fish in 9-compartment aquarium}]
\label{ex:fish-maze}
Xét một con cá bơi trong bể có $9$ ngăn xếp thành lưới $3 \times 3$, đánh số từ $1$ đến $9$.
Tại mỗi bước, cá di chuyển ngẫu nhiên đều sang một ngăn kề (chia sẻ cạnh chung).

Đây là chuỗi Markov trên $S = \{1, 2, \ldots, 9\}$. Ta có thể áp dụng phân tích bước đầu tiên
để tính:
\begin{itemize}[nosep]
  \item Xác suất cá đến ngăn thức ăn (ví dụ ngăn $9$) trước khi bị bắt (ví dụ ngăn $1$).
  \item Thời gian trung bình để cá đến một ngăn cho trước.
  \item Số lần trung bình cá thăm mỗi ngăn trước khi bị hấp thụ.
\end{itemize}

Ví dụ, nếu ngăn $1$ và ngăn $9$ là hấp thụ, ta đặt $A = \{1, 9\}$ và giải hệ phương trình
tuyến tính $7 \times 7$ cho các trạng thái còn lại $\{2, 3, 4, 5, 6, 7, 8\}$.

Ma trận chuyển được xây dựng từ cấu trúc kề của lưới: mỗi ngăn góc có $2$ ngăn kề,
mỗi ngăn cạnh có $3$ ngăn kề, và ngăn trung tâm có $4$ ngăn kề.
\end{example}

% ============================================================================
\section{Tổng kết chương (Chapter Summary)}
\label{sec:ch5-summary}

\begin{tomtat}
\textbf{Phân tích bước đầu tiên} (\en{first step analysis}) là kỹ thuật quan trọng nhất
để giải các bài toán về chuỗi Markov hữu hạn. Ý tưởng cốt lõi:
\emph{điều kiện hóa theo bước đầu tiên $Z_1$, rồi dùng tính Markov để quy về bài toán cùng dạng.}

\medskip
\begin{center}
\renewcommand{\arraystretch}{1.5}
\begin{tabular}{|l|c|c|}
\hline
\textbf{Đại lượng} & \textbf{Phương trình} & \textbf{Điều kiện biên} \\
\hline
Xác suất chạm $g_l(k)$ &
$g_l(k) = \displaystyle\sum_m P_{k,m}\, g_l(m)$ &
$g_l(k) = \ind_{\{k=l\}}$, $k \in A$ \\
\hline
Thời gian chạm TB $h_A(k)$ &
$h_A(k) = 1 + \displaystyle\sum_l P_{k,l}\, h_A(l)$ &
$h_A(l) = 0$, $l \in A$ \\
\hline
Thời gian quay lại TB $\mu_j(i)$ &
$\mu_j(i) = 1 + \displaystyle\sum_{l \neq j} P_{i,l}\, \mu_j(l)$ &
--- \\
\hline
\end{tabular}
\end{center}

\medskip
Dạng ma trận cho chuỗi hấp thụ $P = \begin{bsmallmatrix} Q & R \\ 0 & I \end{bsmallmatrix}$:
\[
  \mathbf{g}_l = (I-Q)^{-1} R\, \mathbf{e}_l, \qquad
  \mathbf{h}_A = (I-Q)^{-1}\, \mathbf{1}.
\]
\end{tomtat}

% ============================================================================
\section{Bài toán con bạc -- Gambler's Ruin}
\label{sec:gamblers-ruin}

\textbf{Câu chuyện:} Một người chơi bắt đầu với $i$ đô la. Mỗi ván, anh ta thắng \$1 (xác suất $p$)
hoặc thua \$1 (xác suất $q = 1-p$). Trò chơi dừng khi cháy túi ($0$ đô la) hoặc đạt mục tiêu ($N$ đô la).

Đây là bài toán kinh điển nhất của phân tích bước đầu tiên.

\subsection{Xác suất phá sản (Ruin Probability)}

Đặt $g(i) = \PP(\text{đến } N \text{ trước } 0 \mid Z_0 = i)$ = xác suất thắng.

Phân tích bước đầu tiên:
\[
  g(i) = p\, g(i+1) + q\, g(i-1), \qquad 1 \leq i \leq N-1,
\]
với $g(0) = 0$ (cháy túi = thua) và $g(N) = 1$ (đạt mục tiêu = thắng).

\textbf{Giải phương trình:} Viết lại: $p[g(i+1) - g(i)] = q[g(i) - g(i-1)]$.

Đặt $d_i = g(i) - g(i-1)$. Ta có $d_{i+1} = (q/p)\, d_i$, suy ra:
\[
  d_k = d_1 \left(\frac{q}{p}\right)^{k-1}, \qquad k \geq 1.
\]

Từ $g(N) - g(0) = \sum_{k=1}^{N} d_k = 1$:

\begin{congthuc}
\textbf{Xác suất thắng trong bài toán con bạc:}

\textbf{Trường hợp $p \neq q$} (đặt $\rho = q/p$):
\begin{equation}\label{eq:gambler-win}
  g(i) = \frac{1 - \rho^i}{1 - \rho^N}, \qquad 0 \leq i \leq N.
\end{equation}

\textbf{Trường hợp $p = q = 1/2$} (trò chơi ``công bằng''):
\begin{equation}\label{eq:gambler-fair}
  g(i) = \frac{i}{N}.
\end{equation}
\end{congthuc}

\begin{example}[Con bạc với \$3, mục tiêu \$5]
$N = 5$, $i = 3$, $p = 0.4$ (bất lợi), $\rho = q/p = 0.6/0.4 = 1.5$.
\[
  g(3) = \frac{1 - 1.5^3}{1 - 1.5^5} = \frac{1 - 3.375}{1 - 7.594} = \frac{-2.375}{-6.594} \approx 0.360.
\]
Chỉ $36\%$ cơ hội thắng! So với trò chơi công bằng: $g(3) = 3/5 = 60\%$.
\end{example}

\subsection{Thời gian chơi trung bình (Expected Duration)}

Đặt $h(i) = \E[T \mid Z_0 = i]$ = thời gian trung bình cho đến khi dừng.

Phân tích bước đầu tiên:
\[
  h(i) = 1 + p\, h(i+1) + q\, h(i-1), \qquad 1 \leq i \leq N-1,
\]
với $h(0) = h(N) = 0$.

\textbf{Nghiệm cho $p = q = 1/2$:}
\begin{equation}\label{eq:gambler-time-fair}
  h(i) = i(N - i).
\end{equation}

\textbf{Nghiệm cho $p \neq q$:}
\begin{equation}\label{eq:gambler-time-biased}
  h(i) = \frac{1}{q - p}\left[i - N \cdot \frac{1 - \rho^i}{1 - \rho^N}\right], \qquad \rho = q/p.
\end{equation}

\begin{example}[Thời gian chơi -- trò chơi công bằng]
$N = 10$, $p = q = 1/2$, xuất phát $i = 5$:
\[
  h(5) = 5 \times (10 - 5) = 25 \text{ ván}.
\]
Xuất phát $i = 1$: $h(1) = 1 \times 9 = 9$ ván.
Xuất phát $i = 9$: $h(9) = 9 \times 1 = 9$ ván.

Trò chơi kéo dài nhất khi xuất phát ở giữa!
\end{example}

\begin{tomtat}
Bài toán con bạc cho thấy sức mạnh của phân tích bước đầu tiên:
\begin{itemize}[nosep]
  \item Xác suất thắng phụ thuộc vào $\rho = q/p$ và vị trí ban đầu.
  \item Nếu $p < 1/2$ (bất lợi), xác suất phá sản rất cao khi $N$ lớn: $g(i) \to 0$.
  \item Trò chơi ``công bằng'' ($p = 1/2$): xác suất thắng tỷ lệ với vốn ban đầu.
  \item Đây cũng là mô hình cho bước đi ngẫu nhiên với biên hấp thụ.
\end{itemize}
\end{tomtat}

% ============================================================================
\section{Bài tập (Exercises)}
\label{sec:ch5-exercises}

\begin{exercise}
Xét chuỗi Markov trên $S = \{0, 1, 2, 3, 4\}$ với các trạng thái hấp thụ $A = \{0, 4\}$
và ma trận chuyển:
\[
  P = \begin{bmatrix}
    1 & 0 & 0 & 0 & 0\\
    1/3 & 0 & 2/3 & 0 & 0\\
    0 & 1/2 & 0 & 1/2 & 0\\
    0 & 0 & 1/3 & 0 & 2/3\\
    0 & 0 & 0 & 0 & 1
  \end{bmatrix}.
\]
\begin{enumerate}[label=(\alph*)]
  \item Tính xác suất bị hấp thụ vào trạng thái $0$, xuất phát từ mỗi trạng thái $k = 1, 2, 3$.
  \item Tính thời gian hấp thụ trung bình $h_A(k)$ với $k = 1, 2, 3$.
\end{enumerate}
\end{exercise}

\begin{exercise}
Cho chuỗi Markov hai trạng thái với ma trận chuyển
$P = \begin{bsmallmatrix} 1-a & a \\ b & 1-b \end{bsmallmatrix}$,
trong đó $a, b \in (0,1)$.
\begin{enumerate}[label=(\alph*)]
  \item Chứng minh rằng $\mu_0(0) = (a+b)/b$ và $\mu_1(1) = (a+b)/a$.
  \item Xác minh rằng $\pi_j = 1/\mu_j(j)$ với phân phối dừng $\pi$ đã tìm ở Chương~\ref{ch:discrete}.
\end{enumerate}
\end{exercise}

\begin{exercise}
Một con chuột ở trong mê cung gồm $5$ phòng $\{1, 2, 3, 4, 5\}$, xếp thành hàng ngang.
Tại mỗi bước, chuột di chuyển sang phòng kề bên trái hoặc bên phải với xác suất bằng nhau.
Phòng $1$ có phô mai và phòng $5$ có bẫy (cả hai đều là trạng thái hấp thụ).
\begin{enumerate}[label=(\alph*)]
  \item Viết ma trận chuyển.
  \item Tính xác suất chuột đến phô mai trước bẫy, xuất phát từ phòng $k = 2, 3, 4$.
  \item Tính thời gian trung bình để chuột bị hấp thụ, xuất phát từ phòng $3$.
\end{enumerate}
\end{exercise}

\begin{exercise}
Cho chuỗi Markov trên $S = \{0, 1, 2\}$ với ma trận chuyển:
\[
  P = \begin{bmatrix}
    0 & 1 & 0\\
    1/4 & 1/2 & 1/4\\
    0 & 1 & 0
  \end{bmatrix}.
\]
\begin{enumerate}[label=(\alph*)]
  \item Tính $p_{00}$, $p_{11}$, $p_{22}$ (xác suất quay lại).
  \item Tính $\mu_0(0)$, $\mu_1(1)$, $\mu_2(2)$ (thời gian quay lại trung bình).
  \item Tính $\E[R_1 \mid X_0 = 0]$ (kỳ vọng số lần thăm trạng thái $1$ xuất phát từ $0$).
\end{enumerate}
\end{exercise}

\begin{exercise}
Chứng minh công thức \eqref{eq:Pn-fn}:
\[
  P_{i,i}^{(n)} = \sum_{k=1}^{n} f_{ii}^{(k)}\, P_{i,i}^{(n-k)}, \qquad n \geq 1,
\]
bằng cách phân hoạch sự kiện $\{X_n = i\}$ theo giá trị của $\tau_i$.
\end{exercise}

\begin{exercise}
Xét bể cá $9$ ngăn (Ví dụ~\ref{ex:fish-maze}). Giả sử cá xuất phát ở ngăn $5$ (trung tâm),
ngăn $1$ có thức ăn (hấp thụ) và ngăn $9$ là bẫy (hấp thụ).
\begin{enumerate}[label=(\alph*)]
  \item Viết ma trận chuyển $9 \times 9$.
  \item Xác định ma trận $Q$ (phần không hấp thụ) và $R$.
  \item Tính (hoặc thiết lập hệ phương trình) xác suất cá đến thức ăn trước bẫy.
\end{enumerate}
\end{exercise}

\chapter{Phân loại trạng thái\\{\normalsize\textit{Classification of States}}}
\label{ch:classification}

\textbf{Câu chuyện:} Hãy tưởng tượng bạn đi du lịch giữa các thành phố.
Có những thành phố bạn chắc chắn sẽ quay lại (quán café quen --- ``hồi quy''),
và có những thành phố bạn chỉ ghé qua một lần rồi không bao giờ quay lại (``tạm thời'').
Chương này giúp phân loại: trạng thái nào ``quen'' (hồi quy), trạng thái nào ``lạ'' (tạm thời)?

% ============================================================================
\section{Trạng thái liên thông (Communicating States)}
\label{sec:communicating-states}

\textbf{Câu chuyện:} Hãy tưởng tượng bạn đang ở thành phố $i$. Câu hỏi đặt ra:
``Liệu bạn có thể đi đến thành phố $j$ được không?''
Nếu có một con đường (dù phải đi qua nhiều thành phố trung gian),
tức là tồn tại số bước $n$ sao cho $P_{i,j}^{(n)} > 0$,
thì ta nói $j$ \textbf{có thể đạt đến} từ $i$.

Ví dụ đời thường: Từ Hà Nội bạn có thể bay đến Tokyo (qua Seoul chuyển tiếp)
--- vậy Tokyo ``có thể đạt đến'' từ Hà Nội.

\begin{dinhnghia}[(Trạng thái có thể đạt đến -- \en{Accessible state})]
\label{def:accessible}
Cho chuỗi Markov $(X_n)_{n \in \N}$ với không gian trạng thái $S$ và ma trận chuyển $P$.
Trạng thái $j$ được gọi là \term{có thể đạt đến} (\en{accessible}) từ trạng thái $i$,
ký hiệu $i \to j$, nếu tồn tại $n \geq 0$ sao cho:
\[
  P_{i,j}^{(n)} := \PP(X_n = j \mid X_0 = i) > 0.
\]
Nói cách khác, xuất phát từ $i$, chuỗi Markov có xác suất dương để đến được $j$
sau một số bước hữu hạn nào đó.
\end{dinhnghia}

\begin{dinhnghia}[(Trạng thái liên thông -- \en{Communicating states})]
\label{def:communicating}
Hai trạng thái $i$ và $j$ được gọi là \term{liên thông} (\en{communicate}),
ký hiệu $i \leftrightarrow j$, nếu $i \to j$ và $j \to i$, tức là:
\[
  \exists\, a \geq 0: \; P_{i,j}^{(a)} > 0
  \qquad \text{và} \qquad
  \exists\, b \geq 0: \; P_{j,i}^{(b)} > 0.
\]
\end{dinhnghia}

Ví dụ: Hà Nội và Tokyo ``liên thông'' vì bạn có thể bay đi lẫn bay về.
Nhưng nếu đảo $X$ chỉ có chuyến bay đến mà không có chuyến bay đi,
thì $X$ ``có thể đạt đến'' nhưng không ``liên thông'' với các thành phố khác.

Quan hệ ``liên thông'' $\leftrightarrow$ là một \term{quan hệ tương đương} (\en{equivalence relation})
trên $S$, vì nó thỏa mãn ba tính chất:

\begin{itemize}[nosep]
  \item \textbf{Phản xạ} (\en{Reflexive}): $i \leftrightarrow i$ vì $P_{i,i}^{(0)} = 1 > 0$.
  \item \textbf{Đối xứng} (\en{Symmetric}): $i \leftrightarrow j \Rightarrow j \leftrightarrow i$ (theo định nghĩa).
  \item \textbf{Bắc cầu} (\en{Transitive}): Nếu $i \leftrightarrow j$ và $j \leftrightarrow k$ thì $i \leftrightarrow k$.
        Thật vậy, nếu $P_{i,j}^{(a)} > 0$ và $P_{j,k}^{(b)} > 0$ thì theo phương trình Chapman--Kolmogorov:
        \[
          P_{i,k}^{(a+b)} = \sum_{l \in S} P_{i,l}^{(a)} P_{l,k}^{(b)}
          \geq P_{i,j}^{(a)} P_{j,k}^{(b)} > 0.
        \]
\end{itemize}

Vì $\leftrightarrow$ là quan hệ tương đương, nó \term{phân hoạch} (\en{partition}) không gian trạng thái
$S$ thành các \term{lớp tương đương} (\en{equivalence classes}), gọi là \term{lớp liên thông}
(\en{communicating classes}):
\begin{equation}\label{eq:partition}
  S = A_1 \cup A_2 \cup \cdots \cup A_M,
  \qquad A_k \cap A_l = \emptyset \;\text{với}\; k \neq l.
\end{equation}

\begin{tomtat}
Quan hệ $\leftrightarrow$ phân chia không gian trạng thái $S$ thành các lớp liên thông rời nhau.
Bên trong mỗi lớp, mọi cặp trạng thái đều có thể đạt đến nhau.
Giữa các lớp khác nhau, việc chuyển đổi có thể chỉ xảy ra theo một chiều hoặc không xảy ra.

Ví dụ: trong mạng lưới giao thông, các thành phố trong cùng khu vực đều liên thông.
Nhưng ``vùng hấp thụ'' (ví dụ: hố đen) chỉ có đường vào mà không có đường ra.
\end{tomtat}

% ============================================================================
\section{Lớp liên thông (Communicating Class)}
\label{sec:communicating-class}

\begin{dinhnghia}[(Chuỗi bất khả quy -- \en{Irreducible chain})]
\label{def:irreducible}
Một chuỗi Markov được gọi là \term{bất khả quy} (\en{irreducible}) nếu toàn bộ không gian
trạng thái $S$ chỉ gồm duy nhất một lớp liên thông, tức là mọi cặp trạng thái $i, j \in S$
đều liên thông: $i \leftrightarrow j$.

Ngược lại, nếu $S$ có ít nhất hai lớp liên thông thì chuỗi được gọi là
\term{khả quy} (\en{reducible}).
\end{dinhnghia}

\textbf{Ví dụ trực quan:}
\begin{itemize}[nosep]
  \item \textbf{Bất khả quy:} Bạn có thể đi từ bất kỳ thành phố nào đến bất kỳ thành phố nào (dù qua trung chuyển).
  \item \textbf{Khả quy:} Có ``vùng cấm'' --- một khi vào rồi, không ra được.
\end{itemize}

\begin{example}[Chuỗi khả quy -- \en{Reducible chain}]
\label{ex:reducible-chain}
Xét chuỗi Markov trên $S = \{0, 1, 2, 3\}$ với ma trận chuyển:
\begin{equation}\label{eq:reducible-matrix}
  P = \begin{bmatrix}
    0   & 0.2 & 0.8 & 0   \\
    0.3 & 0.1 & 0   & 0.6 \\
    0.5 & 0   & 0   & 0.5 \\
    0   & 0   & 0   & 1
  \end{bmatrix}.
\end{equation}

\textbf{Phân tích:}
\begin{itemize}[nosep]
  \item Từ trạng thái $0$: có thể đi đến $1$ (qua $P_{0,1} = 0.2$) và đến $2$ (qua $P_{0,2} = 0.8$).
  \item Từ trạng thái $1$: có thể đi đến $0$ (qua $P_{1,0} = 0.3$) và đến $3$ (qua $P_{1,3} = 0.6$).
  \item Từ trạng thái $2$: có thể đi đến $0$ (qua $P_{2,0} = 0.5$) và đến $3$ (qua $P_{2,3} = 0.5$).
  \item Từ trạng thái $3$: $P_{3,3} = 1$, trạng thái $3$ là \term{hấp thụ} (\en{absorbing}).
\end{itemize}

Các trạng thái $0, 1, 2$ liên thông với nhau, nhưng từ $3$ không thể quay lại $0, 1, 2$.
Do đó chuỗi có hai lớp liên thông:
\[
  A_1 = \{0, 1, 2\}, \qquad A_2 = \{3\}.
\]
Chuỗi này là \term{khả quy} (\en{reducible}).
\end{example}

\begin{tomtat}
Chuỗi \term{bất khả quy}: chỉ có một lớp liên thông, mọi trạng thái đều đạt đến được từ bất kỳ
trạng thái nào khác.

Chuỗi \term{khả quy}: có nhiều lớp liên thông; một khi vào một lớp ``đóng'' (\en{closed class}),
chuỗi không thể thoát ra.
\end{tomtat}

% ============================================================================
\section{Trạng thái hồi quy (Recurrent States)}
\label{sec:recurrent-states}

\textbf{Câu chuyện:} Quán café quen là ``hồi quy'' --- dù bạn đi đâu, chắc chắn một ngày sẽ quay lại.
Quán ghé qua lần đầu ở thành phố lạ là ``tạm thời'' --- có thể bạn không bao giờ quay lại nữa.

Nhớ lại từ Chương~\ref{ch:first-step}, \term{thời gian quay lại lần đầu} (\en{first return time}) trạng thái $i$
được định nghĩa là:
\[
  \tau_i := \inf\{n \geq 1 : X_n = i\},
\]
và \term{xác suất quay lại} (\en{return probability}):
\[
  p_{ii} := \PP(\tau_i < \infty \mid X_0 = i).
\]

\begin{dinhnghia}[(Trạng thái hồi quy -- \en{Recurrent state})]
\label{def:recurrent}
Trạng thái $i \in S$ được gọi là \term{hồi quy} (\en{recurrent}) nếu:
\begin{equation}\label{eq:recurrent-def}
  p_{ii} = \PP(\tau_i < \infty \mid X_0 = i) = 1.
\end{equation}
Tức là, xuất phát từ $i$, chuỗi Markov \emph{chắc chắn} sẽ quay lại $i$.
\end{dinhnghia}

Tính hồi quy có thể được đặc trưng bằng các điều kiện tương đương sau.

\begin{dinhly}[(Đặc trưng của trạng thái hồi quy)]
\label{thm:recurrent-equiv}
Các mệnh đề sau là tương đương:
\begin{enumerate}[label=(\roman*)]
  \item $i$ là trạng thái hồi quy, tức $p_{ii} = 1$.
  \item \term{Số lần quay lại trung bình} (\en{expected number of returns}) là vô hạn:
        \begin{equation}\label{eq:recurrent-Ri}
          \E[R_i \mid X_0 = i] = +\infty,
        \end{equation}
        trong đó $R_i = \sum_{n=1}^{\infty} \ind_{\{X_n = i\}}$ là tổng số lần chuỗi thăm $i$.
  \item $\PP(R_i = +\infty \mid X_0 = i) = 1$, tức chuỗi quay lại $i$ \emph{vô hạn lần} với xác suất~$1$.
  \item Chuỗi các xác suất quay lại phân kỳ:
        \begin{equation}\label{eq:recurrent-series}
          \sum_{n=0}^{\infty} P_{i,i}^{(n)} = +\infty.
        \end{equation}
\end{enumerate}
\end{dinhly}

\begin{congthuc}
\textbf{Tiêu chuẩn hồi quy (bằng chuỗi):}
\[
  i \text{ hồi quy} \quad\Longleftrightarrow\quad
  \sum_{n=0}^{\infty} P_{i,i}^{(n)} = +\infty.
\]
Trực quan: nếu ``tổng xác suất quay lại'' phân kỳ, thì chuỗi quay lại vô hạn lần.
\end{congthuc}

% ============================================================================
\section{Trạng thái tạm thời (Transient States)}
\label{sec:transient-states}

\begin{dinhnghia}[(Trạng thái tạm thời -- \en{Transient state})]
\label{def:transient}
Trạng thái $i \in S$ được gọi là \term{tạm thời} (\en{transient}) nếu:
\begin{equation}\label{eq:transient-def}
  \PP(R_i = +\infty \mid X_0 = i) < 1,
\end{equation}
hay tương đương, $p_{ii} < 1$.
\end{dinhnghia}

Đối với trạng thái tạm thời:
\begin{itemize}[nosep]
  \item $\E[R_i \mid X_0 = i] = \dfrac{p_{ii}}{1 - p_{ii}} < +\infty$: số lần quay lại trung bình là hữu hạn.
  \item $\displaystyle\sum_{n=0}^{\infty} P_{i,i}^{(n)} < +\infty$: chuỗi các xác suất quay lại hội tụ.
  \item Với xác suất $1 - p_{ii} > 0$, chuỗi sẽ \emph{rời khỏi} $i$ và không bao giờ quay lại.
\end{itemize}

\begin{tomtat}
\begin{center}
\renewcommand{\arraystretch}{1.4}
\begin{tabular}{lcc}
\hline
\textbf{Tính chất} & \textbf{Hồi quy} (\en{Recurrent}) & \textbf{Tạm thời} (\en{Transient}) \\
\hline
$p_{ii}$ & $= 1$ & $< 1$ \\
$\E[R_i \mid X_0 = i]$ & $= +\infty$ & $< +\infty$ \\
$\sum_n P_{i,i}^{(n)}$ & $= +\infty$ & $< +\infty$ \\
Số lần thăm $i$ & Vô hạn lần (h.c.c.) & Hữu hạn lần (h.c.c.) \\
Ví dụ đời thường & Quán café quen & Quán ghé qua khi du lịch \\
\hline
\end{tabular}
\end{center}
\end{tomtat}

% --- Corollary: recurrence is a class property ---

\begin{corollary}[(Hồi quy là tính chất lớp -- \en{Recurrence is a class property})]
\label{cor:recurrence-class}
Nếu trạng thái $i$ là hồi quy và $j$ liên thông với $i$ (tức $i \leftrightarrow j$),
thì $j$ cũng là hồi quy.
\end{corollary}

\begin{proof}
Vì $i \leftrightarrow j$, tồn tại $a, b \geq 0$ sao cho $P_{j,i}^{(a)} > 0$ và $P_{i,j}^{(b)} > 0$.
Theo phương trình Chapman--Kolmogorov, với mọi $n \geq 0$:
\[
  P_{j,j}^{(a+n+b)} \geq P_{j,i}^{(a)} \, P_{i,i}^{(n)} \, P_{i,j}^{(b)}.
\]
Do đó:
\[
  \sum_{m=0}^{\infty} P_{j,j}^{(m)}
  \geq P_{j,i}^{(a)} \, P_{i,j}^{(b)} \sum_{n=0}^{\infty} P_{i,i}^{(n)}
  = +\infty,
\]
vì $i$ là hồi quy nên $\sum_n P_{i,i}^{(n)} = +\infty$, và $P_{j,i}^{(a)} P_{i,j}^{(b)} > 0$.
Vậy $j$ cũng là hồi quy.
\end{proof}

\textbf{Hệ quả:} Tạm thời cũng là tính chất lớp. Nếu $i$ là tạm thời và $i \leftrightarrow j$,
thì $j$ cũng là tạm thời.

% --- Example: Simple random walk ---

\begin{example}[Bước đi ngẫu nhiên đơn giản trên $\Z$ -- \en{Simple random walk on $\Z$}]
\label{ex:random-walk-recurrence}
Xét bước đi ngẫu nhiên $(S_n)_{n \geq 0}$ trên $\Z$ với $S_0 = 0$,
$\PP(X_k = +1) = p$ và $\PP(X_k = -1) = q = 1-p$.

Xác suất trở về gốc tại thời điểm chẵn $n = 2k$:
\[
  P_{0,0}^{(2k)} = \binom{2k}{k} p^k q^k, \qquad P_{0,0}^{(2k+1)} = 0.
\]

\textbf{Trường hợp $p = q = 1/2$:} $4pq = 1$, suy ra $P_{0,0}^{(2k)} \sim \frac{1}{\sqrt{\pi k}}$ và:
\[
  \sum_{k=0}^{\infty} P_{0,0}^{(2k)} = +\infty
  \qquad \Longrightarrow \qquad
  \text{trạng thái $0$ là \textbf{hồi quy}.}
\]

\textbf{Trường hợp $p \neq q$:} $4pq < 1$, suy ra chuỗi hội tụ:
\[
  \sum_{k=0}^{\infty} P_{0,0}^{(2k)} < +\infty
  \qquad \Longrightarrow \qquad
  \text{trạng thái $0$ là \textbf{tạm thời}.}
\]

Trực quan: với $p \neq q$, bước đi ngẫu nhiên có ``xu hướng'' đi về một phía,
nên càng lúc càng xa gốc và khó quay lại.
\end{example}

\begin{congthuc}
\textbf{Bước đi ngẫu nhiên đơn giản trên $\Z$:}
\[
  \text{Hồi quy} \;\Longleftrightarrow\; p = q = \frac{1}{2}.
\]
\end{congthuc}

% --- Polya's theorem ---
\begin{dinhly}[(Định lý Polya về hồi quy -- \en{Polya's recurrence theorem})]
\label{thm:polya}
Bước đi ngẫu nhiên đơn giản đối xứng ($p = 1/2$ mỗi chiều) trên mạng lưới $\Z^d$:
\begin{itemize}[nosep]
  \item $d = 1$: \textbf{hồi quy} (chắc chắn quay lại gốc).
  \item $d = 2$: \textbf{hồi quy} (``A drunk man will find his way home'' --- Polya).
  \item $d \geq 3$: \textbf{tạm thời} (``A drunk bird will NOT find its way home'').
\end{itemize}
\end{dinhly}

Trực quan: trong không gian 1D và 2D, ``không có đủ chỗ để trốn'' nên bước đi phải quay lại.
Trong 3D trở lên, không gian quá rộng, bước đi ``lạc'' mãi mãi với xác suất dương.

% --- Theorem 4: Finite chain has at least one recurrent state ---

\begin{dinhly}[(Tồn tại trạng thái hồi quy trong chuỗi hữu hạn)]
\label{thm:finite-recurrent}
Mọi chuỗi Markov với không gian trạng thái hữu hạn $S$ có ít nhất một trạng thái hồi quy.
\end{dinhly}

\begin{proof}
Giả sử bằng phản chứng rằng \emph{tất cả} trạng thái đều tạm thời.
Khi đó, với mọi $j \in S$:
\[
  \sum_{n=0}^{\infty} P_{i,j}^{(n)} < +\infty
  \qquad \Longrightarrow \qquad
  \lim_{n \to \infty} P_{i,j}^{(n)} = 0.
\]
Lấy tổng theo $j$ trên $S$ hữu hạn:
\[
  \lim_{n \to \infty} \sum_{j \in S} P_{i,j}^{(n)}
  = \sum_{j \in S} \lim_{n \to \infty} P_{i,j}^{(n)} = 0.
\]
Nhưng $\sum_{j \in S} P_{i,j}^{(n)} = 1$ với mọi $n$, nên giới hạn phải bằng $1 \neq 0$.
Mâu thuẫn. Vậy phải tồn tại ít nhất một trạng thái hồi quy.
\end{proof}

% ============================================================================
\section{Hồi quy dương và hồi quy không (Positive and Null Recurrence)}
\label{sec:pos-null-recurrence}

\textbf{Câu chuyện:} Biết rằng bạn chắc chắn quay lại quán café quen --- nhưng MẤT BAO LÂU?
\begin{itemize}[nosep]
  \item \textbf{Hồi quy dương:} Thời gian quay lại trung bình hữu hạn (ví dụ: trung bình 5 ngày quay lại quán).
  \item \textbf{Hồi quy không:} Chắc chắn quay lại, nhưng thời gian trung bình vô hạn (``chờ mãi mãi'').
\end{itemize}

\begin{dinhnghia}[(Thời gian quay lại trung bình -- \en{Mean recurrence time})]
\[
  \mu_i(i) := \E[\tau_i \mid X_0 = i].
\]
\end{dinhnghia}

\begin{dinhnghia}[(Hồi quy dương và hồi quy không)]
Cho trạng thái $i$ là hồi quy ($p_{ii} = 1$).
\begin{itemize}[nosep]
  \item $i$ là \term{hồi quy dương} (\en{positive recurrent}) nếu $\mu_i(i) < +\infty$.
  \item $i$ là \term{hồi quy không} (\en{null recurrent}) nếu $\mu_i(i) = +\infty$.
\end{itemize}
\end{dinhnghia}

\begin{example}[Bước đi ngẫu nhiên đối xứng là hồi quy không]
Bước đi ngẫu nhiên đơn giản trên $\Z$ với $p = q = 1/2$ là \term{hồi quy không}.
Thời gian quay lại trung bình $\mu_0(0) = +\infty$.

Trực quan: bạn \emph{chắc chắn} quay lại gốc, nhưng không ai biết ``trung bình bao lâu'' ---
có thể mất 2 bước, có thể mất 2 triệu bước.
\end{example}

\begin{dinhly}[(Chuỗi hữu hạn và hồi quy dương)]
\label{thm:finite-positive-recurrent}
Trong chuỗi Markov với không gian trạng thái hữu hạn, mọi trạng thái hồi quy đều là
\term{hồi quy dương}.
\end{dinhly}

\begin{corollary}[(Chuỗi bất khả quy hữu hạn)]
\label{cor:finite-irreducible}
Nếu chuỗi Markov là \term{bất khả quy} và có không gian trạng thái hữu hạn,
thì \emph{tất cả} các trạng thái đều là \term{hồi quy dương}.
\end{corollary}

\begin{tomtat}
\begin{center}
\renewcommand{\arraystretch}{1.4}
\begin{tabular}{lccc}
\hline
\textbf{Tính chất} & \textbf{Tạm thời} & \textbf{Hồi quy không} & \textbf{Hồi quy dương} \\
\hline
$p_{ii}$ & $< 1$ & $= 1$ & $= 1$ \\
Quay lại? & Không chắc & Chắc chắn & Chắc chắn \\
$\mu_i(i)$ & --- & $= +\infty$ & $< +\infty$ \\
Ví dụ & RW trên $\Z$, $p \neq q$ & RW trên $\Z$, $p = 1/2$ & Chuỗi hữu hạn \\
\hline
\end{tabular}
\end{center}

\textbf{Quy tắc nhớ:} Chuỗi Markov hữu hạn bất khả quy $\Rightarrow$ mọi trạng thái hồi quy dương.
Hồi quy không chỉ xảy ra khi không gian trạng thái vô hạn.
\end{tomtat}

% ============================================================================
\section{Tính tuần hoàn và phi tuần hoàn (Periodicity and Aperiodicity)}
\label{sec:periodicity}

\textbf{Câu chuyện:} Một số chuỗi Markov ``dao động'' có quy luật. Ví dụ: nếu bạn ở ô trắng trên bàn cờ,
sau 1 bước chắc chắn ở ô đen, sau 2 bước lại ở ô trắng. Chuỗi này có ``nhịp'' (chu kỳ) = 2.
Chuỗi tuần hoàn không hội tụ vì nó cứ ``dao động'' mãi.

\begin{dinhnghia}[(Chu kỳ của trạng thái -- \en{Period of a state})]
\term{Chu kỳ} (\en{period}) của trạng thái $i$ được định nghĩa là:
\begin{equation}\label{eq:period}
  d(i) := \gcd\{n \geq 1 : P_{i,i}^{(n)} > 0\}.
\end{equation}
\end{dinhnghia}

\begin{dinhnghia}[(Phi tuần hoàn -- \en{Aperiodic})]
Trạng thái $i$ được gọi là \term{phi tuần hoàn} (\en{aperiodic}) nếu $d(i) = 1$.

Một chuỗi Markov được gọi là \term{phi tuần hoàn} nếu \emph{tất cả} các trạng thái đều phi tuần hoàn.
\end{dinhnghia}

\textbf{Nhận xét:}
\begin{itemize}[nosep]
  \item Nếu $P_{i,i} > 0$ thì $d(i) = 1$, tức trạng thái $i$ phi tuần hoàn.
        (Mẹo: thêm ``vòng lặp'' tại $i$ là phá vỡ tuần hoàn!)
  \item \term{Trạng thái hấp thụ} ($P_{k,k} = 1$) là đồng thời phi tuần hoàn và hồi quy.
  \item Chu kỳ là tính chất lớp: nếu $i \leftrightarrow j$ thì $d(i) = d(j)$.
\end{itemize}

\begin{example}[Chuỗi có chu kỳ $2$]
Chuỗi Markov trên $S = \{0, 1\}$ với $P = \begin{bmatrix} 0 & 1 \\ 1 & 0 \end{bmatrix}$
có $d(0) = d(1) = 2$: chuỗi ``dao động'' giữa $0$ và $1$.

$P_{0,0}^{(n)} > 0$ chỉ khi $n$ chẵn, nên $\gcd\{2, 4, 6, \ldots\} = 2$.
\end{example}

\begin{example}[Chuỗi có chu kỳ $3$]
Trên $S = \{0, 1, 2\}$ với $P = \begin{bmatrix} 0 & 1 & 0 \\ 0 & 0 & 1 \\ 1 & 0 & 0 \end{bmatrix}$.
Chuỗi quay vòng: $0 \to 1 \to 2 \to 0 \to \cdots$, chu kỳ $d = 3$.
\end{example}

% --- Ergodic definition ---

\begin{dinhnghia}[(Trạng thái ergodic -- \en{Ergodic state})]
Trạng thái $i$ được gọi là \term{ergodic} nếu nó đồng thời:
\begin{enumerate}[label=(\roman*)]
  \item \term{Hồi quy dương}: $p_{ii} = 1$ và $\mu_i(i) < +\infty$.
  \item \term{Phi tuần hoàn}: $d(i) = 1$.
\end{enumerate}
\end{dinhnghia}

\begin{tomtat}
Tính ergodic là điều kiện then chốt để chuỗi Markov hội tụ đến phân phối dừng (Chương~\ref{ch:limiting}).
\[
  \boxed{\text{Ergodic} = \text{Hồi quy dương} + \text{Phi tuần hoàn}}
\]
Chuỗi Markov bất khả quy, hữu hạn và phi tuần hoàn luôn là ergodic.

\medskip
\textbf{Checklist để kiểm tra hội tụ:}
\begin{enumerate}[nosep]
  \item Bất khả quy? (Mọi trạng thái liên thông?)
  \item Phi tuần hoàn? (Có $P_{i,i} > 0$ cho trạng thái nào đó?)
  \item Hữu hạn? (Nếu có $\Rightarrow$ tự động hồi quy dương.)
\end{enumerate}
Nếu cả 3: phân phối dừng tồn tại, duy nhất, và bằng phân phối giới hạn.
\end{tomtat}

% ============================================================================
\section{Bài tập (Exercises)}
\label{sec:ch6-exercises}

\begin{exercise}
\label{ex:ch6-1}
Xét chuỗi Markov trên $S = \{0, 1, 2, 3\}$ với ma trận chuyển:
\[
  P = \begin{bmatrix}
    1/3 & 1/3 & 1/3 & 0 \\
    0   & 0   & 0   & 1 \\
    0   & 1   & 0   & 0 \\
    0   & 0   & 1   & 0
  \end{bmatrix}.
\]
\begin{enumerate}[label=(\alph*)]
  \item Tìm các lớp liên thông.
  \item Tìm chu kỳ của mỗi trạng thái.
  \item Xác định trạng thái nào hồi quy, tạm thời, hấp thụ.
  \item Chuỗi có bất khả quy không?
\end{enumerate}
\end{exercise}

\begin{exercise}
\label{ex:ch6-2}
Chứng minh rằng chu kỳ là tính chất lớp: nếu $i \leftrightarrow j$ thì $d(i) = d(j)$.

(Gợi ý: nếu $P_{i,j}^{(a)} > 0$ và $P_{j,i}^{(b)} > 0$, thì $P_{i,i}^{(a+b)} > 0$,
nên $d(i)$ chia hết $a + b$. Dùng điều này để chỉ $d(i) \mid d(j)$ và ngược lại.)
\end{exercise}

\begin{exercise}
\label{ex:ch6-3}
Xét chuỗi Markov trên $S = \{0, 1, 2, 3, 4\}$ với ma trận chuyển:
\[
  P = \begin{bmatrix}
    0 & 1 & 0 & 0 & 0 \\
    0 & 0 & 1 & 0 & 0 \\
    1/2 & 0 & 0 & 1/2 & 0 \\
    0 & 0 & 0 & 0 & 1 \\
    0 & 0 & 0 & 1 & 0
  \end{bmatrix}.
\]
\begin{enumerate}[label=(\alph*)]
  \item Tìm các lớp liên thông và xác định lớp nào đóng.
  \item Phân loại mỗi trạng thái (hồi quy/tạm thời, tuần hoàn/phi tuần hoàn).
  \item Từ trạng thái 0, chuỗi sẽ kết thúc ở lớp nào?
\end{enumerate}
\end{exercise}

\begin{exercise}
\label{ex:ch6-4}
Cho chuỗi Markov bất khả quy trên không gian trạng thái hữu hạn $S$ với $|S| = N$.
\begin{enumerate}[label=(\alph*)]
  \item Chứng minh rằng mọi trạng thái đều hồi quy.
  \item Chứng minh rằng mọi trạng thái đều hồi quy dương. (Gợi ý: dùng Định lý~\ref{thm:finite-positive-recurrent}.)
  \item Nếu thêm điều kiện phi tuần hoàn, kết luận gì về sự tồn tại phân phối giới hạn?
\end{enumerate}
\end{exercise}

\chapter{Phân phối giới hạn và dừng\\{\normalsize\textit{Limiting and Stationary Distributions}}}
\label{ch:limiting}

\textbf{Câu chuyện:} Hãy tưởng tượng một thành phố với nhiều khu dân cư. Mỗi năm, người dân di cư giữa các khu.
Ban đầu có thể khu A đông, khu B vắng. Nhưng sau nhiều năm, tỷ lệ dân ở mỗi khu sẽ \textbf{ổn định} ---
không phụ thuộc vào phân bố ban đầu. Đây chính là \term{phân phối giới hạn}: trạng thái ``cân bằng dài hạn'' của hệ thống.

% ============================================================================
\section{Phân phối giới hạn (Limiting Distributions)}
\label{sec:limiting-dist}

\begin{dinhnghia}[(Phân phối giới hạn -- \en{Limiting distribution})]
Một chuỗi Markov $(X_n)_{n\in\N}$ với ma trận chuyển $P$ trên không gian trạng thái $S$
được gọi là có \term{phân phối giới hạn} (\en{limiting distribution}) nếu giới hạn
\[
  \pi_j := \lim_{n\to\infty} \PP(X_n = j \mid X_0 = i)
\]
tồn tại với mọi $i, j \in S$, và \emph{không phụ thuộc} vào trạng thái ban đầu $i$.
\end{dinhnghia}

\begin{tomtat}
Phân phối giới hạn $\pi = (\pi_j)_{j\in S}$ mô tả xác suất ``dài hạn'' của chuỗi Markov:
sau thời gian đủ lớn, xác suất ở trạng thái $j$ tiến tới $\pi_j$,
bất kể chuỗi xuất phát từ đâu.

Ví dụ: sau nhiều năm di cư, dù ban đầu tất cả sống ở khu A, tỷ lệ dân vẫn tiến tới $\pi$.
\end{tomtat}

% --- Regular matrices ---
\subsection{Ma trận chính quy (Regular Matrices)}
\label{subsec:regular-matrix}

\begin{dinhnghia}[(Ma trận chính quy -- \en{Regular matrix})]
Ma trận chuyển $P$ được gọi là \term{chính quy} (\en{regular}) nếu tồn tại
lũy thừa $P^n$ (với $n$ nào đó) mà tất cả các phần tử đều dương.
\end{dinhnghia}

Nếu $P$ là chính quy, thì chuỗi Markov tương ứng có phân phối giới hạn
$\pi_j = \lim_{n\to\infty} \PP(X_n = j \mid X_0 = i)$, và giới hạn này \emph{không phụ thuộc} vào $i$.

% --- Two-state example ---
\subsection{Ví dụ: Chuỗi hai trạng thái (Two-State Chain Example)}

Xét chuỗi Markov hai trạng thái $S = \{0, 1\}$ với ma trận chuyển:
\[
  P = \begin{bmatrix} 1-a & a \\ b & 1-b \end{bmatrix},
  \qquad a, b \in (0,1].
\]

Từ Chương~\ref{ch:discrete}, ta đã biết phân phối giới hạn là:
\begin{equation}\label{eq:ch7-pi-two-state}
  (\pi_0, \pi_1) = \left(\frac{b}{a+b},\; \frac{a}{a+b}\right).
\end{equation}

\textbf{Thời gian quay lại trung bình} (\en{mean return times}).
Với chuỗi hai trạng thái, thời gian quay lại trung bình của các trạng thái là:
\begin{equation}\label{eq:mean-return-two-state}
  \mu_0(0) = \frac{a+b}{b}, \qquad
  \mu_1(1) = \frac{a+b}{a}.
\end{equation}

\begin{tomtat}
\textbf{Quan sát quan trọng:}
So sánh \eqref{eq:ch7-pi-two-state} và \eqref{eq:mean-return-two-state}, ta thấy:
\[
  \pi_0 = \frac{b}{a+b} = \frac{1}{\mu_0(0)}, \qquad
  \pi_1 = \frac{a}{a+b} = \frac{1}{\mu_1(1)}.
\]
Xác suất giới hạn của trạng thái $j$ bằng \emph{nghịch đảo} của thời gian quay lại trung bình $\mu_j(j)$.
Đây là kết quả tổng quát, không chỉ đúng cho chuỗi hai trạng thái.

Trực quan: nếu trung bình cứ 5 bước bạn quay lại quán quen, thì ``tỷ lệ thời gian'' ở quán là $1/5 = 20\%$.
\end{tomtat}

% --- Main theorem ---
\subsection{Định lý chính về phân phối giới hạn}

\begin{dinhly}[{Định lý 6 -- Phân phối giới hạn (\en{Theorem IV.4.1 in Feller})}]
\label{thm:limiting-dist}
Cho chuỗi Markov $(X_n)_{n\in\N}$ \term{bất khả quy} (\en{irreducible}),
\term{hồi quy} (\en{recurrent}), và \term{không tuần hoàn} (\en{aperiodic})
trên không gian trạng thái $S$. Khi đó:
\begin{equation}\label{eq:limiting-thm}
  \lim_{n\to\infty} \PP(X_n = i \mid X_0 = j)
  = \frac{1}{\mu_i(i)}, \qquad \text{với mọi } i, j \in S,
\end{equation}
trong đó $\mu_i(i) = \E[\tau_i \mid X_0 = i]$ là \term{thời gian hồi quy trung bình}
(\en{mean recurrence time}) của trạng thái $i$, và
$\tau_i = \inf\{n \geq 1 : X_n = i\}$ là \term{thời điểm quay lại} (\en{return time}).
\end{dinhly}

\begin{congthuc}
\textbf{Phân phối giới hạn và thời gian hồi quy:}
\[
  \boxed{\pi_j = \frac{1}{\mu_j(j)}}
\]
Xác suất dài hạn ở trạng thái $j$ tỷ lệ nghịch với thời gian trung bình để quay lại $j$.
Trạng thái có thời gian hồi quy ngắn sẽ được ``ghé thăm'' thường xuyên hơn.
\end{congthuc}

\textbf{Giải thích trực quan:} Nếu trung bình cứ $\mu_j(j)$ bước ta quay lại trạng thái $j$ một lần,
thì tỷ lệ thời gian ở trạng thái $j$ xấp xỉ $1/\mu_j(j)$.

\textbf{Các điều kiện cần thiết:}
\begin{itemize}[nosep]
  \item \textbf{Bất khả quy} (\en{irreducible}): mọi trạng thái đều liên thông với nhau.
  \item \textbf{Hồi quy} (\en{recurrent}): chuỗi chắc chắn quay lại mỗi trạng thái.
  \item \textbf{Không tuần hoàn} (\en{aperiodic}): chu kỳ $d(i) = 1$ với mọi $i$.
\end{itemize}

Nếu chuỗi có \term{chu kỳ} (\en{periodic}) $d > 1$, thì giới hạn
$\lim_{n\to\infty} \PP(X_n = j \mid X_0 = i)$ \emph{không tồn tại} (xem ví dụ ở Mục~\ref{subsec:periodic-no-limit}).

% ============================================================================
\section{Phân phối dừng (Stationary Distributions)}
\label{sec:stationary-dist}

\textbf{Câu chuyện:} Hãy tưởng tượng bạn đổ nước vào các bình thông nhau.
Sau một lúc, mực nước ổn định --- đó là trạng thái ``dừng''.
Phân phối dừng $\pi$ giống vậy: nếu ban đầu phân phối đã là $\pi$,
thì sau bao nhiêu bước đi nữa, phân phối vẫn là $\pi$. Hệ thống ``đứng yên''.

\begin{dinhnghia}[(Phân phối dừng -- \en{Stationary distribution})]
\label{def:stationary}
Vectơ xác suất $\pi = (\pi_i)_{i\in S}$ với $\pi_i \geq 0$ và $\sum_{i\in S} \pi_i = 1$
được gọi là \term{phân phối dừng} (\en{stationary distribution}) (hay \term{phân phối bất biến}
-- \en{invariant distribution}) của chuỗi Markov nếu:
khi $X_0$ có phân phối $\pi$ (tức $\PP(X_0 = i) = \pi_i$),
thì $X_1$ cũng có phân phối $\pi$.
\end{dinhnghia}

Điều kiện trên tương đương với:
\begin{equation}\label{eq:stationary-condition}
  \pi_j = \sum_{i \in S} \pi_i \, P_{i,j}, \qquad \text{với mọi } j \in S,
\end{equation}
hay viết dưới dạng ma trận:
\begin{equation}\label{eq:pi-P}
  \boxed{\pi = \pi P.}
\end{equation}

\begin{tomtat}
$\pi$ là phân phối dừng nếu và chỉ nếu $\pi P = \pi$:
$\pi$ là \term{vectơ riêng trái} (\en{left eigenvector}) của $P$ ứng với trị riêng $\lambda = 1$.

Vì quá trình Markov là \term{đồng nhất theo thời gian} (\en{time homogeneous}),
nếu $X_0 \sim \pi$ thì $X_k \sim \pi$ với \emph{mọi} $k \geq 1$.
Do đó, phân phối $\pi$ được ``bảo toàn'' qua thời gian -- vì vậy gọi là ``dừng'' hay ``bất biến''.
\end{tomtat}

% --- Proposition 1 ---
\begin{proposition}[Phân phối giới hạn là phân phối dừng]
\label{prop:limiting-is-stationary}
Cho chuỗi Markov trên $S = \{0, \ldots, N\}$ hữu hạn. Nếu phân phối giới hạn
\[
  \pi_j = \lim_{n\to\infty} \PP(X_n = j \mid X_0 = i)
\]
tồn tại và không phụ thuộc vào $i$, thì $\pi = (\pi_j)_{j\in S}$ là phân phối dừng: $\pi = \pi P$.
\end{proposition}

\begin{proof}
Với mọi $j \in S$, bằng tính chất Markov:
\begin{align*}
  \pi_j
  &= \lim_{n\to\infty} \PP(X_{n+1} = j \mid X_0 = i) \\
  &= \lim_{n\to\infty} \sum_{l \in S} \PP(X_{n+1} = j \mid X_n = l) \, \PP(X_n = l \mid X_0 = i) \\
  &= \sum_{l \in S} P_{l,j} \lim_{n\to\infty} \PP(X_n = l \mid X_0 = i) \\
  &= \sum_{l \in S} P_{l,j} \, \pi_l.
\end{align*}
Đây chính là $\pi_j = \sum_{l} \pi_l P_{l,j}$, hay $\pi = \pi P$.
\end{proof}

% --- Important remarks ---
\subsection{Lưu ý quan trọng}
\label{subsec:periodic-no-limit}

\textbf{1. Phân phối dừng có thể tồn tại mà không có phân phối giới hạn.}

\begin{example}[Chuỗi tuần hoàn -- \en{Periodic chain}]
Xét chuỗi Markov hai trạng thái với $a = b = 1$:
\[
  P = \begin{bmatrix} 0 & 1 \\ 1 & 0 \end{bmatrix}.
\]
Chuỗi này có chu kỳ $d = 2$: từ trạng thái $0$ luôn chuyển sang $1$ và ngược lại.
Do đó $\lim_{n\to\infty} P^n$ \emph{không tồn tại}.

Tuy nhiên, $\pi = (1/2, 1/2)$ là phân phối dừng vì:
\[
  \begin{bmatrix} 1/2 & 1/2 \end{bmatrix}
  \begin{bmatrix} 0 & 1 \\ 1 & 0 \end{bmatrix}
  = \begin{bmatrix} 1/2 & 1/2 \end{bmatrix}.
\]
\end{example}

\textbf{2. Phân phối giới hạn (khi tồn tại và không phụ thuộc $i$) luôn là phân phối dừng}
(theo Mệnh đề~\ref{prop:limiting-is-stationary}).

\begin{tomtat}
\begin{center}
\begin{tabular}{ll}
  \textbf{Phân phối giới hạn tồn tại} & $\Rightarrow$ \textbf{Phân phối dừng tồn tại} \\
  \textbf{Phân phối dừng tồn tại} & $\not\Rightarrow$ \textbf{Phân phối giới hạn tồn tại}
\end{tabular}
\end{center}
Chuỗi tuần hoàn có phân phối dừng nhưng không có phân phối giới hạn.
\end{tomtat}

% ============================================================================
\section{Tính phân phối dừng (Computing Stationary Distributions)}
\label{sec:computing-stationary}

Để tìm phân phối dừng $\pi = (\pi_i)_{i\in S}$, ta giải hệ phương trình tuyến tính:

\begin{congthuc}
\textbf{Hệ phương trình cho phân phối dừng:}
\[
  \begin{cases}
    \pi = \pi P & \text{(phương trình bất biến)}, \\
    \displaystyle\sum_{i\in S} \pi_i = 1 & \text{(điều kiện chuẩn hóa)}.
  \end{cases}
\]
Phương trình $\pi = \pi P$ tương đương với $\pi (P - I) = 0$, hay $\pi$ nằm trong
\term{không gian riêng trái} (\en{left eigenspace}) ứng với trị riêng $1$ của $P$.
\end{congthuc}

% --- Example 1: Two-state ---
\begin{example}[Chuỗi hai trạng thái]
\label{ex:stationary-two-state}
Với $P = \begin{bmatrix} 1-a & a \\ b & 1-b \end{bmatrix}$,
hệ $\pi = \pi P$ cho:
\begin{align*}
  \pi_0 &= (1-a)\pi_0 + b\pi_1, \\
  \pi_1 &= a\pi_0 + (1-b)\pi_1.
\end{align*}
Cả hai phương trình đều cho $a\pi_0 = b\pi_1$.
Kết hợp với $\pi_0 + \pi_1 = 1$:
\[
  \pi_0 = \frac{b}{a+b}, \qquad \pi_1 = \frac{a}{a+b}.
\]
Kiểm tra: $\pi P = \pi$. Đúng! Kết quả trùng với phân phối giới hạn đã tìm trước đó.
\end{example}

% --- Example 2: Three-state ---
\begin{example}[Chuỗi ba trạng thái -- chi tiết]
\label{ex:stationary-three-state}
Xét chuỗi Markov trên $S = \{0, 1, 2\}$ với ma trận chuyển:
\[
  P = \begin{bmatrix}
    0 & 1/2 & 1/2 \\
    1/4 & 1/2 & 1/4 \\
    1/4 & 1/4 & 1/2
  \end{bmatrix}.
\]
Hệ phương trình $\pi = \pi P$:
\begin{align*}
  \pi_0 &= \tfrac{1}{4}\pi_1 + \tfrac{1}{4}\pi_2, \\
  \pi_1 &= \tfrac{1}{2}\pi_0 + \tfrac{1}{2}\pi_1 + \tfrac{1}{4}\pi_2, \\
  \pi_2 &= \tfrac{1}{2}\pi_0 + \tfrac{1}{4}\pi_1 + \tfrac{1}{2}\pi_2.
\end{align*}
Từ phương trình thứ nhất: $4\pi_0 = \pi_1 + \pi_2 = 1 - \pi_0$, suy ra $\pi_0 = 1/5$.
Từ phương trình thứ hai: $\pi_1 = \pi_0 + \tfrac{1}{2}\pi_2$, kết hợp $\pi_1 + \pi_2 = 4/5$:
\[
  \pi_1 = \frac{2}{5}, \qquad \pi_2 = \frac{2}{5}.
\]
Vậy $\pi = (1/5,\; 2/5,\; 2/5)$.

\textbf{Thời gian hồi quy:} $\mu_0(0) = 1/\pi_0 = 5$, $\mu_1(1) = 1/\pi_1 = 5/2$, $\mu_2(2) = 5/2$.

Trạng thái 0 được ghé thăm ít nhất (cứ 5 bước quay lại 1 lần).
\end{example}

% --- Example 3: Cover graph (5 states) ---
\begin{example}[Đồ thị bìa sách -- \en{Cover graph chain}]
\label{ex:cover-graph}
Xét chuỗi Markov trên $S = \{1, 2, 3, 4, 5\}$ với ma trận chuyển:
\[
  P = \begin{bmatrix}
    0 & 1/2 & 0 & 1/2 & 0 \\
    1/3 & 0 & 1/3 & 0 & 1/3 \\
    0 & 1/2 & 0 & 0 & 1/2 \\
    1/2 & 0 & 0 & 0 & 1/2 \\
    0 & 1/3 & 1/3 & 1/3 & 0
  \end{bmatrix}.
\]
Trong đồ thị này, mỗi đỉnh chuyển sang các đỉnh kề với xác suất bằng nhau.

Đối với \term{bước đi ngẫu nhiên trên đồ thị} (\en{random walk on a graph}),
phân phối dừng tỷ lệ với \term{bậc} (\en{degree}) của mỗi đỉnh:
\[
  \pi_i = \frac{d_i}{2|E|},
\]
trong đó $d_i$ là bậc của đỉnh $i$ và $|E|$ là số cạnh.

Ở đây, các bậc là $d_1 = 2$, $d_2 = 3$, $d_3 = 2$, $d_4 = 2$, $d_5 = 3$,
với tổng $\sum d_i = 12$, tức $|E| = 6$. Do đó:
\[
  \pi = \left(\frac{2}{12},\; \frac{3}{12},\; \frac{2}{12},\; \frac{2}{12},\; \frac{3}{12}\right)
  = \left(\frac{1}{6},\; \frac{1}{4},\; \frac{1}{6},\; \frac{1}{6},\; \frac{1}{4}\right).
\]

Trực quan: đỉnh có nhiều ``đường vào'' (bậc cao) được ghé thăm thường xuyên hơn.
\end{example}

% ============================================================================
\section{Điều kiện cân bằng chi tiết (Detailed Balance)}
\label{sec:detailed-balance}

\textbf{Câu chuyện:} Ở trạng thái cân bằng, ``dòng người'' từ khu $i$ sang khu $j$
bằng ``dòng người'' từ khu $j$ sang khu $i$.
Giống như ở một ngã tư lúc giờ cao điểm: số xe chạy từ Đông sang Tây bằng số xe từ Tây sang Đông
--- đó là ``cân bằng chi tiết''. Nếu không bằng nhau, sẽ có kẹt xe (hệ thống chưa cân bằng).

\begin{dinhnghia}[(Điều kiện cân bằng chi tiết -- \en{Detailed balance condition})]
\label{def:detailed-balance}
Phân phối $\pi = (\pi_i)_{i\in S}$ thỏa mãn \term{điều kiện cân bằng chi tiết}
(\en{detailed balance condition}) với ma trận chuyển $P$ nếu:
\begin{equation}\label{eq:detailed-balance}
  \boxed{\pi_i P_{i,j} = \pi_j P_{j,i}, \qquad \text{với mọi } i, j \in S.}
\end{equation}
\end{dinhnghia}

\textbf{Ý nghĩa:} $\pi_i P_{i,j}$ = ``dòng xác suất'' (\en{probability flow}) từ $i$ sang $j$ ở trạng thái cân bằng.
Điều kiện này nói: dòng từ $i \to j$ bằng dòng từ $j \to i$, với mọi cặp $(i, j)$.

\begin{proposition}[(Cân bằng chi tiết $\Rightarrow$ phân phối dừng)]
\label{prop:db-implies-stationary}
Nếu $\pi$ thỏa mãn điều kiện cân bằng chi tiết \eqref{eq:detailed-balance} và $\sum_i \pi_i = 1$,
thì $\pi$ là phân phối dừng ($\pi P = \pi$).
\end{proposition}

\begin{proof}
Lấy tổng \eqref{eq:detailed-balance} theo $i$:
\[
  \sum_{i \in S} \pi_i P_{i,j} = \sum_{i \in S} \pi_j P_{j,i} = \pi_j \sum_{i \in S} P_{j,i} = \pi_j \cdot 1 = \pi_j.
\]
Đây chính là $(\pi P)_j = \pi_j$, tức $\pi P = \pi$.
\end{proof}

\textbf{Lưu ý:} Chiều ngược lại KHÔNG đúng. Có phân phối dừng không thỏa cân bằng chi tiết.
Chuỗi thỏa mãn cân bằng chi tiết được gọi là \term{thuận nghịch} (\en{reversible}).

\begin{example}[Kiểm tra cân bằng chi tiết]
Với chuỗi hai trạng thái $P = \begin{bmatrix} 1-a & a \\ b & 1-b \end{bmatrix}$ và $\pi = (b/(a+b), a/(a+b))$:
\[
  \pi_0 P_{0,1} = \frac{b}{a+b} \cdot a = \frac{ab}{a+b}, \qquad
  \pi_1 P_{1,0} = \frac{a}{a+b} \cdot b = \frac{ab}{a+b}.
\]
Bằng nhau! Chuỗi hai trạng thái luôn thỏa cân bằng chi tiết.
\end{example}

\begin{example}[Bước đi ngẫu nhiên trên đồ thị -- thuận nghịch]
\label{ex:rw-graph-reversible}
Bước đi ngẫu nhiên trên đồ thị (nhảy đều sang đỉnh kề) luôn thỏa cân bằng chi tiết với $\pi_i = d_i/(2|E|)$.

Kiểm tra: $\pi_i P_{i,j} = \frac{d_i}{2|E|} \cdot \frac{1}{d_i} = \frac{1}{2|E|}$ (nếu $i \sim j$).
Và $\pi_j P_{j,i} = \frac{d_j}{2|E|} \cdot \frac{1}{d_j} = \frac{1}{2|E|}$. Bằng nhau!
\end{example}

\begin{congthuc}
\textbf{Tìm phân phối dừng bằng cân bằng chi tiết:}

Nếu chuỗi là thuận nghịch, ta có thể tìm $\pi$ nhanh hơn bằng cách giải \eqref{eq:detailed-balance}:
\[
  \frac{\pi_j}{\pi_i} = \frac{P_{i,j}}{P_{j,i}}, \qquad \text{với mọi } i \sim j.
\]
Chọn một trạng thái tham chiếu (ví dụ $\pi_0$), tính tất cả $\pi_j$ theo $\pi_0$, rồi chuẩn hóa.
\end{congthuc}

% ============================================================================
\section{Định lý hội tụ (Convergence Theorem)}
\label{sec:convergence}

Định lý sau tổng hợp các kết quả chính về phân phối giới hạn và phân phối dừng.

\begin{dinhly}[(Hội tụ -- \en{Convergence theorem})]
\label{thm:convergence}
Cho chuỗi Markov $(X_n)_{n\in\N}$ \term{bất khả quy} (\en{irreducible}),
\term{không tuần hoàn} (\en{aperiodic}), và \term{hồi quy dương} (\en{positive recurrent})
trên không gian trạng thái $S$. Khi đó:
\begin{enumerate}[label=(\roman*)]
  \item Tồn tại \emph{duy nhất} một phân phối dừng $\pi = (\pi_j)_{j\in S}$
    thỏa mãn $\pi = \pi P$ và $\sum_{j} \pi_j = 1$.
  \item Phân phối dừng \emph{trùng} với phân phối giới hạn:
    \[
      \lim_{n\to\infty} \PP(X_n = j \mid X_0 = i) = \pi_j, \qquad \text{với mọi } i, j \in S.
    \]
  \item Phân phối dừng liên hệ với thời gian hồi quy trung bình:
    \[
      \pi_j = \frac{1}{\mu_j(j)}, \qquad \text{với mọi } j \in S.
    \]
\end{enumerate}
\end{dinhly}

\begin{tomtat}
Với chuỗi Markov bất khả quy, không tuần hoàn, hồi quy dương:
\[
  \underbrace{\text{Phân phối dừng}}_{\pi = \pi P}
  = \underbrace{\text{Phân phối giới hạn}}_{\lim P^n_{i,j}}
  = \underbrace{\text{Nghịch đảo thời gian hồi quy}}_{1/\mu_j(j)}.
\]
Ba khái niệm này \emph{trùng nhau} trong trường hợp tốt nhất.
\end{tomtat}

\textbf{Lưu ý về không gian trạng thái hữu hạn:}
Nếu $S$ hữu hạn và chuỗi bất khả quy, thì chuỗi \emph{tự động} hồi quy dương.
Do đó, với chuỗi Markov hữu hạn bất khả quy và không tuần hoàn,
phân phối dừng tồn tại, duy nhất, và bằng phân phối giới hạn.

% ============================================================================
\section{Tốc độ hội tụ (Rate of Convergence)}
\label{sec:convergence-rate}

\textbf{Câu chuyện:} Ta biết $P^n_{i,j} \to \pi_j$ khi $n \to \infty$.
Nhưng ``nhanh'' thế nào? Cần bao nhiêu bước để ``quên'' trạng thái ban đầu?

\begin{dinhly}[(Tốc độ hội tụ hình học)]
\label{thm:convergence-rate}
Cho chuỗi Markov hữu hạn bất khả quy và không tuần hoàn với ma trận chuyển $P$ có
các trị riêng $1 = \lambda_1 > |\lambda_2| \geq |\lambda_3| \geq \cdots \geq |\lambda_N|$.
Khi đó tồn tại hằng số $C > 0$ sao cho:
\begin{equation}\label{eq:convergence-rate}
  |P^n_{i,j} - \pi_j| \leq C\, |\lambda_2|^n, \qquad \text{với mọi } i, j \in S, \; n \geq 0.
\end{equation}
\end{dinhly}

\textbf{Ý nghĩa:} Tốc độ hội tụ được quyết định bởi $|\lambda_2|$ --- trị riêng lớn thứ hai (về giá trị tuyệt đối).
\begin{itemize}[nosep]
  \item $|\lambda_2|$ gần 0: hội tụ rất nhanh (vài bước).
  \item $|\lambda_2|$ gần 1: hội tụ chậm (nhiều bước).
\end{itemize}

\textbf{Khoảng cách phổ} (\en{spectral gap}): $\gamma = 1 - |\lambda_2|$.
$\gamma$ lớn $\Rightarrow$ hội tụ nhanh. Số bước cần thiết xấp xỉ $\sim 1/\gamma$.

\begin{example}[Tốc độ hội tụ chuỗi hai trạng thái]
Với $P = \begin{bmatrix} 1-a & a \\ b & 1-b \end{bmatrix}$:
$\lambda_2 = 1 - a - b$, nên $|\lambda_2| = |1-a-b|$.

\begin{itemize}[nosep]
  \item $a = b = 0.5$: $|\lambda_2| = 0$, $\gamma = 1$. Hội tụ trong 1 bước!
  \item $a = b = 0.1$: $|\lambda_2| = 0.8$, $\gamma = 0.2$. Cần $\sim 5$ bước.
  \item $a = b = 0.01$: $|\lambda_2| = 0.98$, $\gamma = 0.02$. Cần $\sim 50$ bước.
\end{itemize}
\end{example}

% ============================================================================
\section{Bảng tổng hợp phương pháp (Summary of Methods)}
\label{sec:summary-methods}

\begin{congthuc}
\textbf{Bảng tổng hợp các phương pháp tính toán cho chuỗi Markov:}

\medskip
\begin{center}
\renewcommand{\arraystretch}{1.5}
\begin{tabular}{|l|l|}
  \hline
  \textbf{Đại lượng cần tính} & \textbf{Phương pháp} \\
  \hline
  Kỳ vọng $\E[X]$ &
    $\displaystyle\sum_{x} x \cdot \PP(X = x)$ \\
  \hline
  Xác suất chạm (\en{hitting prob.}) $g(k)$ &
    Giải $g = Pg$ trên trạng thái không hấp thụ \\
    & với điều kiện biên $g = 0$ hoặc $g = 1$ \\
  \hline
  Thời gian chạm trung bình $h(k)$ &
    Giải $h = 1 + Ph$ trên trạng thái không hấp thụ \\
    & với điều kiện biên $h = 0$ \\
  \hline
  Phân phối dừng $\pi$ &
    Giải $\pi = \pi P$ với $\displaystyle\sum_i \pi_i = 1$ \\
  \hline
  Phân phối dừng (thuận nghịch) &
    Dùng $\pi_i P_{i,j} = \pi_j P_{j,i}$ \\
  \hline
\end{tabular}
\end{center}
\end{congthuc}

\begin{tomtat}
Tất cả các bài toán trên đều quy về \textbf{hệ phương trình tuyến tính}.
\begin{itemize}[nosep]
  \item Xác suất chạm và thời gian chạm: giải hệ với điều kiện biên tại trạng thái hấp thụ.
  \item Phân phối dừng: giải hệ thuần nhất $\pi(P - I) = 0$ với ràng buộc chuẩn hóa $\sum \pi_i = 1$.
  \item Nếu chuỗi thuận nghịch: dùng cân bằng chi tiết --- đơn giản hơn, không cần giải hệ.
\end{itemize}
\end{tomtat}

% ============================================================================
\section{Bài tập (Exercises)}
\label{sec:ch7-exercises}

\begin{exercise}
\label{ex:ch7-1}
Cho chuỗi Markov trên $S = \{0, 1, 2\}$ với ma trận chuyển:
\[
  P = \begin{bmatrix}
    1/2 & 1/4 & 1/4 \\
    1/3 & 1/3 & 1/3 \\
    1/4 & 1/4 & 1/2
  \end{bmatrix}.
\]
\begin{enumerate}[label=(\alph*)]
  \item Kiểm tra rằng $P$ là ma trận chính quy.
  \item Tìm phân phối dừng $\pi$.
  \item Tính thời gian hồi quy trung bình $\mu_i(i)$ với mỗi $i \in S$.
\end{enumerate}
\end{exercise}

\begin{exercise}
\label{ex:ch7-2}
Xét chuỗi Markov hai trạng thái với ma trận chuyển:
\[
  P = \begin{bmatrix} 0.7 & 0.3 \\ 0.4 & 0.6 \end{bmatrix}.
\]
\begin{enumerate}[label=(\alph*)]
  \item Tính $P^n$ bằng phương pháp chéo hóa.
  \item Tìm $\lim_{n\to\infty} P^n$ và xác định phân phối giới hạn.
  \item Xác minh rằng phân phối giới hạn thỏa mãn $\pi = \pi P$.
  \item Xác minh điều kiện cân bằng chi tiết.
\end{enumerate}
\end{exercise}

\begin{exercise}
\label{ex:ch7-3}
Cho chuỗi Markov trên $S = \{0, 1\}$ với $P = \begin{bmatrix} 0 & 1 \\ 1 & 0 \end{bmatrix}$.
\begin{enumerate}[label=(\alph*)]
  \item Tìm phân phối dừng $\pi$.
  \item Chuỗi có phân phối giới hạn không? Giải thích.
  \item Tính $P^n$ cho $n$ chẵn và $n$ lẻ. Chu kỳ của chuỗi bằng bao nhiêu?
\end{enumerate}
\end{exercise}

\begin{exercise}
\label{ex:ch7-4}
Chuỗi Ehrenfest với $N = 3$ viên bi có ma trận chuyển:
\[
  P = \begin{bmatrix}
    0 & 1 & 0 & 0 \\
    1/3 & 0 & 2/3 & 0 \\
    0 & 2/3 & 0 & 1/3 \\
    0 & 0 & 1 & 0
  \end{bmatrix}.
\]
\begin{enumerate}[label=(\alph*)]
  \item Chuỗi này có bất khả quy không? Có tuần hoàn không?
  \item Tìm phân phối dừng $\pi$.
  \item So sánh $\pi$ với phân phối nhị thức $\mathrm{Bin}(3, 1/2)$. Nhận xét?
  \item Kiểm tra điều kiện cân bằng chi tiết.
\end{enumerate}
\end{exercise}

\begin{exercise}
\label{ex:ch7-5}
Cho chuỗi Markov trên $S = \{1, 2, 3, 4\}$ với ma trận chuyển:
\[
  P = \begin{bmatrix}
    0 & 1/2 & 1/2 & 0 \\
    1/2 & 0 & 0 & 1/2 \\
    1/3 & 0 & 1/3 & 1/3 \\
    0 & 1/2 & 1/2 & 0
  \end{bmatrix}.
\]
\begin{enumerate}[label=(\alph*)]
  \item Tìm phân phối dừng $\pi$ bằng cách giải hệ $\pi = \pi P$.
  \item Sử dụng kết quả $\pi_j = 1/\mu_j(j)$ để tính thời gian hồi quy trung bình
    của mỗi trạng thái.
  \item Trạng thái nào được ``ghé thăm'' nhiều nhất trong dài hạn?
\end{enumerate}
\end{exercise}

\begin{exercise}
\label{ex:ch7-6}
Cho chuỗi Markov trên $S = \{0, 1, 2\}$ với ma trận chuyển:
\[
  P = \begin{bmatrix}
    0 & 2/3 & 1/3 \\
    1/2 & 0 & 1/2 \\
    1/3 & 2/3 & 0
  \end{bmatrix}.
\]
\begin{enumerate}[label=(\alph*)]
  \item Tìm phân phối dừng $\pi$.
  \item Kiểm tra xem chuỗi có thỏa mãn điều kiện cân bằng chi tiết không.
  \item Tìm trị riêng thứ hai $\lambda_2$ của $P$ và ước lượng tốc độ hội tụ.
\end{enumerate}
\end{exercise}

\begin{exercise}
\label{ex:ch7-7}
Chứng minh rằng bước đi ngẫu nhiên trên đồ thị vô hướng liên thông (nhảy đều sang đỉnh kề)
luôn thỏa mãn điều kiện cân bằng chi tiết với $\pi_i = d_i / (2|E|)$.
\end{exercise}

\chapter{Chuỗi Markov thời gian liên tục\\{\normalsize\textit{Continuous-Time Markov Chains}}}
\label{ch:continuous}

% ============================================================================
\section{Quá trình Poisson (The Poisson Process)}
\label{sec:poisson}

\textbf{Câu chuyện:} Hãy tưởng tượng bạn ngồi ở quán café quan sát xe buýt đến trạm.
Trung bình cứ 10 phút có một chuyến, nhưng chính xác khi nào thì xe đến --- bạn không biết trước.
Có khi 5 phút đã có xe, có khi phải đợi 15 phút. Đây chính là bản chất của \term{quá trình Poisson}:
sự kiện xảy ra ``ngẫu nhiên'' theo thời gian, với một tốc độ trung bình cố định.

Trong chương này ta bắt đầu nghiên cứu các \term{quá trình thời gian liên tục} (\en{continuous-time processes}),
là các họ $(X_t)_{t\in\R_+}$ gồm các biến ngẫu nhiên được đánh chỉ số bởi $\R_+$.
Khác với chuỗi Markov rời rạc (thời gian nhảy tại $n = 0, 1, 2, \ldots$),
ở đây thời gian chạy liên tục và quá trình có thể ``nhảy'' vào bất kỳ lúc nào.

\begin{tomtat}
\textbf{Ý tưởng chính:} Quá trình Poisson đếm số sự kiện xảy ra theo thời gian.
Nó là ``viên gạch nền tảng'' cho mọi chuỗi Markov liên tục,
giống như bước đi ngẫu nhiên là nền tảng cho chuỗi Markov rời rạc.
\end{tomtat}

\begin{dinhnghia}[(Quá trình Poisson chuẩn -- \en{Standard Poisson process})]
\term{Quá trình Poisson chuẩn} (\en{standard Poisson process}) là quá trình ngẫu nhiên $(N_t)_{t\in\R_+}$ chỉ có bước nhảy $+1$, và đường đi (\en{paths}) không đổi giữa các bước nhảy. Tại thời điểm $t$:
\[
  N_t = \sum_{k=1}^{\infty} \ind_{(T_k, \infty)}(t), \qquad t \in \R_+,
\]
trong đó $(T_k)_{k\geq 1}$ là dãy tăng các thời điểm nhảy với $\lim_{k\to\infty} T_k = +\infty$.

$(N_t)_{t\in\R_+}$ thỏa mãn:
\begin{enumerate}[nosep]
  \item \textbf{Gia số độc lập} (\en{independence of increments}): với mọi $0 \leq t_0 < t_1 < \cdots < t_n$, các biến $N_{t_1} - N_{t_0}, \ldots, N_{t_n} - N_{t_{n-1}}$ là độc lập.
  \item \textbf{Gia số dừng} (\en{stationarity of increments}): $N_{t+h} - N_{s+h}$ có cùng phân phối với $N_t - N_s$ với mọi $h > 0$ và $0 \leq s \leq t$.
\end{enumerate}
\end{dinhnghia}

\textbf{Ví dụ đời thường:}
\begin{itemize}[nosep]
  \item Số cuộc gọi đến tổng đài trong khoảng thời gian $[0, t]$.
  \item Số khách hàng vào cửa hàng từ 8h sáng đến giờ $t$.
  \item Số hạt phóng xạ phát ra bởi một mẫu chất trong thời gian $t$.
\end{itemize}

\begin{dinhly}[(Định lý 9 -- Phân phối Poisson)]
Giả sử quá trình đếm $(N_t)_{t\in\R_+}$ thỏa mãn các điều kiện (1) và (2). Khi đó:
\[
  \PP(N_t - N_s = k) = e^{-\lambda(t-s)} \frac{(\lambda(t-s))^k}{k!}, \qquad k \in \N, \quad 0 \leq s \leq t,
\]
với hằng số $\lambda > 0$.
\end{dinhly}

Tham số $\lambda > 0$ được gọi là \term{cường độ} (\en{intensity}) của quá trình:
\begin{equation}\label{eq:intensity}
  \lambda := \lim_{h \to 0} \frac{1}{h} \PP(N_h = 1).
\end{equation}

Đặc biệt, $N_t$ có \term{phân phối Poisson} (\en{Poisson distribution}) với tham số $\lambda t$:
\[
  \PP(N_t = k) = \frac{(\lambda t)^k}{k!}\, e^{-\lambda t}, \qquad t > 0.
\]

\textbf{Xác suất chuyển vi mô} (\en{infinitesimal transition probabilities}):
\begin{align}
  \PP(N_{t+h} - N_t = 1) &\simeq \lambda h, \qquad h \to 0, \label{eq:poisson-inf-1}\\
  \PP(N_{t+h} - N_t = 0) &\simeq 1 - \lambda h, \qquad h \to 0. \label{eq:poisson-inf-0}
\end{align}

Điều này có nghĩa: trong khoảng thời gian ``ngắn'' $[t, t+h]$, gia số $N_{t+h} - N_t$ xử sự như biến ngẫu nhiên Bernoulli với tham số $\lambda h$.

\begin{proposition}[(Phân phối thời gian nhảy)]
Luật của $T_n$ có mật độ $t \mapsto \lambda^n e^{-\lambda t} \frac{t^{n-1}}{(n-1)!}$ trên $\R_+$, $n \geq 1$.
\end{proposition}

Thời gian $\tau_k := T_{k+1} - T_k$ giữa các bước nhảy tạo thành dãy các biến ngẫu nhiên \term{độc lập cùng phân phối} (\en{i.i.d.}) có phân phối mũ với tham số $\lambda > 0$:
\[
  \PP(\tau_0 > t_0, \ldots, \tau_n > t_n) = e^{-\lambda(t_0 + \cdots + t_n)} = \prod_{k=0}^n e^{-\lambda t_k},
\]
và $\E[\tau_k] = 1/\lambda$.

\begin{example}[Xe buýt đến trạm]
Giả sử xe buýt đến trạm theo quá trình Poisson với cường độ $\lambda = 6$ chuyến/giờ (trung bình 10 phút/chuyến).

\textbf{Câu hỏi 1:} Xác suất không có xe nào trong 15 phút?
\[
  \PP(N_{0.25} = 0) = e^{-6 \times 0.25} = e^{-1.5} \approx 0.223.
\]

\textbf{Câu hỏi 2:} Xác suất có đúng 2 xe trong 20 phút?
\[
  \PP(N_{1/3} = 2) = e^{-2} \frac{2^2}{2!} = 2e^{-2} \approx 0.271.
\]

\textbf{Câu hỏi 3:} Thời gian đợi trung bình cho xe tiếp theo?
\[
  \E[\tau] = \frac{1}{\lambda} = \frac{1}{6} \text{ giờ} = 10 \text{ phút}.
\]
\end{example}

\begin{tomtat}
Quá trình Poisson: gia số độc lập và dừng, $N_t \sim \text{Poisson}(\lambda t)$.
Cường độ $\lambda$ cao hơn $\Rightarrow$ nhảy thường xuyên hơn $\Rightarrow$ thời gian ở mỗi trạng thái ngắn hơn ($\E[\tau_k] = 1/\lambda$).

Tính chất ``không nhớ'' (\en{memoryless}): dù bạn đã đợi 5 phút rồi, thời gian đợi thêm vẫn có phân phối mũ --- quá khứ không ảnh hưởng tương lai.
\end{tomtat}

% ============================================================================
\section{Quá trình sinh-tử (Birth and Death Processes)}
\label{sec:birth-death}

\textbf{Câu chuyện:} Hãy tưởng tượng một quán café. Khách đến (``sinh'') và khách đi (``tử'').
Nếu quán đang có $i$ khách:
\begin{itemize}[nosep]
  \item Khách mới đến với tốc độ $\lambda_i$ (số khách tăng lên 1).
  \item Một khách hiện tại rời đi với tốc độ $\mu_i$ (số khách giảm đi 1).
\end{itemize}
Số khách trong quán thay đổi liên tục theo thời gian --- đây là \term{quá trình sinh-tử}.

\begin{dinhnghia}[(Quá trình sinh thuần túy -- \en{Pure birth process})]
Chuỗi Markov thời gian liên tục $(X_t)_{t\in\R_+}$ thỏa mãn:
\[
  \PP(X_{t+h} - X_t = 1 \mid X_t = i) \simeq \lambda_i h, \qquad h \to 0,\; i \in S,
\]
\[
  \PP(X_{t+h} - X_t = 0 \mid X_t = i) \simeq 1 - \lambda_i h, \qquad h \to 0,\; i \in S,
\]
được gọi là \term{quá trình sinh thuần túy} (\en{pure birth process}) với tốc độ sinh $\lambda_i \geq 0$.
Quá trình Poisson là trường hợp đặc biệt với $\lambda_n = \lambda$ không phụ thuộc trạng thái.
\end{dinhnghia}

\begin{dinhnghia}[(Quá trình tử thuần túy -- \en{Pure death process})]
Tương tự, $(X_t)_{t\in\R_+}$ với:
\[
  \PP(X_{t+h} - X_t = -1 \mid X_t = i) \simeq \mu_i h, \quad
  \PP(X_{t+h} - X_t = 0 \mid X_t = i) \simeq 1 - \mu_i h,
\]
là \term{quá trình tử thuần túy} (\en{pure death process}) với tốc độ tử $\mu_i \geq 0$.
\end{dinhnghia}

\begin{dinhnghia}[(Quá trình sinh-tử -- \en{Birth and death process})]
Chuỗi Markov thời gian liên tục $(X_t)_{t\in\R_+}$ với:
\begin{align*}
  \PP(X_{t+h} - X_t = 1 \mid X_t = i)  &\simeq \lambda_i h, \\
  \PP(X_{t+h} - X_t = -1 \mid X_t = i) &\simeq \mu_i h, \\
  \PP(X_{t+h} - X_t = 0 \mid X_t = i)  &\simeq 1 - (\lambda_i + \mu_i)h,
\end{align*}
khi $h \to 0$, được gọi là \term{quá trình sinh-tử} (\en{birth and death process}) với tốc độ sinh $\lambda_i \geq 0$ và tốc độ tử $\mu_i \geq 0$, $i \in S$.
\end{dinhnghia}

Thời gian $\tau_i$ ở trạng thái $i$ trước khi chuyển sang trạng thái khác là biến ngẫu nhiên có phân phối mũ:
\[
  \tau_i = \min(\tau_{i,i+1},\; \tau_{i,i-1}),
\]
với tham số $\lambda_i + \mu_i$, và $\E[\tau_i] = \frac{1}{\lambda_i + \mu_i}$.

\begin{example}[Quán café có sức chứa $N = 4$]
Xét quán café có tối đa $4$ chỗ ngồi. Khách đến với tốc độ $\lambda = 2$ người/giờ,
khách rời đi với tốc độ $\mu = 3$ người/giờ (mỗi khách ngồi trung bình 20 phút).

Khi quán có $i$ khách ($0 \leq i \leq 4$):
\begin{itemize}[nosep]
  \item Tốc độ sinh: $\lambda_i = \lambda = 2$ (nếu $i < 4$), $\lambda_4 = 0$ (quán đầy).
  \item Tốc độ tử: $\mu_i = i \cdot \mu = 3i$ (mỗi khách rời đi độc lập với tốc độ $\mu$).
\end{itemize}

Thời gian trung bình ở trạng thái $i = 2$:
\[
  \E[\tau_2] = \frac{1}{\lambda_2 + \mu_2} = \frac{1}{2 + 6} = \frac{1}{8} \text{ giờ} = 7.5 \text{ phút}.
\]
\end{example}

% ============================================================================
\section{Chuỗi Markov thời gian liên tục (Continuous-Time Markov Chains)}
\label{sec:ctmc-definition}

\textbf{Câu chuyện:} Giống như con kiến trên bàn cờ (Chương~\ref{ch:discrete}) chỉ quan tâm ô hiện tại,
chuỗi Markov liên tục cũng ``không nhớ'' --- nhưng bây giờ con kiến di chuyển trong thời gian liên tục.
Nó ở một ô trong thời gian ngẫu nhiên (phân phối mũ), rồi nhảy sang ô khác.

\begin{dinhnghia}[(Tính chất Markov liên tục -- \en{Continuous-time Markov property})]
Quá trình ngẫu nhiên $(X_t)_{t\in\R_+}$ nhận giá trị trong $\Z$ có \term{tính chất Markov} nếu với mọi
$0 < s_1 < \cdots < s_{n-1} < s < t$:
\[
  \PP(X_t = j \mid X_s = i_n,\; X_{s_{n-1}} = i_{n-1},\; \ldots,\; X_{s_1} = i_0) = \PP(X_t = j \mid X_s = i_n).
\]
\end{dinhnghia}

\textbf{Ví dụ:} Quá trình Poisson là chuỗi Markov thời gian liên tục vì nó có gia số độc lập.
Tổng quát hơn, mọi quá trình có gia số độc lập đều là chuỗi Markov.

\subsection{Nửa nhóm chuyển (Transition Semigroup)}

Xác suất chuyển phụ thuộc thời gian:
\[
  P_{i,j}(t) := \PP(X_{t+s} = j \mid X_s = i), \qquad i,j \in \Z,\; t \in \R_+,
\]
trong đó ta giả sử quá trình là \term{đồng nhất theo thời gian} (\en{time homogeneous}): $\PP(X_{t+s} = j \mid X_s = i)$ không phụ thuộc $s$. Lưu ý $P(0) = I_d$.

Ma trận $P(t) = [P_{i,j}(t)]_{i,j\in S}$ được gọi là \term{nửa nhóm chuyển} (\en{transition semigroup}).

\begin{congthuc}
\textbf{Tính chất nửa nhóm} (\en{semigroup property}):
\begin{equation}\label{eq:semigroup}
  P(s+t) = P(s)\,P(t) = P(t)\,P(s), \qquad s, t \in \R_+.
\end{equation}
Đây chính là phương trình \term{Chapman--Kolmogorov} phiên bản thời gian liên tục.
\end{congthuc}

% ============================================================================
\section{Ma trận sinh vô cùng bé (Infinitesimal Generator)}
\label{sec:generator}

\textbf{Câu chuyện:} Trong trường hợp rời rạc, ma trận chuyển $P$ cho bạn biết ``xác suất đi đâu sau 1 bước''.
Nhưng trong thời gian liên tục, không có ``1 bước'' --- thời gian chạy liên tục!
Thay vào đó, ta cần biết hệ thống thay đổi \textbf{nhanh thế nào} --- đây là vai trò của ma trận sinh $Q$.

Hãy nghĩ $Q$ như ``tốc độ'' thay vì ``xác suất'':
\begin{itemize}[nosep]
  \item $\lambda_{i,j} > 0$ (ngoài đường chéo): tốc độ chuyển từ $i$ sang $j$. Càng lớn $\Rightarrow$ càng nhanh.
  \item $\lambda_{i,i} < 0$ (đường chéo): tốc độ rời khỏi $i$. $|\lambda_{i,i}|$ lớn $\Rightarrow$ ở $i$ rất ngắn.
\end{itemize}

Bằng cách lấy đạo hàm quan hệ nửa nhóm \eqref{eq:semigroup} theo $t$ rồi cho $t = 0$:
\[
  P'(t) = \lim_{s \to 0} \frac{P(t+s) - P(t)}{s} = P(t) \lim_{s \to 0} \frac{P(s) - P(0)}{s} = P(t)\,Q,
\]
trong đó:

\begin{dinhnghia}[(Ma trận sinh -- \en{Infinitesimal generator})]
\[
  Q := P'(0) = \lim_{s \to 0} \frac{P(s) - P(0)}{s}
\]
được gọi là \term{ma trận sinh vô cùng bé} (\en{infinitesimal generator}) của $(X_t)_{t\in\R_+}$.
\end{dinhnghia}

Ký hiệu $Q = [\lambda_{i,j}]_{i,j\in S}$. Các hàng của $Q$ luôn có tổng bằng 0:
\[
  \sum_{j \in S} \lambda_{i,j} = 0, \qquad i \in S.
\]
Do đó: $\lambda_{i,i} = -\sum_{l \neq i} \lambda_{i,l}$ (phần tử đường chéo luôn âm hoặc bằng 0).

\begin{congthuc}
\textbf{Phương trình Kolmogorov tiến} (\en{forward Kolmogorov equation}):
\begin{equation}\label{eq:forward-kolmogorov}
  P'(t) = P(t)\,Q, \qquad t > 0.
\end{equation}
\textbf{Phương trình Kolmogorov lùi} (\en{backward Kolmogorov equation}):
\begin{equation}\label{eq:backward-kolmogorov}
  P'(t) = Q\,P(t), \qquad t > 0.
\end{equation}
\textbf{Nghiệm} (dùng \term{hàm mũ ma trận} -- \en{matrix exponential}):
\[
  P(t) = P(0)\,e^{tQ} = e^{tQ}, \qquad \text{trong đó } e^{tQ} = I_d + \sum_{n=1}^{\infty} \frac{t^n}{n!} Q^n.
\]
\end{congthuc}

\textbf{Xấp xỉ bậc nhất} khi $h \to 0$:
\[
  P(h) \simeq I_d + hQ + o(h), \qquad h \to 0.
\]

Xác suất chuyển vi mô:
\[
  P_{i,j}(h) = \PP(X_{t+h} = j \mid X_t = i) \simeq
  \begin{cases}
    \lambda_{i,j} h, & i \neq j, \\[3pt]
    1 + \lambda_{i,i} h, & i = j,
  \end{cases}
  \qquad h \to 0.
\]

\textbf{Ý nghĩa vật lý:} Quá trình chuyển từ trạng thái $i$ sang $j \neq i$ giống như quá trình Poisson với cường độ $\lambda_{i,j}$. Thời gian $\tau_{i,j}$ trước khi chuyển từ $i$ sang $j$ có phân phối mũ:
\begin{equation}\label{eq:sojourn-exp}
  \PP(\tau_{i,j} > t) = e^{-\lambda_{i,j} t}, \qquad t \in \R_+, \quad i \neq j.
\end{equation}
và $\E[\tau_{i,j}] = 1/\lambda_{i,j}$.

Thời gian ở trạng thái $i$ (trước khi rời đi bất kỳ đâu):
\[
  \tau_i = \min_{j \in S,\, j \neq i} \tau_{i,j},
\]
là biến mũ với tham số $\sum_{j \neq i} \lambda_{i,j} = -\lambda_{i,i}$, nên $\E[\tau_i] = -1/\lambda_{i,i}$.

\subsection{Ví dụ các ma trận sinh}

\begin{example}[Quá trình Poisson]
Ma trận sinh:
\[
  Q = [\lambda_{i,j}]_{i,j\in\N} = \begin{bmatrix}
    -\lambda & \lambda & 0 & \cdots \\
    0 & -\lambda & \lambda & \cdots \\
    0 & 0 & -\lambda & \cdots \\
    \vdots & \vdots & \vdots & \ddots
  \end{bmatrix}.
\]
Thời gian ở mỗi trạng thái: $\E[\tau_i] = 1/\lambda$ (không phụ thuộc $i$).
\end{example}

\begin{example}[Quá trình sinh-tử trên $\{0,1,\ldots,N\}$]
Ma trận sinh (với $\mu_0 = \lambda_N = 0$):
\[
  Q = \begin{bmatrix}
    -\lambda_0 & \lambda_0 & 0 & \cdots & 0 \\
    \mu_1 & -\lambda_1-\mu_1 & \lambda_1 & \cdots & 0 \\
    \vdots & \ddots & \ddots & \ddots & \vdots \\
    0 & \cdots & \mu_{N-1} & -\lambda_{N-1}-\mu_{N-1} & \lambda_{N-1} \\
    0 & \cdots & 0 & \mu_N & -\mu_N
  \end{bmatrix}.
\]
\end{example}

\begin{example}[Ma trận sinh 5 trạng thái -- tính toán chi tiết]
\label{ex:5state-generator}
Xét quá trình Markov liên tục trên $S = \{0, 1, 2, 3, 4\}$ với ma trận sinh:
\[
  Q = \begin{bmatrix}
    -3 & 2 & 1 & 0 & 0 \\
    1 & -4 & 2 & 1 & 0 \\
    0 & 1 & -3 & 1 & 1 \\
    0 & 0 & 2 & -5 & 3 \\
    0 & 0 & 0 & 4 & -4
  \end{bmatrix}.
\]

\textbf{Đọc thông tin từ $Q$:}
\begin{itemize}[nosep]
  \item Từ trạng thái $0$: chuyển sang $1$ với tốc độ $2$, sang $2$ với tốc độ $1$. Tổng tốc độ rời = $3$.
  \item Thời gian trung bình ở trạng thái $0$: $\E[\tau_0] = 1/3$ (đơn vị thời gian).
  \item Thời gian trung bình ở trạng thái $3$: $\E[\tau_3] = 1/5$ (trạng thái ``bận rộn'' nhất).
  \item Từ trạng thái $4$: chỉ có thể quay lại $3$ (tốc độ $4$), $\E[\tau_4] = 1/4$.
\end{itemize}
\end{example}

\begin{tomtat}
Ma trận sinh $Q$ đóng vai trò tương tự ma trận chuyển $P$ trong trường hợp rời rạc.
Quan hệ: $P(h) \approx I_d + hQ$ cho $h$ nhỏ --- ``rời rạc hóa'' chuỗi liên tục.
Khác biệt chính: phần tử ngoài đường chéo $\lambda_{i,j} \geq 0$ (tốc độ), phần tử đường chéo $\lambda_{i,i} \leq 0$.

\medskip
\textbf{So sánh $P$ (rời rạc) vs $Q$ (liên tục):}
\begin{center}
\renewcommand{\arraystretch}{1.3}
\begin{tabular}{lcc}
\hline
 & \textbf{Ma trận $P$} & \textbf{Ma trận $Q$} \\
\hline
Phần tử ngoài đường chéo & $\geq 0$ (xác suất) & $\geq 0$ (tốc độ) \\
Phần tử đường chéo & $\geq 0$ & $\leq 0$ \\
Tổng mỗi hàng & $= 1$ & $= 0$ \\
Ý nghĩa & ``đi đâu'' & ``đi nhanh thế nào'' \\
\hline
\end{tabular}
\end{center}
\end{tomtat}

% ============================================================================
\section{Chuỗi Markov hai trạng thái liên tục (Two-State Continuous-Time Chain)}
\label{sec:two-state-continuous}

Xét quá trình Markov liên tục trên $S = \{0, 1\}$ với ma trận sinh:
\[
  Q = \begin{bmatrix} -\alpha & \alpha \\ \beta & -\beta \end{bmatrix}, \qquad \alpha, \beta \geq 0.
\]

\textbf{Ví dụ:} Một máy tính có thể ở trạng thái ``hoạt động'' (0) hoặc ``hỏng'' (1).
Máy hỏng với tốc độ $\alpha$ và được sửa với tốc độ $\beta$.

Phương trình Kolmogorov tiến $P'(t) = P(t) Q$ cho hệ phương trình vi phân:
\[
  \begin{cases}
    P'_{0,0}(t) = -\alpha P_{0,0}(t) + \beta P_{0,1}(t), \\
    P'_{0,1}(t) = \alpha P_{0,0}(t) - \beta P_{0,1}(t),
  \end{cases}
\]
với điều kiện đầu $P(0) = I_d$.

\textbf{Nghiệm bằng hàm mũ ma trận:}

Ma trận $Q$ có trị riêng $\lambda_1 = 0$ và $\lambda_2 = -(\alpha + \beta)$.

Chéo hóa và tính $e^{tQ}$:
\begin{equation}\label{eq:Pt-two-state}
  P(t) = e^{tQ} = \frac{1}{\alpha+\beta}
  \begin{bmatrix}
    \beta + \alpha e^{-t(\alpha+\beta)} & \alpha - \alpha e^{-t(\alpha+\beta)} \\
    \beta - \beta e^{-t(\alpha+\beta)} & \alpha + \beta e^{-t(\alpha+\beta)}
  \end{bmatrix}, \quad t > 0.
\end{equation}

\textbf{So sánh với trường hợp rời rạc:} Biểu thức \eqref{eq:Pt-two-state} tương tự công thức rời rạc (Chương~\ref{ch:discrete}), nhưng thay $\lambda_2^n = (1-a-b)^n$ bằng $e^{-t(\alpha+\beta)}$. Sự hội tụ theo hàm mũ thay vì lũy thừa.

\subsection{Phân phối giới hạn}

Khi $t \to \infty$ (với $\alpha > 0$ hoặc $\beta > 0$):
\begin{align}
  \lim_{t\to\infty} \PP(X_t = 1 \mid X_0 = 0) &= \lim_{t\to\infty} \PP(X_t = 1 \mid X_0 = 1) = \frac{\alpha}{\alpha+\beta}, \label{eq:ct-limit-1}\\
  \lim_{t\to\infty} \PP(X_t = 0 \mid X_0 = 0) &= \lim_{t\to\infty} \PP(X_t = 0 \mid X_0 = 1) = \frac{\beta}{\alpha+\beta}. \label{eq:ct-limit-0}
\end{align}

Phân phối giới hạn:
\begin{equation}\label{eq:ct-pi}
  \pi = (\pi_0, \pi_1) = \left(\frac{\beta}{\alpha+\beta},\; \frac{\alpha}{\alpha+\beta}\right).
\end{equation}

\begin{tomtat}
Chuỗi hai trạng thái liên tục có cùng phân phối giới hạn $\pi = (\beta/(\alpha+\beta),\; \alpha/(\alpha+\beta))$ như trường hợp rời rạc (thay $a \to \alpha$, $b \to \beta$).
Sự hội tụ luôn xảy ra trong thời gian liên tục (trừ $\alpha = \beta = 0$), và nhanh hơn khi $\alpha + \beta$ lớn hơn.
\end{tomtat}

% ============================================================================
\section{Phân phối giới hạn và dừng (Limiting and Stationary Distributions)}
\label{sec:ct-stationary}

\textbf{Câu chuyện:} Quay lại ví dụ quán café. Sau nhiều giờ hoạt động, tỷ lệ thời gian quán có $0, 1, 2, \ldots$ khách
sẽ ổn định --- không phụ thuộc vào lúc mở cửa có bao nhiêu khách.
Đây chính là phân phối dừng: ``dân số ổn định'' sau thời gian dài.

\begin{dinhnghia}[(Phân phối dừng liên tục -- \en{Stationary distribution, continuous-time})]
Phân phối $\pi = (\pi_i)_{i\in S}$ được gọi là \term{dừng} (\en{stationary}) nếu $\pi P(t) = \pi$ với mọi $t \in \R_+$.
\end{dinhnghia}

\begin{proposition}[(Mệnh đề 4)]
Phân phối $\pi = (\pi_i)_{i\in S}$ là dừng khi và chỉ khi:
\begin{equation}\label{eq:pi-Q}
  \pi Q = 0.
\end{equation}
\end{proposition}

\textit{Chứng minh.} Nếu $\pi Q = 0$, thì $\pi P(t) = \pi e^{tQ} = \pi \sum_{n=0}^{\infty} \frac{t^n}{n!} Q^n = \pi$.
Ngược lại, lấy đạo hàm $\pi = \pi P(t)$ tại $t = 0$: $0 = \pi P'(0) = \pi Q P(0) = \pi Q$. \qed

\begin{proposition}[(Mệnh đề 5)]
Nếu chuỗi Markov liên tục $(X_t)_{t\in\R_+}$ có phân phối giới hạn
\[
  \pi_j := \lim_{t\to\infty} \PP(X_t = j \mid X_0 = i) = \lim_{t\to\infty} P_{i,j}(t), \qquad j \in S,
\]
không phụ thuộc vào trạng thái ban đầu $i \in S$, thì $\pi Q = 0$, tức $\pi$ là phân phối dừng.
\end{proposition}

\begin{congthuc}
\textbf{Tìm phân phối dừng:} Giải hệ $\pi Q = 0$ kết hợp $\sum_i \pi_i = 1$.

Tương đương: $\sum_{i \in S} \pi_i \lambda_{i,j} = 0$ với mọi $j \in S$.

Lưu ý: điều kiện $\pi Q = 0$ tương đương với $\pi = \pi(I_d + hQ)$ cho mọi $h > 0$, chính là phân phối dừng của chuỗi rời rạc có ma trận chuyển $P(h) \simeq I_d + hQ$.
\end{congthuc}

\subsection{Phân phối dừng của quá trình sinh-tử}

Với quá trình sinh-tử trên $\N$ có ma trận sinh $Q = [\lambda_{i,j}]$, hệ $\pi Q = 0$ cho:
\[
  \begin{cases}
    0 = -\lambda_0 \pi_0 + \mu_1 \pi_1 \\
    0 = \lambda_{j-1}\pi_{j-1} - (\lambda_j + \mu_j)\pi_j + \mu_{j+1}\pi_{j+1}, \quad j \geq 1
  \end{cases}
\]
Giải đệ quy:
\[
  \pi_1 = \frac{\lambda_0}{\mu_1}\pi_0, \qquad
  \pi_{j+1} = \frac{\lambda_j \cdots \lambda_0}{\mu_{j+1} \cdots \mu_1} \pi_0, \qquad j \geq 0,
\]
với $\pi_0$ được xác định bởi $\sum_{i=0}^{\infty} \pi_i = 1$:
\[
  \pi_0 = \frac{1}{\displaystyle\sum_{i=0}^{\infty} \frac{\lambda_{i-1} \cdots \lambda_0}{\mu_i \cdots \mu_1}}.
\]

\textbf{Trường hợp đặc biệt:} Khi $\lambda_i = \lambda$ và $\mu_i = \mu$ với mọi $i \geq 1$:
\[
  \pi_j = \left(1 - \frac{\lambda}{\mu}\right) \left(\frac{\lambda}{\mu}\right)^j, \qquad j \in \N,
\]
với điều kiện $\lambda < \mu$ (phân phối hình học với tham số $\lambda/\mu$).

\begin{example}[Hàng đợi M/M/1]
\label{ex:mm1-queue}
Một quầy phục vụ có khách đến theo Poisson với tốc độ $\lambda$ và thời gian phục vụ mũ với tốc độ $\mu$.
Đây là quá trình sinh-tử với $\lambda_i = \lambda$ và $\mu_i = \mu$ cho mọi $i \geq 1$, $\mu_0 = 0$.

Nếu $\rho = \lambda/\mu < 1$ (hệ thống ổn định):
\[
  \pi_j = (1 - \rho)\rho^j, \qquad j \in \N.
\]

Số khách trung bình trong hệ thống:
\[
  \E[X] = \sum_{j=0}^{\infty} j\, \pi_j = \frac{\rho}{1-\rho} = \frac{\lambda}{\mu - \lambda}.
\]

Ví dụ: nếu $\lambda = 4$ khách/giờ, $\mu = 5$ khách/giờ, thì $\rho = 0.8$ và $\E[X] = 4$ khách.
\end{example}

% ============================================================================
\section{Chuỗi nhúng (The Embedded Chain)}
\label{sec:embedded}

\textbf{Câu chuyện:} Hãy tưởng tượng bạn quay video quán café cả ngày. Chuỗi Markov liên tục là toàn bộ video.
Nhưng nếu bạn chỉ chụp ảnh mỗi khi có sự thay đổi (khách đến hoặc đi),
bạn thu được một ``bộ ảnh'' --- đây chính là \term{chuỗi nhúng}.

Chuỗi nhúng ``bỏ qua thời gian'', chỉ ghi lại \emph{trạng thái nào} xuất hiện, không quan tâm \emph{bao lâu} ở mỗi trạng thái.

\begin{dinhnghia}[(Chuỗi nhúng -- \en{Embedded chain})]
Xét dãy $(T_n)_{n\in\N}$ các thời điểm nhảy của quá trình Markov liên tục $(X_t)_{t\in\R_+}$:
\[
  T_1 = \inf\{t > 0 : X_t \neq X_0\}, \qquad
  T_{n+1} = \inf\{t > T_n : X_t \neq X_{T_n}\}, \quad n \in \N.
\]
\term{Chuỗi Markov nhúng} (\en{embedded Markov chain}) là chuỗi rời rạc $(Z_n)_{n\in\N}$ với:
\[
  Z_0 = X_0, \qquad Z_n = X_{T_n}, \quad n \in \N.
\]
\end{dinhnghia}

\subsection{Ma trận chuyển của chuỗi nhúng}

Với quá trình sinh-tử, ma trận chuyển của chuỗi nhúng:
\[
  P_{i,i+1} = \frac{\lambda_i}{\lambda_i + \mu_i}, \qquad
  P_{i,i-1} = \frac{\mu_i}{\lambda_i + \mu_i}, \qquad i \in \N.
\]

Điều này là do:
\[
  \PP(\tau_{i,i+1} < \tau_{i,i-1}) = \frac{\lambda_i}{\lambda_i + \mu_i},
\]
theo tính chất cạnh tranh giữa hai biến mũ độc lập.

\textbf{Tổng quát:} Với chuỗi Markov liên tục có ma trận sinh $Q = [\lambda_{i,j}]$,
ma trận chuyển của chuỗi nhúng là:
\begin{equation}\label{eq:embedded-general}
  \hat{P}_{i,j} = \frac{\lambda_{i,j}}{-\lambda_{i,i}} = \frac{\lambda_{i,j}}{\sum_{l \neq i} \lambda_{i,l}}, \qquad i \neq j.
\end{equation}

\begin{example}[Từ ma trận $Q$ đến chuỗi nhúng -- 5 trạng thái]
Xét quá trình ở Ví dụ~\ref{ex:5state-generator} với ma trận sinh:
\[
  Q = \begin{bmatrix}
    -3 & 2 & 1 & 0 & 0 \\
    1 & -4 & 2 & 1 & 0 \\
    0 & 1 & -3 & 1 & 1 \\
    0 & 0 & 2 & -5 & 3 \\
    0 & 0 & 0 & 4 & -4
  \end{bmatrix}.
\]

Ma trận chuyển của chuỗi nhúng $\hat{P}$:
\[
  \hat{P} = \begin{bmatrix}
    0 & 2/3 & 1/3 & 0 & 0 \\
    1/4 & 0 & 2/4 & 1/4 & 0 \\
    0 & 1/3 & 0 & 1/3 & 1/3 \\
    0 & 0 & 2/5 & 0 & 3/5 \\
    0 & 0 & 0 & 1 & 0
  \end{bmatrix}.
\]

\textbf{Cách tính:} Hàng $i$ của $\hat{P}$ = hàng $i$ của $Q$ (bỏ phần tử đường chéo), chia cho $|\lambda_{i,i}|$.

Ví dụ hàng 1 (trạng thái 1): $\lambda_{1,0} = 1$, $\lambda_{1,2} = 2$, $\lambda_{1,3} = 1$, $|\lambda_{1,1}| = 4$.
\[
  \hat{P}_{1,0} = \frac{1}{4}, \quad \hat{P}_{1,2} = \frac{2}{4}, \quad \hat{P}_{1,3} = \frac{1}{4}.
\]
\end{example}

\begin{tomtat}
Chuỗi nhúng ``rời rạc hóa'' quá trình liên tục bằng cách chỉ ghi lại các trạng thái tại thời điểm nhảy.
Ma trận chuyển của chuỗi nhúng chỉ phụ thuộc vào \emph{tỷ lệ} các tốc độ, không phụ thuộc vào tốc độ tuyệt đối.

Để phục hồi quá trình liên tục từ chuỗi nhúng, cần thêm thông tin về \emph{thời gian ở mỗi trạng thái}
$\tau_i \sim \text{Exp}(-\lambda_{i,i})$.
\end{tomtat}

% ============================================================================
\section{Xác suất hấp thụ (Absorption Probabilities)}
\label{sec:absorption-prob}

\textbf{Câu chuyện:} Hãy tưởng tượng một con kiến bắt đầu ở vị trí $i$ trên một thanh gỗ.
Hai đầu thanh gỗ ($0$ và $N$) là ``tường'' --- khi chạm tường, kiến dừng lại vĩnh viễn.
Câu hỏi: ``Con kiến sẽ chạm tường nào? Với xác suất bao nhiêu?''

Đây là bài toán \term{xác suất hấp thụ} cho quá trình sinh-tử liên tục.
Vì chuỗi nhúng quyết định ``đi đâu'' (không quan tâm ``bao lâu''),
xác suất hấp thụ của quá trình liên tục \emph{bằng} xác suất hấp thụ của chuỗi nhúng.

\subsection{Thiết lập bài toán}

Xét quá trình sinh-tử trên $S = \{0, 1, \ldots, N\}$ với:
\begin{itemize}[nosep]
  \item Trạng thái $0$ và $N$ là hấp thụ: $\lambda_0 = 0$ (không sinh khi ở 0) và $\mu_N = 0$ (không tử khi ở $N$).
  \item Khi ở trạng thái $i$ ($1 \leq i \leq N-1$): sinh với tốc độ $\lambda_i$, tử với tốc độ $\mu_i$.
\end{itemize}

Đặt $g_0(i) = \PP(\text{bị hấp thụ vào trạng thái } 0 \mid X_0 = i)$.

\subsection{Phương trình cho xác suất hấp thụ}

Vì xác suất hấp thụ chỉ phụ thuộc chuỗi nhúng, ta áp dụng phân tích bước đầu tiên trên chuỗi nhúng:
\begin{equation}\label{eq:absorption-fsa}
  g_0(i) = \hat{P}_{i,i-1}\, g_0(i-1) + \hat{P}_{i,i+1}\, g_0(i+1),
  \qquad 1 \leq i \leq N-1,
\end{equation}
với điều kiện biên $g_0(0) = 1$ và $g_0(N) = 0$.

Thay $\hat{P}_{i,i-1} = \frac{\mu_i}{\lambda_i + \mu_i}$ và $\hat{P}_{i,i+1} = \frac{\lambda_i}{\lambda_i + \mu_i}$:
\begin{equation}\label{eq:absorption-bd}
  g_0(i) = \frac{\mu_i}{\lambda_i + \mu_i}\, g_0(i-1) + \frac{\lambda_i}{\lambda_i + \mu_i}\, g_0(i+1).
\end{equation}

Đây là phương trình ``con bạc'' (\en{gambler's equation}) ở dạng liên tục.

\subsection{Nghiệm tường minh}

Đặt $\rho_i = \mu_i / \lambda_i$ (tỷ số tử/sinh). Viết lại \eqref{eq:absorption-bd}:
\[
  g_0(i) - g_0(i-1) = \frac{\lambda_i}{\mu_i}\left[g_0(i+1) - g_0(i)\right],
\]
hay:
\[
  g_0(i+1) - g_0(i) = \rho_i \left[g_0(i) - g_0(i-1)\right].
\]

Đặt $d_i = g_0(i) - g_0(i-1)$. Ta có $d_{i+1} = \rho_i\, d_i$, suy ra:
\[
  d_k = d_1 \prod_{j=1}^{k-1} \rho_j = d_1 \prod_{j=1}^{k-1} \frac{\mu_j}{\lambda_j}, \qquad k \geq 2.
\]

Từ $g_0(0) = 1$ và $g_0(N) = 0$:
\[
  g_0(N) - g_0(0) = \sum_{k=1}^{N} d_k = -1.
\]

Đặt $R_k = \prod_{j=1}^{k-1} \frac{\mu_j}{\lambda_j}$ (với $R_1 = 1$). Ta được $d_1 \sum_{k=1}^{N} R_k = -1$, nên:

\begin{congthuc}
\textbf{Xác suất hấp thụ vào trạng thái $0$:}
\begin{equation}\label{eq:absorption-formula}
  g_0(i) = \frac{\displaystyle\sum_{k=i+1}^{N} R_k}{\displaystyle\sum_{k=1}^{N} R_k},
  \qquad \text{với } R_k = \prod_{j=1}^{k-1} \frac{\mu_j}{\lambda_j},\; R_1 = 1.
\end{equation}
\end{congthuc}

\textbf{Trường hợp đặc biệt} $\lambda_i = \lambda$, $\mu_i = \mu$ (tốc độ không đổi):

Nếu $\rho = \mu/\lambda \neq 1$: $R_k = \rho^{k-1}$, suy ra:
\begin{equation}\label{eq:absorption-constant}
  g_0(i) = \frac{\rho^i - \rho^N}{1 - \rho^N}.
\end{equation}

Nếu $\rho = 1$ (tức $\lambda = \mu$):
\[
  g_0(i) = 1 - \frac{i}{N}.
\]

\begin{example}[Xác suất hấp thụ với $N = 4$]
\label{ex:absorption-N4}
Xét quá trình sinh-tử trên $\{0, 1, 2, 3, 4\}$ với $\lambda_i = 2$, $\mu_i = 3$ cho $1 \leq i \leq 3$.

Ta có $\rho = \mu/\lambda = 3/2$. Áp dụng \eqref{eq:absorption-constant}:
\[
  g_0(i) = \frac{(3/2)^i - (3/2)^4}{1 - (3/2)^4}
  = \frac{(3/2)^i - 81/16}{1 - 81/16}
  = \frac{(3/2)^i - 81/16}{-65/16}.
\]

Tính cụ thể:
\begin{align*}
  g_0(1) &= \frac{3/2 - 81/16}{-65/16} = \frac{24/16 - 81/16}{-65/16} = \frac{-57/16}{-65/16} = \frac{57}{65} \approx 0.877, \\
  g_0(2) &= \frac{9/4 - 81/16}{-65/16} = \frac{36/16 - 81/16}{-65/16} = \frac{45}{65} = \frac{9}{13} \approx 0.692, \\
  g_0(3) &= \frac{27/8 - 81/16}{-65/16} = \frac{54/16 - 81/16}{-65/16} = \frac{27}{65} \approx 0.415.
\end{align*}

\textbf{Nhận xét:} Vì $\mu > \lambda$ (tốc độ tử $>$ tốc độ sinh), quá trình ``kéo về 0'' mạnh hơn,
nên xác suất chạm 0 khá cao ngay cả từ trạng thái 3.
\end{example}

% ============================================================================
\section{Thời gian hấp thụ trung bình (Mean Absorption Times)}
\label{sec:absorption-time}

\textbf{Câu chuyện:} Ta đã biết con kiến sẽ chạm tường. Nhưng \textbf{mất bao lâu}?
Trong thời gian liên tục, đây là câu hỏi quan trọng vì thời gian có ý nghĩa thực tế
(ví dụ: ``Trung bình bao lâu thì hệ thống tắt?'').

Khác với chuỗi rời rạc (đếm số bước), ở đây ta cần cộng các \emph{thời gian lưu trú} (\en{sojourn times})
tại mỗi trạng thái.

\subsection{Thiết lập}

Đặt $\eta(i) = \E[T_{\text{abs}} \mid X_0 = i]$, trong đó $T_{\text{abs}}$ là thời gian đến khi bị hấp thụ
(chạm $0$ hoặc $N$).

Với quá trình sinh-tử, phân tích bước đầu tiên \emph{trên chuỗi nhúng} cho:
\begin{equation}\label{eq:mean-abs-embedded}
  \eta(i) = \E[\tau_i] + \hat{P}_{i,i-1}\, \eta(i-1) + \hat{P}_{i,i+1}\, \eta(i+1),
  \qquad 1 \leq i \leq N-1,
\end{equation}
với điều kiện biên $\eta(0) = \eta(N) = 0$.

Ở đây $\E[\tau_i] = \frac{1}{\lambda_i + \mu_i}$ là thời gian trung bình ở trạng thái $i$ trước khi nhảy.

\textbf{Sự khác biệt so với rời rạc:} Trong trường hợp rời rạc, ta có ``$+1$'' (mỗi bước mất 1 đơn vị thời gian).
Ở đây, ta thay bằng ``$+\E[\tau_i]$'' --- mỗi bước mất thời gian $\text{Exp}(\lambda_i + \mu_i)$.

\subsection{Nghiệm}

Thay $\hat{P}_{i,i\pm 1}$ và $\E[\tau_i]$ vào \eqref{eq:mean-abs-embedded}:
\[
  \eta(i) = \frac{1}{\lambda_i + \mu_i} + \frac{\mu_i}{\lambda_i + \mu_i}\, \eta(i-1) + \frac{\lambda_i}{\lambda_i + \mu_i}\, \eta(i+1).
\]

Nhân cả hai vế với $(\lambda_i + \mu_i)$:
\[
  (\lambda_i + \mu_i)\eta(i) = 1 + \mu_i\, \eta(i-1) + \lambda_i\, \eta(i+1).
\]

Viết lại:
\begin{equation}\label{eq:mean-abs-recurrence}
  \lambda_i[\eta(i) - \eta(i+1)] = 1 + \mu_i[\eta(i-1) - \eta(i)].
\end{equation}

Đặt $\Delta_i = \eta(i) - \eta(i+1)$. Ta có quan hệ đệ quy:
\[
  \lambda_i \Delta_i = 1 - \mu_i \Delta_{i-1},
\]
hay:
\[
  \Delta_i = \frac{1}{\lambda_i} - \frac{\mu_i}{\lambda_i} \Delta_{i-1}
  = \frac{1}{\lambda_i} + \rho_i \Delta_{i-1}.
\]

\begin{example}[Thời gian hấp thụ trung bình với $N = 4$]
Tiếp tục Ví dụ~\ref{ex:absorption-N4} với $\lambda_i = 2$, $\mu_i = 3$.

Thời gian trung bình ở mỗi trạng thái: $\E[\tau_i] = 1/(\lambda_i + \mu_i) = 1/5$ giờ = 12 phút.

Hệ phương trình (nhân cả hai vế với 5):
\[
  \begin{cases}
    5\eta(1) = 1 + 3\eta(0) + 2\eta(2) = 1 + 2\eta(2), \\
    5\eta(2) = 1 + 3\eta(1) + 2\eta(3), \\
    5\eta(3) = 1 + 3\eta(2) + 2\eta(4) = 1 + 3\eta(2).
  \end{cases}
\]

Từ PT3: $\eta(3) = \frac{1 + 3\eta(2)}{5}$.

Thay vào PT2: $5\eta(2) = 1 + 3\eta(1) + 2 \cdot \frac{1 + 3\eta(2)}{5}$
$\Rightarrow 25\eta(2) = 5 + 15\eta(1) + 2 + 6\eta(2)$
$\Rightarrow 19\eta(2) = 7 + 15\eta(1)$.

Từ PT1: $\eta(1) = \frac{1 + 2\eta(2)}{5}$.

Thay: $19\eta(2) = 7 + 15 \cdot \frac{1 + 2\eta(2)}{5} = 7 + 3(1 + 2\eta(2)) = 10 + 6\eta(2)$.

Suy ra: $13\eta(2) = 10$, $\eta(2) = 10/13$.

Khi đó: $\eta(1) = \frac{1 + 20/13}{5} = \frac{33/13}{5} = \frac{33}{65}$,
$\eta(3) = \frac{1 + 30/13}{5} = \frac{43/13}{5} = \frac{43}{65}$.

\textbf{Kết quả:}
\[
  \eta(1) = \frac{33}{65} \approx 0.508, \quad
  \eta(2) = \frac{10}{13} \approx 0.769, \quad
  \eta(3) = \frac{43}{65} \approx 0.662.
\]

(Đơn vị: giờ, nếu $\lambda, \mu$ tính theo giờ.)

\textbf{Nhận xét:} $\eta(2) > \eta(3) > \eta(1)$. Trạng thái $2$ (ở giữa) mất lâu nhất
vì nó xa cả hai biên hấp thụ.
\end{example}

% ============================================================================
\section{Quan hệ giữa phân phối dừng liên tục và rời rạc}
\label{sec:continuous-vs-discrete}

Một câu hỏi tự nhiên: phân phối dừng $\pi$ của quá trình liên tục
và phân phối dừng $\hat{\pi}$ của chuỗi nhúng có quan hệ gì?

\begin{dinhly}[(Quan hệ $\pi$ và $\hat{\pi}$)]
Nếu $\pi$ là phân phối dừng của quá trình Markov liên tục (tức $\pi Q = 0$),
và $\hat{\pi}$ là phân phối dừng của chuỗi nhúng (tức $\hat{\pi} \hat{P} = \hat{\pi}$), thì:
\begin{equation}\label{eq:pi-hat-relation}
  \pi_i \propto \frac{\hat{\pi}_i}{-\lambda_{i,i}}, \qquad i \in S.
\end{equation}
Tức là: $\pi_i = \frac{\hat{\pi}_i / (-\lambda_{i,i})}{\sum_j \hat{\pi}_j / (-\lambda_{j,j})}$.
\end{dinhly}

\textbf{Ý nghĩa trực quan:} Phân phối dừng liên tục $\pi_i$ tỷ lệ với cả hai yếu tố:
\begin{enumerate}[nosep]
  \item Tần suất ``ghé thăm'' trạng thái $i$ (đo bằng $\hat{\pi}_i$).
  \item Thời gian ở lại trạng thái $i$ mỗi lần ghé (đo bằng $1/(-\lambda_{i,i})$).
\end{enumerate}
Trạng thái được ghé thăm thường xuyên VÀ giữ chân lâu sẽ có $\pi_i$ lớn.

% ============================================================================
\section{Bài tập (Exercises)}
\label{sec:ch9-exercises}

\begin{exercise}
Xét quá trình sinh-tử trên $\{0, 1, 2, 3\}$ với tốc độ sinh $\lambda_i = 2$ và tốc độ tử $\mu_i = 3$ cho $1 \leq i \leq 2$, $\lambda_0 = 2$, $\mu_3 = 3$, $\lambda_3 = \mu_0 = 0$.
\begin{enumerate}[nosep]
  \item Viết ma trận sinh $Q$.
  \item Tìm phân phối dừng $\pi$.
  \item Viết ma trận chuyển của chuỗi nhúng.
\end{enumerate}
\end{exercise}

\begin{exercise}
Cho chuỗi Markov liên tục hai trạng thái với $\alpha = 3$ và $\beta = 6$.
\begin{enumerate}[nosep]
  \item Tính $P(t)$ cho mọi $t > 0$.
  \item Tính thời gian trung bình ở mỗi trạng thái.
  \item Tìm phân phối giới hạn.
\end{enumerate}
\end{exercise}

\begin{exercise}[Hàng đợi M/M/1]
Khách đến quầy theo Poisson với tốc độ $\lambda = 3$ khách/giờ.
Thời gian phục vụ mỗi khách là biến mũ với tốc độ $\mu = 5$ phục vụ/giờ.
\begin{enumerate}[nosep]
  \item Viết ma trận sinh $Q$ (vô hạn chiều).
  \item Tìm phân phối dừng. Hệ thống có ổn định không?
  \item Tính số khách trung bình trong hệ thống.
  \item Tính xác suất hệ thống rỗng (không có khách nào).
\end{enumerate}
\end{exercise}

\begin{exercise}[Quá trình điện tín -- \en{Telegraph process}]
Tín hiệu điện tín chuyển giữa ``bật'' và ``tắt'' theo chuỗi Markov liên tục hai trạng thái
với tốc độ bật$\to$tắt là $\alpha = 2$ và tắt$\to$bật là $\beta = 3$.
\begin{enumerate}[nosep]
  \item Tìm xác suất tín hiệu đang ``bật'' tại thời điểm $t$, biết ban đầu ``bật''.
  \item Tìm phân phối giới hạn. Bao nhiêu phần trăm thời gian tín hiệu ``bật''?
  \item Thời gian trung bình một chu kỳ ``bật'' kéo dài bao lâu?
\end{enumerate}
\end{exercise}

\begin{exercise}[Xác suất hấp thụ]
Xét quá trình sinh-tử trên $\{0, 1, 2, 3, 4, 5\}$ với $\lambda_i = 1$ và $\mu_i = 2$ cho $1 \leq i \leq 4$,
trạng thái $0$ và $5$ hấp thụ.
\begin{enumerate}[nosep]
  \item Tính xác suất bị hấp thụ vào trạng thái $0$ xuất phát từ $i = 1, 2, 3, 4$.
  \item So sánh với trường hợp $\lambda_i = \mu_i = 1$ (tốc độ bằng nhau).
\end{enumerate}
\end{exercise}

\begin{exercise}[Quá trình Yule -- \en{Yule process}]
Quá trình Yule là quá trình sinh thuần túy với $\lambda_i = i\lambda$ (tốc độ sinh tỷ lệ với số cá thể hiện tại).
Bắt đầu với $X_0 = 1$.
\begin{enumerate}[nosep]
  \item Tính thời gian trung bình ở trạng thái $i$ (tức $\E[\tau_i]$).
  \item Tính thời gian trung bình để quần thể tăng từ $1$ lên $n$ cá thể.
  \item Giải thích tại sao quá trình Yule mô hình hóa sự tăng trưởng dân số.
\end{enumerate}
\end{exercise}

\begin{exercise}[Thời gian hấp thụ trung bình]
Xét quá trình sinh-tử trên $\{0, 1, 2, 3\}$ với $\lambda_i = 1$, $\mu_i = 1$ cho $i = 1, 2$,
trạng thái $0$ và $3$ hấp thụ.
\begin{enumerate}[nosep]
  \item Tính thời gian hấp thụ trung bình $\eta(1)$ và $\eta(2)$.
  \item So sánh với thời gian hấp thụ trong chuỗi rời rạc tương ứng (gambler's ruin với $p = q = 1/2$, $N = 3$).
\end{enumerate}
\end{exercise}


% --- Appendix ---
\appendix
\chapter{Nền tảng Xác suất\\{\normalsize\textit{Probability Background}}}
\label{ch:probability}

Phụ lục này tóm tắt các kiến thức xác suất cần thiết cho việc nghiên cứu chuỗi Markov.

% ============================================================================
\section{Không gian xác suất (Probability Spaces)}

\begin{dinhnghia}[(Không gian xác suất -- \en{Probability space})]
\term{Không gian xác suất} là bộ ba $(\Omega, \mathcal{F}, \PP)$ trong đó:
\begin{itemize}[nosep]
  \item $\Omega$: \term{không gian mẫu} (\en{sample space}) --- tập hợp tất cả các kết quả có thể.
  \item $\mathcal{F}$: \term{sigma-đại số} (\en{$\sigma$-algebra}) --- tập các \term{biến cố} (\en{events}).
  \item $\PP$: \term{độ đo xác suất} (\en{probability measure}) --- ánh xạ $\PP: \mathcal{F} \to [0,1]$.
\end{itemize}
\end{dinhnghia}

\textbf{Ví dụ:}
\begin{itemize}[nosep]
  \item Tung đồng xu: $\Omega = \{H, T\}$.
  \item Gieo xúc xắc: $\Omega = \{1, 2, 3, 4, 5, 6\}$.
  \item Kết quả nguyên: $\Omega = \N$.
  \item Kết quả thực không âm: $\Omega = \R_+$.
\end{itemize}

\term{Độ đo xác suất} $\PP: \mathcal{F} \to [0,1]$ thỏa mãn:
\begin{enumerate}[nosep]
  \item $\PP(\Omega) = 1$.
  \item $\PP\!\left(\bigcup_{n=1}^{\infty} A_n\right) = \sum_{n=1}^{\infty} \PP(A_n)$ khi $A_k \cap A_l = \emptyset$ với $k \neq l$ (\term{tính cộng đếm được} -- \en{countable additivity}).
\end{enumerate}

Tính chất quan trọng:
\[
  \PP(A \cup B) = \PP(A) + \PP(B) - \PP(A \cap B).
\]

% ============================================================================
\section{Xác suất có điều kiện và tính độc lập (Conditional Probability and Independence)}

\begin{dinhnghia}[(Xác suất có điều kiện -- \en{Conditional probability})]
Với hai biến cố $A, B \subset \Omega$ và $\PP(B) \neq 0$:
\[
  \PP(A \mid B) := \frac{\PP(A \cap B)}{\PP(B)}.
\]
\end{dinhnghia}

\textbf{Quy tắc xác suất toàn phần} (\en{law of total probability}): Nếu $\{A_n\}_{n\geq 1}$ là phân hoạch của $\Omega$ với $\PP(A_n) > 0$:
\[
  \PP(B) = \sum_{n=1}^{\infty} \PP(B \mid A_n)\,\PP(A_n) = \sum_{n=1}^{\infty} \PP(B \cap A_n).
\]

\textbf{Công thức Bayes} (\en{Bayes' theorem}):
\[
  \PP(A_k \mid B) = \frac{\PP(B \mid A_k)\,\PP(A_k)}{\sum_{n} \PP(B \mid A_n)\,\PP(A_n)}.
\]

\begin{dinhnghia}[(Tính độc lập -- \en{Independence})]
Hai biến cố $A$ và $B$ được gọi là \term{độc lập} (\en{independent}) nếu:
\[
  \PP(A \cap B) = \PP(A)\,\PP(B).
\]
Tương đương: $\PP(A \mid B) = \PP(A)$ (biết $B$ không ảnh hưởng đến $A$).
\end{dinhnghia}

\textbf{Độc lập nhiều biến cố:} $A_1, \ldots, A_n$ \term{độc lập tương hỗ} (\en{mutually independent}) nếu
với mọi tập con $I \subseteq \{1, \ldots, n\}$:
\[
  \PP\!\left(\bigcap_{i \in I} A_i\right) = \prod_{i \in I} \PP(A_i).
\]

% ============================================================================
\section{Biến ngẫu nhiên (Random Variables)}

\begin{dinhnghia}[(Biến ngẫu nhiên -- \en{Random variable})]
\term{Biến ngẫu nhiên} là ánh xạ $X: \Omega \to \R$ xác định trên không gian xác suất $(\Omega, \mathcal{F}, \PP)$.
\end{dinhnghia}

\term{Hàm chỉ báo} (\en{indicator function}) của biến cố $A$:
\[
  \ind_A(\omega) = \begin{cases} 1 & \text{nếu } \omega \in A, \\ 0 & \text{nếu } \omega \notin A. \end{cases}
\]

Hai biến ngẫu nhiên $X$ và $Y$ được gọi là \term{độc lập} nếu:
\[
  \PP(X \in A,\; Y \in B) = \PP(X \in A)\,\PP(Y \in B)
\]
với mọi tập đo được $A, B \subset \R$.

% ============================================================================
\section{Phân phối xác suất (Probability Distributions)}

\subsection{Phân phối rời rạc (Discrete Distributions)}

Biến ngẫu nhiên \term{rời rạc} (\en{discrete}) nhận giá trị đếm được.
Phân phối được xác định bởi $\PP(X = k)$ với mọi $k$.

\textbf{Các phân phối rời rạc quan trọng:}
\begin{itemize}[nosep]
  \item \textbf{Bernoulli}$(p)$: $\PP(X = 1) = p$, $\PP(X = 0) = 1-p$. $\E[X] = p$, $\text{Var}(X) = p(1-p)$.
  \item \textbf{Nhị thức} (\en{Binomial}) $B(n,p)$: $\PP(X = k) = \binom{n}{k} p^k (1-p)^{n-k}$.
    $\E[X] = np$, $\text{Var}(X) = np(1-p)$.
  \item \textbf{Hình học} (\en{Geometric}) $\text{Geo}(p)$: $\PP(X = k) = (1-p)^k p$, $k \in \N$.
    $\E[X] = (1-p)/p$, $\text{Var}(X) = (1-p)/p^2$.
  \item \textbf{Poisson}$(\lambda)$: $\PP(X = k) = e^{-\lambda} \frac{\lambda^k}{k!}$, $k \in \N$.
    $\E[X] = \lambda$, $\text{Var}(X) = \lambda$.
\end{itemize}

\subsection{Phân phối liên tục (Continuous Distributions)}

Biến ngẫu nhiên \term{liên tục} (\en{continuous}) có \term{hàm mật độ} (\en{density function}) $f: \R \to \R_+$ sao cho:
\[
  \PP(a \leq X \leq b) = \int_a^b f(x)\,dx, \qquad \int_{-\infty}^{\infty} f(x)\,dx = 1.
\]

\textbf{Các phân phối liên tục quan trọng:}
\begin{itemize}[nosep]
  \item \textbf{Đều} (\en{Uniform}) $U(a,b)$:
    $f(x) = \frac{1}{b-a}$ với $x \in [a,b]$. $\E[X] = (a+b)/2$.
  \item \textbf{Chuẩn} (\en{Gaussian/Normal}) $\mathcal{N}(\mu, \sigma^2)$:
    $f(x) = \frac{1}{\sqrt{2\pi\sigma^2}} e^{-(x-\mu)^2/(2\sigma^2)}$.
  \item \textbf{Mũ} (\en{Exponential}) $\text{Exp}(\lambda)$:
    $f(x) = \lambda e^{-\lambda x}$ với $x \geq 0$, và $\PP(X > t) = e^{-\lambda t}$.
    $\E[X] = 1/\lambda$, $\text{Var}(X) = 1/\lambda^2$.
  \item \textbf{Gamma} $\Gamma(n, \lambda)$:
    $f(x) = \frac{\lambda^n}{(n-1)!} x^{n-1} e^{-\lambda x}$ với $x \geq 0$, $n \in \N^*$.
    $\E[X] = n/\lambda$.
\end{itemize}

\begin{congthuc}
\textbf{Tính chất quan trọng của phân phối mũ:}

Nếu $X_1, \ldots, X_n$ độc lập với $X_i \sim \text{Exp}(\lambda_i)$, thì:
\begin{equation}\label{eq:min-exp}
  \min(X_1, \ldots, X_n) \sim \text{Exp}(\lambda_1 + \cdots + \lambda_n),
\end{equation}
và:
\begin{equation}\label{eq:compete-exp}
  \PP(X_1 < X_2) = \frac{\lambda_1}{\lambda_1 + \lambda_2}.
\end{equation}
Hai tính chất này rất quan trọng trong chuỗi Markov thời gian liên tục (Chương~\ref{ch:continuous}).

\textbf{Tính ``không nhớ''} (\en{memoryless property}): nếu $X \sim \text{Exp}(\lambda)$:
\[
  \PP(X > s + t \mid X > s) = \PP(X > t), \qquad s, t \geq 0.
\]
Đã đợi $s$ rồi, thời gian đợi thêm vẫn có cùng phân phối --- không ``nhớ'' đã đợi bao lâu.
\end{congthuc}

% ============================================================================
\section{Kỳ vọng toán học (Mathematical Expectation)}

\begin{dinhnghia}[(Kỳ vọng -- \en{Expectation})]
\term{Kỳ vọng} (\en{expected value}) của biến ngẫu nhiên $X$:
\begin{itemize}[nosep]
  \item Trường hợp rời rạc: $\E[X] = \sum_{k} k \cdot \PP(X = k)$.
  \item Trường hợp liên tục: $\E[X] = \int_{-\infty}^{\infty} x\, f(x)\,dx$.
\end{itemize}
\end{dinhnghia}

\textbf{Tính chất:}
\begin{itemize}[nosep]
  \item \term{Tuyến tính} (\en{linearity}): $\E[aX + bY] = a\E[X] + b\E[Y]$ (luôn đúng, kể cả khi không độc lập).
  \item Nếu $X, Y$ độc lập: $\E[XY] = \E[X]\,\E[Y]$.
  \item $\E[g(X)] = \sum_k g(k)\,\PP(X = k)$ (rời rạc) hoặc $\int g(x) f(x)\,dx$ (liên tục).
  \item \term{Đơn điệu} (\en{monotonicity}): nếu $X \leq Y$ h.c.c., thì $\E[X] \leq \E[Y]$.
\end{itemize}

\term{Phương sai} (\en{variance}):
\[
  \text{Var}(X) = \E[(X - \E[X])^2] = \E[X^2] - (\E[X])^2.
\]

Nếu $X, Y$ độc lập: $\text{Var}(X + Y) = \text{Var}(X) + \text{Var}(Y)$.

% ============================================================================
\section{Bất đẳng thức quan trọng (Important Inequalities)}

\begin{dinhly}[(Bất đẳng thức Markov -- \en{Markov's inequality})]
Nếu $X \geq 0$ và $a > 0$:
\[
  \PP(X \geq a) \leq \frac{\E[X]}{a}.
\]
\end{dinhly}

\begin{dinhly}[(Bất đẳng thức Chebyshev -- \en{Chebyshev's inequality})]
Với mọi $k > 0$:
\[
  \PP(|X - \E[X]| \geq k) \leq \frac{\text{Var}(X)}{k^2}.
\]
\end{dinhly}

\textbf{Ứng dụng:} Bất đẳng thức Chebyshev cho ước lượng ``nhanh'' về xác suất
biến ngẫu nhiên lệch xa kỳ vọng, mà không cần biết phân phối chính xác.

% ============================================================================
\section{Hàm sinh xác suất (Probability Generating Functions)}
\label{sec:pgf}

\begin{dinhnghia}[(Hàm sinh xác suất -- \en{Probability generating function})]
Cho biến ngẫu nhiên $X$ nhận giá trị trong $\N$. \term{Hàm sinh xác suất} (\en{p.g.f.}) của $X$ là:
\[
  G_X(s) = \E[s^X] = \sum_{k=0}^{\infty} \PP(X = k)\, s^k, \qquad |s| \leq 1.
\]
\end{dinhnghia}

\textbf{Tính chất:}
\begin{itemize}[nosep]
  \item $G_X(1) = 1$ (vì tổng xác suất bằng 1).
  \item $G_X'(1) = \E[X]$.
  \item $G_X''(1) = \E[X(X-1)]$, nên $\text{Var}(X) = G_X''(1) + G_X'(1) - [G_X'(1)]^2$.
  \item Nếu $X, Y$ độc lập: $G_{X+Y}(s) = G_X(s) \cdot G_Y(s)$.
  \item $\PP(X = k) = \frac{G_X^{(k)}(0)}{k!}$ (khôi phục phân phối từ hàm sinh).
\end{itemize}

\textbf{Ứng dụng trong chuỗi Markov:} Hàm sinh rất hữu ích để phân tích:
\begin{itemize}[nosep]
  \item Phân phối thời gian quay lại (Chương~\ref{ch:first-step}): $F(s) = \sum_{n=1}^{\infty} f_{ii}^{(n)} s^n$.
  \item Quan hệ $P(s) = F(s) \cdot P(s) + 1$ giữa hàm sinh của $P_{ii}^{(n)}$ và $f_{ii}^{(n)}$.
  \item Số lần thăm trạng thái: phân phối hình học có $G(s) = \frac{p\, s}{1 - (1-p)s}$.
\end{itemize}

\begin{example}[Hàm sinh của phân phối hình học]
Nếu $X \sim \text{Geo}(p)$ (tức $\PP(X = k) = (1-p)^k p$, $k \geq 0$):
\[
  G_X(s) = \sum_{k=0}^{\infty} (1-p)^k p\, s^k = \frac{p}{1 - (1-p)s}.
\]
Kiểm tra: $G_X'(1) = \frac{p(1-p)}{[1-(1-p)]^2} = \frac{1-p}{p} = \E[X]$. Đúng!
\end{example}

% ============================================================================
\section{Kỳ vọng có điều kiện (Conditional Expectation)}

\begin{dinhnghia}[(Kỳ vọng có điều kiện -- \en{Conditional expectation})]
\[
  \E[X \mid B] = \sum_k k\,\PP(X = k \mid B), \qquad \text{khi } \PP(B) > 0.
\]
\end{dinhnghia}

\textbf{Quy tắc kỳ vọng lặp} (\en{law of iterated expectations}/\en{tower property}):
\[
  \E[X] = \sum_{n} \E[X \mid A_n]\,\PP(A_n),
\]
khi $\{A_n\}$ là phân hoạch của $\Omega$.

Tổng quát hơn, với hai biến ngẫu nhiên:
\[
  \E[X] = \E[\E[X \mid Y]] = \sum_y \E[X \mid Y = y]\,\PP(Y = y).
\]

\textbf{Ứng dụng quan trọng --- Phân tích bước đầu tiên:}

Đặt $h(i) = \E[T \mid X_0 = i]$ (thời gian chạm). Điều kiện hóa theo $X_1$:
\[
  h(i) = \E[T \mid X_0 = i] = 1 + \sum_j P_{i,j}\, h(j).
\]
Đây chính là kỹ thuật cốt lõi của Chương~\ref{ch:first-step}.

% ============================================================================
\section{Phân phối đồng thời và mật độ có điều kiện}

\begin{dinhnghia}[(Phân phối đồng thời -- \en{Joint distribution})]
\term{Phân phối đồng thời} của $(X, Y)$ rời rạc:
\[
  p_{X,Y}(x, y) = \PP(X = x,\; Y = y).
\]
\term{Phân phối biên} (\en{marginal}):
\[
  \PP(X = x) = \sum_y p_{X,Y}(x, y), \qquad
  \PP(Y = y) = \sum_x p_{X,Y}(x, y).
\]
\end{dinhnghia}

Trong trường hợp liên tục, \term{mật độ đồng thời} (\en{joint density}) $f_{X,Y}(x,y)$ thỏa mãn:
\[
  \PP(X \in A,\; Y \in B) = \int_A \int_B f_{X,Y}(x,y)\,dy\,dx.
\]

\term{Mật độ có điều kiện} (\en{conditional density}):
\[
  f_{X|Y}(x \mid y) = \frac{f_{X,Y}(x,y)}{f_Y(y)}, \qquad f_Y(y) > 0.
\]

\textbf{Kỳ vọng có điều kiện} (trường hợp liên tục):
\[
  \E[X \mid Y = y] = \int_{-\infty}^{\infty} x\, f_{X|Y}(x \mid y)\,dx.
\]

% ============================================================================
\section{Tổng kết (Summary)}

\begin{tomtat}
Các công cụ xác suất chính dùng trong tài liệu:
\begin{itemize}[nosep]
  \item \textbf{Xác suất có điều kiện} và \textbf{quy tắc xác suất toàn phần}: nền tảng cho phân tích bước đầu tiên (Chương~\ref{ch:first-step}).
  \item \textbf{Tính độc lập}: cần cho bước đi ngẫu nhiên và quá trình Poisson.
  \item \textbf{Phân phối mũ} và tính chất min/cạnh tranh: then chốt cho chuỗi Markov liên tục (Chương~\ref{ch:continuous}).
  \item \textbf{Kỳ vọng có điều kiện}: dùng để tính thời gian chạm và xác suất hấp thụ.
  \item \textbf{Hàm sinh xác suất}: công cụ đại số để phân tích phân phối thời gian quay lại.
  \item \textbf{Bất đẳng thức Markov/Chebyshev}: ước lượng xác suất không cần biết phân phối chính xác.
\end{itemize}
\end{tomtat}


\end{document}
